\section{Question and results}

Paper G8 assigns a measure to sets of walks of one length on a gain graph, and leaves three items as assumptions: the second order of the rule, the exclusion of a first-order weighting of returning walks, and any relation between the measures at different lengths \cite{g8}. Paper G9 takes the simplest such relation, consistency under continuation, as a requirement and finds what it forces on the carrier \cite{g9}. It leaves open a reason for the requirement, an error bound for the position of a ray, and whether gains of bounded order can imitate a weak force. The present paper asks what the same requirement decides about the rule itself, and returns to the last two of these open items.

\begin{enumerate}[label=(\roman*)]
\item On a digraph in which every vertex has an outgoing arc, a family of set functions on walks of each length from one vertex defines one function on sets of unending walks exactly when it is consistent under continuation, the function at each length being rescaled by a power of the factor (Lemma~\ref{lem:cyl}). The requirement of G9 is thus the statement that events are sets of histories.
\item On a digraph without loops and of constant out-degree, an additive switching-invariant rule that is consistent under continuation is a multiple of the number of walks: the weighting of returning walks by their holonomy, which switching invariance allows, is excluded (Proposition~\ref{prop:add}).
\item On the signed arc digraph of a line, a rule that adds over end vertices a function of the modulus of the summed gains is consistent under continuation only if that function is the square (Theorem~\ref{thm:quad}); the same holds on every finite connected signed digraph that is norm-preserving of degree two (Theorem~\ref{thm:deg2}). On any norm-preserving gain digraph with a merging pair of walks, the only consistent power of the modulus is the second (Theorem~\ref{thm:power}). On a norm-preserving carrier of degree two with gains in the $q$-th roots of unity, $q$ even, on which every gain reaches both vertices of a block in abundance, the consistent rules of this form are exactly the combinations of the squared moduli in the conjugate characters; they form a space of dimension $\varphi(q)/2$ for $q>2$ and of dimension one for $q=2$, and the rule is unique up to a factor exactly for $q\in\{2,4,6\}$ (Theorem~\ref{thm:conj}). Such a rule is non-negative exactly when its weights are, and if it is multiplicative on products of carriers it is the squared modulus in one conjugate character (Proposition~\ref{prop:compo}). If the gain of a walk is itself recorded, so that the carrier is the signed lift with one sheet for each pair of opposite gains, the rule is unique for every $q$: it is the square of the integer amplitude, the mean of the squared moduli over the characters in which $-1$ acts as $-1$ (Corollary~\ref{cor:lift}). Without any choice of records, for any finite abelian group of gains and any degree, and under a hypothesis on the quotients of gains that reach two vertices with a common successor, the consistent switching-invariant pair rules of \cite{g8} are exactly the real combinations of the squared-norm rules of the characters for which the step is norm-preserving (Theorem~\ref{thm:kernel}). In degree two these are the characters in which $-1$ acts as $-1$, the kernel changes sign under $-1$, and for the $q$-th roots of unity the space has dimension $\lceil q/4\rceil$ and contains both families (Corollary~\ref{cor:kernel2}). For blocks of degree three and four with gains of order $3$, $6$ or $4$, with disjoint successor sets and abundant gains, a consistent rule that adds a function of the modulus over end vertices is again the square (Proposition~\ref{prop:deg34}); for blocks $\rho_z\kappa_yw^{zy}$ of every degree $d$ dividing $q$, under the same conditions, the consistent rules of that form are exactly the combinations of the conjugate squared moduli, as in degree two (Theorem~\ref{thm:fourier}). Finally the form of the blocks is not needed at all: on every norm-preserving carrier, of any degree, in which two vertices have a common successor and the events at them are rich enough, the consistent rules of that form are exactly the combinations of the conjugate squared moduli (Theorem~\ref{thm:vsl}, Corollary~\ref{cor:vsl}). On a finite carrier whose derived digraph is strongly connected no hypothesis on events is needed: the order that enters is that of the gains that can meet at one length (Proposition~\ref{prop:abund}, Corollary~\ref{cor:abund}). On carriers with a full block, that is, $d$ vertices with one common set of $d$ successors and a matrix of gains $B$ with $B^*B=d$ for the faithful character, and with abundant gains, the modulus is not needed either; every finite norm-preserving digraph of degree two consists of such blocks. A consistent rule that adds over end vertices any function of the summed gains that vanishes at zero is a quadratic form, and if it is unchanged by switching it is a combination of the conjugate squared moduli (Theorem~\ref{thm:block}, Corollary~\ref{cor:block}). Nor is the sum: for any $d\times d$ matrix of gains in a finite abelian group, $d\ge2$, a function of the multiset of gains that satisfies the block law is an additive function plus a quadratic form, and if the entries generate the group it is a combination of the count and of squared moduli of gains summed through characters; on a carrier with such a block and abundant gains this holds for the consistent rules (Theorem~\ref{thm:multiset}, Corollaries~\ref{cor:multiset}, \ref{cor:multicar}). Nor need the function be the same at every end vertex, or the carrier have a block: on a finite strongly connected gain digraph in which every vertex has $d\ge2$ outgoing arcs, without parallel arcs and with abundant gains, a consistent rule that adds over end vertices a function of the multiset of gains that may depend on the vertex, and vanishes on the empty multiset, is a combination of the count and of the squares $|\sigma_\psi|^2$ for the characters $\psi$ in which the step is norm-preserving, with the same coefficients at every vertex (Theorem~\ref{thm:vertexlaw}, Corollary~\ref{cor:vertexlaw}). The abundance of every gain at every vertex at one length fails on periodic carriers such as a ring with an even number of sites; it is not needed: if the derived digraph on pairs (vertex, gain) is strongly connected, of any period, the consistent rules of this kind are, as functions on events, exactly $a|S|+\sum c\,\lVert v^\psi_S\rVert^2$ over those characters, with constants that depend neither on the vertex nor on the length, and characters whose restrictions to the subgroup of quotients of gains that can meet at one vertex at one length are equal or inverse give the same rule (Theorem~\ref{thm:periodic}, Corollary~\ref{cor:periodic}). If such a consistent rule is in addition multiplicative on the product events of the carrier with itself, at all large lengths, it is zero, the count, or the squared norm for one of these classes of characters: mixtures of the count and a square, and of two squares, are excluded (Corollary~\ref{cor:composition}).
\item A gain that depends on position translates the quasi-momentum by a fixed amount per step exactly only if it is the linear gain; a translation by $\Phi$ therefore needs a gain of order at least $2\pi/\Phi$ (Proposition~\ref{prop:exact}). For gains of order $q\ge3$ no quasi-energy belongs to the spectra of the free step multiplied by all $q$ gain values, except four band edges for $q=4$ (Proposition~\ref{prop:local}). Computed instances of staircases of gains of order $3$, $4$ and $6$ are recorded in Section~\ref{sec:force}; a packet, computed for order $4$, does not follow the ray of the mean field.
\item For the linear gain, the centroid of a packet follows the ray averaged over the quasi-momentum distribution of the packet, displaced by the connection of the branch eigenvector, with a remainder that decreases about as the first power of the field in the instances computed (Section~\ref{sec:ray}). No bound is proved.
\end{enumerate}

Section~\ref{sec:limits} states what is assumed and what is not obtained. The statements are mathematical and concern the constructed trials of the program collection named GH-1 and GH-2 there (gain graphs and gain digraphs with a reversible step, the carriers of \cite{g8} and \cite{g9}). Neither quantum nor classical mechanics is derived.

\section{Setting}\label{sec:setting}

The notation is that of \cite{g9}. A gain digraph $D$ has a finite or countable vertex set $V$, arcs $e$ with initial vertex $s(e)$, terminal vertex $t(e)$ and gain $\varphi(e)$ in a group $G$, at most one arc from a given vertex to a given vertex, and finitely many arcs leaving each vertex; a character $\chi\colon G\to\mathrm{U}(1)$ is fixed and $A_{yx}=\chi(\varphi(e))$ for an arc $e$ from $x$ to $y$, $A_{yx}=0$ if there is none. An event $S$ is a set of walks of one length $j$ from one vertex $x$, with vector $v_S(y)=\sum_{w\in S,\,t(w)=y}\chi(\varphi(w))$, and $\operatorname{Ext}S$ is the event of all walks of length $j+1$ whose first $j$ arcs form a walk of $S$; $v_{\operatorname{Ext}S}=Av_S$. A family $(\mu_j)_{j\ge0}$ of real functions on the events of each length is \emph{consistent under continuation with factor $c>0$} if
\begin{equation}\label{eq:cons}
\mu_{j+1}(\operatorname{Ext}S)=c\,\mu_j(S)\qquad\text{for every event }S\text{ of every length }j .
\end{equation}
A \emph{rule} gives $\mu_j$ by one formula for all $j$. The rule of \cite{g8} is $\mu(S)=\lVert v_S\rVert^2$; by Lemma~3.1 of \cite{g9} it satisfies \eqref{eq:cons} for all events exactly when $A^*A=c$ on the vectors concerned, and $D$ is then called norm-preserving of degree $c$.

\section{Events as sets of histories}\label{sec:cyl}

Suppose every vertex has an outgoing arc. A \emph{history} from $x$ is an unending walk $(e_1,e_2,\dots)$ with $s(e_1)=x$. For a walk $w$ of length $j$ from $x$ let $\operatorname{cyl}(w)$ be the set of histories whose first $j$ arcs are $w$, and for an event $S$ let $\operatorname{cyl}(S)=\bigcup_{w\in S}\operatorname{cyl}(w)$. For a fixed starting vertex the sets $\operatorname{cyl}(S)$ form an algebra of subsets of the histories from that vertex, the algebra of cylinder sets. The assumption on outgoing arcs is needed: a walk that ends at a vertex without one has an empty cylinder.

\begin{lemma}\label{lem:cyl}
Let every vertex have an outgoing arc, and fix a starting vertex; events and histories are those from that vertex.
\begin{enumerate}[label=(\alph*)]
\item $\operatorname{cyl}(\operatorname{Ext}S)=\operatorname{cyl}(S)$, and for events $S,S'$ of lengths $j\le j'$ one has $\operatorname{cyl}(S)=\operatorname{cyl}(S')$ if and only if $S'=\operatorname{Ext}^{\,j'-j}S$.
\item A family $(\nu_j)$ of functions on events defines a function $\nu$ on the algebra of cylinder sets by $\nu(\operatorname{cyl}S)=\nu_j(S)$ if and only if $\nu_{j+1}(\operatorname{Ext}S)=\nu_j(S)$ for all $S$ and $j$.
\item A family $(\mu_j)$ is consistent under continuation with factor $c$ if and only if $\nu_j=c^{-j}\mu_j$ defines a function on the algebra of cylinder sets.
\end{enumerate}
\end{lemma}
\begin{proof}
(a) A history begins with a walk of $S$ exactly when its first $j+1$ arcs form a walk of $\operatorname{Ext}S$. If $\operatorname{cyl}(S)=\operatorname{cyl}(S')$, a walk $w'$ of length $j'$ lies in $S'$ exactly when the histories beginning with $w'$ lie in $\operatorname{cyl}(S)$; since every walk is the beginning of a history, this holds exactly when the first $j$ arcs of $w'$ form a walk of $S$. (b) follows from (a), and (c) from (b).
\end{proof}

Condition \eqref{eq:cons} is therefore not a dynamical law added to the rule. It says that an event is a set of histories, and that the measure of a set of histories does not depend on the length at which the set is described. For a system of finitely many states with a unitary step the decoherence functional of Dowker, Johnston and Surya is defined on exactly this algebra, with cylinder sets of finite strings as generators \cite{djsu}; they show that it does not extend to the generated $\sigma$-algebra in general, which is not used here. That events are sets of histories is an assumption about the event structure. It is supplied. What follows is what it decides once it is made.

\section{Additive rules}\label{sec:add}

By Lemma~3.1 of \cite{g8} a switching-invariant additive rule has the form $\mu(S)=\sum_{w\in S}F(w)$ with $F(w)=a$ for a walk that does not end at its starting vertex and $F(w)=F_{\mathrm c}(\varphi(w))$, a class function of the gain, for a walk that does. Such a rule is consistent under continuation with factor $c$ exactly when
\begin{equation}\label{eq:addcons}
\sum_{e\colon s(e)=t(w)}F(we)=c\,F(w)\qquad\text{for every walk }w\text{ from }x .
\end{equation}

\begin{proposition}\label{prop:add}
Let $D$ be a gain digraph without loops in which every vertex has $d\ge1$ outgoing arcs, and let an additive rule of the form above satisfy \eqref{eq:addcons} with a factor $c>0$. If $a\ne0$, then $c=d$ and $F_{\mathrm c}(g)=a$ for every $g$ that is the gain of a closed walk at $x$. If $a=0$, then $F_{\mathrm c}(g)=0$ for every such $g$. In both cases $\mu$ is $a$ times the number of walks.
\end{proposition}
\begin{proof}
A closed walk $w$ at $x$ with gain $g$ has $d$ continuations, none of them closed, since there are no loops; so \eqref{eq:addcons} reads $da=cF_{\mathrm c}(g)$. If $a=0$ this gives $F_{\mathrm c}(g)=0$. Let $a\ne0$. If some walk from $x$ ends at a vertex $y\ne x$ without an arc to $x$, its continuations are open and \eqref{eq:addcons} reads $da=ca$, so $c=d$ and then $F_{\mathrm c}(g)=a$. Otherwise every walk $w$ ending at $y\ne x$ has exactly one closed continuation, with some gain $g'$, and \eqref{eq:addcons} reads $(d-1)a+F_{\mathrm c}(g')=ca$. Every closed walk of positive length arises in this way, so $F_{\mathrm c}$ has the single value $f=(c-d+1)a$ on their gains; with $da=cf$ this gives $(c-d)(c+1)=0$, hence $c=d$ and $f=a$. The closed walk of length $0$ gives $da=cF_{\mathrm c}(1)$ and the same value.
\end{proof}

The weights $2$ and $1/2$ of the example of Directors A, B and G for $\mathbb{Z}_3$, quoted after Lemma~3.1 of \cite{g8}, are thus not the values of a function on sets of histories: on the complete graph on four vertices with one gain of order three, with $a=1$, $F_{\mathrm c}(1)=2$ and $F_{\mathrm c}(\omega)=F_{\mathrm c}(\omega^{-1})=1/2$, the ratios in \eqref{eq:addcons} take the values $1.5$, $2.5$ and $4$ on walks of length two \cite{comp4s}. For two triangles with a common vertex, where the out-degrees are $2$ and $4$, no factor among those tried makes a rule with $a\ne0$ consistent \cite{comp4s}. A first-order rule on histories is the count of walks, which is the uniform classical measure on histories and does not depend on the gains.

\section{Rules that depend on the modulus}\label{sec:mod}

Let $\phi\colon[0,\infty)\to\mathbb{R}$ with $\phi(0)=0$, and put
\begin{equation}\label{eq:phi}
\mu_\phi(S)=\sum_{y\in V}\phi\bigl(|v_S(y)|\bigr).
\end{equation}
This keeps two assumptions of \cite{g8}: gains of walks with a common end vertex are added, and end vertices are records over which the rule adds. A rule \eqref{eq:phi} is said below to be of the \emph{record-sum form}. It drops the second order. The rule of \cite{g8} is $\phi(t)=t^2$, its rule of fourth order is $\phi(t)=t^4$, and with trivial gains $\phi(t)=t$ is the number of walks.

\subsection{The signed line}
Let $D$ be the signed arc digraph of the line $\mathbb{Z}$ of \cite{g9}: vertices $(x,R)$ and $(x,L)$, arcs from either to $(x+1,R)$ and to $(x-1,L)$, all gains $+1$ except $-1$ on the arcs from $(x,L)$ to $(x-1,L)$. The vectors $v_S$ have integer components, so only the values of $\phi$ at integers enter.

\begin{theorem}\label{thm:quad}
On the signed arc digraph of the line, the rule \eqref{eq:phi} with $\phi(1)\ne0$ is consistent under continuation with a factor $c>0$ if and only if $c=2$ and $\phi(n)=n^2\phi(1)$ for every integer $n\ge0$.
\end{theorem}
\begin{proof}
For one walk $w$, $\mu_\phi(\{w\})=\phi(1)$ and $\operatorname{Ext}\{w\}$ consists of two walks with different end vertices, so $2\phi(1)=c\,\phi(1)$ and $c=2$. Let $a\ge1$. A walk from $(0,R)$ is a sequence of steps to the right and to the left, and its gain is $-1$ to the power of the number of steps to the left that follow a step to the left. Among the walks of length $a+1$ with $a$ steps to the right and one to the left there are $a$ that end with a step to the right and one that ends with the step to the left; all have gain $+1$ and end at site $a-1$. Let $S$ consist of these $a+1$ walks. Then $v_S$ has the components $a$ at $(a-1,R)$ and $1$ at $(a-1,L)$ and no others, and $v_{\operatorname{Ext}S}=Av_S$ has the components $a+1$ at $(a,R)$ and $a-1$ at $(a-2,L)$. Consistency gives
\[
\phi(a+1)+\phi(a-1)=2\phi(a)+2\phi(1),\qquad a\ge1 ,
\]
and with $\phi(0)=0$ induction gives $\phi(n)=n^2\phi(1)$. Conversely $A^*A=2$ and $\phi(t)=t^2$ is consistent by Lemma~3.1 of \cite{g9}.
\end{proof}

The recursion is the parallelogram identity $\phi(|a+b|)+\phi(|a-b|)=2\phi(|a|)+2\phi(|b|)$ at $b=1$; the full identity holds at every site because the two arcs arriving at a site continue to one vertex with equal signs and to another with opposite signs, which is the zigzag gain $-1$ of Theorem~4.3 of \cite{g9}. A numerical solution of the consistency equations over $4000$ sampled events of length $12$ returns $\phi(n)/\phi(1)=n^2$ for $n\le8$ to $3\times10^{-13}$, and rejects $\phi(n)=n$ and $\phi(n)=n^4$ \cite{comp4s}. With all gains $+1$ on the same digraph the rule $\phi(t)=t$ is consistent, being the count.

\subsection{Norm-preserving digraphs}
\begin{theorem}\label{thm:power}
Let $D$ be a gain digraph with a character $\chi$ of modulus one and $A^*A=d$, $d\ge2$. Suppose two walks of equal length from a common vertex end at distinct vertices $y_1,y_2$ that have a common successor. Then for $p>0$ the rule $\mu_p(S)=\sum_y|v_S(y)|^p$ is consistent under continuation with some factor $c>0$ if and only if $p=2$.
\end{theorem}
\begin{proof}
Single walks give $c=d$, as every vertex has $d$ successors. Let $w_1,w_2$ be the two walks, $a,b$ the values of $\chi$ on their gains, and $S=\{w_1,w_2\}$, so $\mu_p(S)=2$. Let $Z$ be the set of the $m\ge1$ common successors of $y_1,y_2$. For $z\in Z$ the component of $Av_S$ is $aA_{zy_1}+bA_{zy_2}$, of modulus $|1+u_z|$ with $u_z=(b/a)\,\overline{A_{zy_1}}A_{zy_2}$; the other $2(d-m)$ successors carry components of modulus one. Orthogonality of the columns $y_1,y_2$ of $A$ gives $\sum_{z\in Z}u_z=0$. Hence $\mu_p(\operatorname{Ext}S)=2(d-m)+h(p)$ with $h(p)=\sum_{z\in Z}|1+u_z|^p$, and consistency requires $h(p)=2m$. Now $h(2)=\sum_z(2+2\operatorname{Re}u_z)=2m$. For $t\ge0$ the inequality between the weighted arithmetic and geometric means of $1$ and $t^2$ gives $t^p\le(1-\tfrac p2)+\tfrac p2t^2$ for $0<p<2$, and that of $1$ and $t^p$ gives $t^2\le(1-\tfrac2p)+\tfrac2pt^p$ for $p>2$, with equality only at $t=1$. Summing over $z\in Z$ with $t=|1+u_z|$ and $h(2)=2m$: $h(p)\le(1+\tfrac p2)m<2m$ for $p<2$, and $h(p)\ge(1+\tfrac p2)m>2m$ for $p>2$. This termwise argument replaces a convexity argument of earlier issues; it was found in a formal proof and is checked in Lean~4 with Mathlib \cite{comp4s}.
\end{proof}

\begin{theorem}[Degree two]\label{thm:deg2}
Let $D$ be a finite gain digraph with gains $\pm1$ that is norm-preserving of degree two and whose underlying graph is connected. The rule \eqref{eq:phi} with $\phi(1)\ne0$ is consistent under continuation with a factor $c>0$ if and only if $c=2$ and $\phi(n)=n^2\phi(1)$ for every integer $n\ge0$.
\end{theorem}
\begin{proof}
Single walks give $c=2$. By Theorem~4.3 of \cite{g9} every vertex has two incoming and two outgoing arcs, so $D$ is strongly connected, and the vertices fall into blocks: pairs $\{x,x'\}$ with common successors $\{y,y'\}$ and signs $s_1=A_{yx}$, $s_2=A_{yx'}$, $s_3=A_{y'x}$, $s_4=A_{y'x'}$ with $s_1s_2s_3s_4=-1$. Fix a block and a starting vertex. Let $B$ be the matrix of the numbers of arcs; $B/2$ is doubly stochastic and irreducible, with some period $r$, so the number of walks of length $j$ from the starting vertex to a given vertex tends to infinity as $j\to\infty$ through the lengths of one residue class modulo $r$ and is zero for the other lengths. Because $x$ and $x'$ have a common successor, their residue classes coincide. Given $n\ge1$ choose a length with at least $2n-1$ walks to $x$ and one walk to $x'$; among the former, $n$ have equal gains. Let $S$ consist of these $n$ walks and one walk to $x'$. Then $v_S$ has the components $\pm n$ at $x$ and $\pm1$ at $x'$ and no others, and $Av_S$ has the components $s_1(\pm n)+s_2(\pm1)$ at $y$ and $s_3(\pm n)+s_4(\pm1)$ at $y'$, of moduli $n+1$ and $n-1$ in some order because $s_1s_2=-s_3s_4$. Consistency gives $\phi(n+1)+\phi(n-1)=2\phi(n)+2\phi(1)$, and the conclusion follows as in Theorem~\ref{thm:quad}. The converse is Lemma~3.1 of \cite{g9}.
\end{proof}

Theorem~\ref{thm:deg2} covers the rings and every other carrier of the constructed trial GH-2 with real gains; the choice of the base digraph does not matter. For $40$ random carriers of this kind on up to twelve vertices, and for $10$ carriers with gains of order four, where the moduli are those of Gaussian integers, a numerical solution of the consistency equations returns the square \cite{comp4s}; gains that are not real are treated in the next subsection. A first version of that computation failed on two carriers through an unknown that occurred in no equation, and is kept.

The hypothesis of Theorem~\ref{thm:power} on the pair of walks is the hypothesis of Lemma~3.1(b) of \cite{g9} together with a common successor; it holds on the ring and the line. It cannot be dropped. On the countable directed binary tree, with every gain $+1$, each vertex has two successors and distinct vertices have none in common, so $A^*A=2$; from the root every walk is the only walk to its end vertex, every event has $\mu_p(S)=|S|$, and every $p>0$ is consistent with factor $2$. This example is due to Director~G \cite{dg}; it shows that branching together with $A^*A=d$ does not suffice, and it corrects a summary sentence of the first issue of this paper, which named ``any branching norm-preserving carrier''. Theorem~\ref{thm:power} is the statement, for the vectors of events, that a unitary matrix which is not a permutation matrix with phases preserves no $\ell_p$ norm other than the Euclidean one; for all vectors this is the classical description of the isometries of $\ell_p$, $p\ne2$, as permutation matrices with phases, attributed to Li and So in the finite-dimensional case \cite{liso}. That reference was not inspected for this paper and is not used in the proof.

\subsection{Gains in the roots of unity}
Let $q\ge2$ be even, $\zeta=e^{2\pi i/q}$ and $R=\mathbb{Z}[\zeta]$; a norm-preserving digraph of degree two in which two distinct vertices have a common successor needs $-1$ in the group generated by its gains, since two columns with two entries of modulus one are orthogonal only if they share both rows and the quotients of their entries in the two rows differ by the factor $-1$; this holds for every finite one \cite{g9}, and not for the binary tree. So the order is even. For $j$ prime to $q$ let $\sigma_j$ be the automorphism of $R$ with $\sigma_j(\zeta)=\zeta^j$. It commutes with complex conjugation, $\sigma_{-j}=\overline{\sigma_j}$, and $|\sigma_j(\alpha)|^2=\sigma_j(|\alpha|^2)$ is a function of $|\alpha|$. Let $J$ contain one of $j,-j$ for each such $j$; it has $\varphi(q)/2$ elements for $q>2$, with $\varphi$ Euler's function, and one for $q=2$. The character $\sigma_j\circ\chi$ is the conjugate character of \cite{g8}, Section~8.

\begin{lemma}[Parallelogram law]\label{lem:par}
Let $\Gamma$ be an abelian group and $f\colon\Gamma\to\mathbb{R}$ with $f(x+y)+f(x-y)=2f(x)+2f(y)$ for all $x,y$. Then $B(x,y)=\tfrac14\bigl(f(x+y)-f(x-y)\bigr)$ is symmetric and additive in each argument, and $f(x)=B(x,x)$.
\end{lemma}
\begin{proof}
The law gives $f(0)=0$, $f(-y)=f(y)$ and $f(2x)=4f(x)$, so $B$ is symmetric, $B(0,y)=0$ and $B(x,x)=f(x)$. Subtracting the law for the pair $(x-y,z)$ from the law for $(x+y,z)$ gives $B(x+z,y)+B(x-z,y)=2B(x,y)$. With $z=x$ this is $B(2x,y)=2B(x,y)$; with $x=u+v$, $z=u-v$ it is $B(2u,y)+B(2v,y)=2B(u+v,y)$, that is, $B(u,y)+B(v,y)=B(u+v,y)$.
\end{proof}
The lemma is the classical characterization of quadratic forms; the proof is included because the group here is a lattice, in which one cannot divide by two.

\begin{theorem}[Conjugate characters]\label{thm:conj}
Let $D$ be a gain digraph that is norm-preserving of degree two, with $\chi$ taking values in the $q$-th roots of unity, $q$ even. A \emph{block} is a pair of distinct vertices $x,x'$ with a common successor; they then have exactly two common successors $y,y'$. Suppose that for some block with vertices $x,x'$ and some starting vertex, and for every $N$, there is a length at which every $q$-th root of unity is the value of $\chi$ on the gains of at least $N$ walks to $x$ and of at least $N$ walks to $x'$. Then the rule \eqref{eq:phi} is consistent under continuation with factor $2$ if and only if there are real numbers $c_j$, $j\in J$, with
\begin{equation}\label{eq:conj}
\phi(|\alpha|)=\sum_{j\in J}c_j\,|\sigma_j(\alpha)|^2\qquad\text{for all }\alpha\in R .
\end{equation}
The consistent rules thus form a real vector space of dimension $\varphi(q)/2$ for $q>2$ and of dimension one for $q=2$; the rule is unique up to a factor exactly for $q\in\{2,4,6\}$.
\end{theorem}
\begin{proof}
\emph{The rules \eqref{eq:conj} are consistent.} For the character $\sigma_j\circ\chi$ the step matrix is $\sigma_j(A)$ and the vector of an event is $\sigma_j(v_S)$. Since $\sigma_j$ commutes with complex conjugation, $\sigma_j(A)^*\sigma_j(A)=\sigma_j(A^*A)=2$, so $\sum_y|\sigma_j(v_S(y))|^2$ is consistent with factor $2$ by Lemma~3.1 of \cite{g9}.

\emph{A consistent rule has the form \eqref{eq:conj}.} Every element of $R$ is a combination of powers of $\zeta$ with non-negative integer coefficients, because $-1=\zeta^{q/2}$. By the hypothesis there is, for any $\alpha,\beta\in R$, an event whose vector has the components $\alpha$ at $x$ and $\beta$ at $x'$ and no others. With the entries $u_1=A_{yx}$, $u_2=A_{yx'}$, $u_3=A_{y'x}$, $u_4=A_{y'x'}$ of the block, $\overline{u_1}u_2+\overline{u_3}u_4=0$, and the vector of the continued event has the components $u_1\alpha+u_2\beta$ and $u_3\alpha+u_4\beta$, of moduli $|\alpha+w\beta|$ and $|\alpha-w\beta|$ with $w=\overline{u_1}u_2$ a root of unity in $R$. Consistency, with $w^{-1}\beta$ in place of $\beta$, says that $f(\alpha)=\phi(|\alpha|)$ satisfies the parallelogram law on the additive group of $R$. By Lemma~\ref{lem:par}, $f(\alpha)=B(\alpha,\alpha)$ with $B$ symmetric and biadditive; $B$ extends to a real symmetric bilinear form on $V=R\otimes\mathbb{R}$, and $B(\zeta\alpha,\zeta\beta)=B(\alpha,\beta)$ because $f(\zeta\alpha)=f(\alpha)$. For $q=2$, $R=\mathbb{Z}$ and $f(n)=n^2f(1)$. For $q>2$ the map $\alpha\mapsto(\sigma_j(\alpha))_{j\in J}$ identifies $V$ with a sum of planes $\mathbb{C}$, on the $j$-th of which multiplication by $\zeta$ is the rotation by $2\pi j/q$. Over the complex numbers this map has the $\varphi(q)$ distinct eigenvalues $\zeta^{\pm j}$, and for eigenvectors $v,w$ with eigenvalues $\lambda,\mu$ invariance gives $B(v,w)(1-\lambda\mu)=0$; so $B$ pairs only the two eigenvectors of one plane, and on each plane it is a real multiple $c_j$ of the Euclidean inner product. Hence $f(\alpha)=\sum_jc_j|\sigma_j(\alpha)|^2$.

\emph{Dimension.} The functions $|\sigma_j(\alpha)|^2$, $j\in J$, are the squared lengths of the components of $\alpha$ in complementary planes and are linearly independent. $\varphi(q)/2=1$ for even $q>2$ exactly when $q\in\{4,6\}$.
\end{proof}

\emph{Example.} Let $B_q$ be the gain digraph on two vertices with step matrix $\bigl(\begin{smallmatrix}1&\zeta\\1&-\zeta\end{smallmatrix}\bigr)$. It is norm-preserving of degree two. The digraph of pairs (vertex, gain) has two arcs into and out of every pair, a loop at each pair with the first vertex, and is strongly connected, so the numbers of walks to each pair grow without bound for all large lengths and the hypothesis of Theorem~\ref{thm:conj} holds. For $q=2,4,6,8,10,12$ a solution of the consistency equations in exact arithmetic in $R$ returns a space of dimension $1,1,1,2,2,2$ that contains every function $|\sigma_j(\alpha)|^2$; for $q=8$ the rule of $\sigma_3$ differs from the squared modulus by more than half its value on sampled moduli \cite{comp4s}.

Theorem~\ref{thm:conj} says where consistency on histories stops. It forces the second order and then leaves exactly the freedom of Proposition~8.1 of \cite{g8}: the choice of the faithful character, that is, of the embedding of the amplitudes in the complex numbers. The orders $2,4,6$, for which that proposition makes the measure independent of the character, are the orders for which the consistent rule is unique. For other orders a combination with unequal $c_j$ is a consistent rule that is the squared modulus of no single character. Nothing in the event structure prefers $\sigma_1$. An instance on five sheets: for a fixed, independently tagged character, the three-step constructive record has the values $(3+\sqrt5)/16$ and $(3-\sqrt5)/16$ under two embeddings; from a localized sheet the same record has probability $\tfrac14$ under both (Director~G \cite{dg,dgn}; the two values for the tagged character were recomputed for this paper, and the third is an instance of \eqref{eq:lift} below).

\begin{proposition}[Positivity and composition]\label{prop:compo}
Let $\phi$ be of the form \eqref{eq:conj}.
\begin{enumerate}[label=(\alph*)]
\item $\phi(|\alpha|)\ge0$ for all $\alpha\in R$ if and only if $c_j\ge0$ for every $j\in J$.
\item Suppose the rule is multiplicative on products: for all gain digraphs with gains in the $q$-th roots of unity, among them one that satisfies the hypothesis of Theorem~\ref{thm:conj}, and events $S_1,S_2$ of the factors, $\mu_\phi(S_1\times S_2)=\mu_\phi(S_1)\,\mu_\phi(S_2)$, where the product is that of \cite{g8}, Section~6, with $v_{S_1\times S_2}=v_{S_1}\otimes v_{S_2}$. Then $\phi=0$ or there is one $j\in J$ with $\phi(|\alpha|)=|\sigma_j(\alpha)|^2$ for all $\alpha\in R$.
\end{enumerate}
\end{proposition}
\begin{proof}
(a) If all $c_j\ge0$ the sum is non-negative. Conversely the image of $R$ in the sum $V$ of the planes is a lattice of full rank, so the rays through its points are dense in $V$; if some $c_j<0$, the form $\sum_jc_j|z_j|^2$ is negative on an open cone of $V$, which contains a point of the lattice. (b) Put $N_j(\alpha)=|\sigma_j(\alpha)|^2$ and $M_j(S)=\sum_yN_j(v_S(y))$. Since $N_j$ is multiplicative, $\mu_\phi(S_1\times S_2)=\sum_jc_jM_j(S_1)M_j(S_2)$, while $\mu_\phi(S_1)\mu_\phi(S_2)=\sum_{j,k}c_jc_kM_j(S_1)M_k(S_2)$. On a carrier that satisfies the hypothesis of Theorem~\ref{thm:conj} the events of walks to one vertex have $M_j(S)=N_j(\alpha)$ with $\alpha$ arbitrary in $R$, and the functions $N_j$ are linearly independent, which means that the vectors $(N_j(\alpha))_{j\in J}$, $\alpha\in R$, span $\mathbb{R}^J$. The events $S_1$ and $S_2$ vary independently of each other, so both sides are bilinear forms that agree on a spanning set in each argument, and their coefficients agree: $c_jc_k=\delta_{jk}c_j$ for all $j,k$. Then every $c_j$ is $0$ or $1$ and at most one is $1$.
\end{proof}

Positivity confines the weights to a cone and selects nothing. Multiplicativity under composition removes the combinations and fixes the normalization: a consistent multiplicative rule is the squared modulus in one conjugate character. Part (b) agrees with a classification obtained independently and issued first by Director~G \cite{dg}: on the field $\mathbb{Q}(\zeta,\sqrt2)$, a weight that passes a marginal test with one balanced coin, is invariant under $\zeta$ and factorizes on product tables is $\lambda(a\bar a)$ for a field embedding $\lambda$ of the real subfield. The hypotheses differ (a marginal condition on records there, consistency on histories here) and the conclusions coincide where both apply. The freedom that remains is the Galois group acting on a single character, the same residue as in Proposition~7.4 of \cite{g7} and Proposition~8.1 of \cite{g8}; neither consistency nor composition breaks it. Director~G describes controls that would: a regular real plane, an order on the real algebraic amplitudes, or a calibrated preparation with a prescribed algebraic entry \cite{dg}. Instances of the proposition for $q=8,10,12$ are in \cite{comp4s}.

\subsection{The signed lift}
In Theorem~\ref{thm:conj} the end vertex of a walk is a record and its gain is not. The other reading makes the gain a record. Let $h=q/2$ and $\Lambda=\mathbb{Z}[X]/(X^h+1)$ with the basis $1,X,\dots,X^{h-1}$; multiplication by $X$ permutes the basis up to the sign $X^h=-1$. For an odd residue $j$ modulo $q$ let $\sigma_j\colon\Lambda\to\mathbb{C}$ be the ring homomorphism with $\sigma_j(X)=\zeta^j$. These are the $h$ characters of the group of $q$-th roots of unity in which $-1$ acts as $-1$; $\sigma_j$ is faithful exactly when $j$ is prime to $q$. Let $D$ be a gain digraph with gains $\zeta^k$, $0\le k<q$, and $\chi(\zeta)=\zeta$. Its \emph{signed lift} $\tilde D$ has the vertices $(y,s)$ with $y\in V$ and $s\in\{0,\dots,h-1\}$, the \emph{sheets}, and for every arc of $D$ from $x$ to $y$ with gain $\zeta^k$ and every $s$ an arc from $(x,s)$ to $(y,t)$ with gain $\varepsilon=\pm1$, where $X^{s+k}=\varepsilon X^t$ in $\Lambda$. A walk of $\tilde D$ from $(x,0)$ is a walk $w$ of $D$ from $x$ together with the record of its gain $\zeta^{k(w)}$: it ends on the sheet $t$ with gain $\varepsilon$, where $X^{k(w)}=\varepsilon X^t$. For a set $S$ of walks of $D$ of one length from $x$ put
\[
\tilde v_S(y)=\sum_{w\in S,\ t(w)=y}X^{k(w)}=\sum_{s=0}^{h-1}n_{y,s}X^s\in\Lambda .
\]
The integers $n_{y,s}$ are the components of the vector of $S$ on $\tilde D$, and $\sigma_j(\tilde v_S(y))$ is the component at $y$ of its vector on $D$ for the character $\zeta\mapsto\zeta^j$.

\begin{corollary}[Signed lift]\label{cor:lift}
Let $D$ be a finite gain digraph that is norm-preserving of degree two, with gains in the $q$-th roots of unity, $q$ even.
\begin{enumerate}[label=(\alph*)]
\item $\tilde D$ has gains $\pm1$ and is norm-preserving of degree two.
\item On each connected component of $\tilde D$ the rule \eqref{eq:phi}, with the pairs of a vertex and a sheet as end vertices and $\phi(1)\ne0$, is consistent under continuation with a factor $c>0$ if and only if $c=2$ and $\phi(n)=n^2\phi(1)$ for every integer $n\ge0$.
\item With $\phi(1)=1$ the value of that rule on a set $S$ of walks of $D$ of one length from one vertex is the integer
\begin{equation}\label{eq:lift}
\tilde\mu(S)=\sum_{y\in V}\sum_{s=0}^{h-1}n_{y,s}^2=\frac2q\sum_{j\ \mathrm{odd}}\ \sum_{y\in V}\bigl|\sigma_j(\tilde v_S(y))\bigr|^2 ,
\end{equation}
the mean, with equal weights, of the squared-modulus rules of the $h$ characters $\sigma_j$, faithful or not.
\item For $q\ge6$ and every odd $j$ the function $|\sigma_j(\cdot)|^2$ on $\Lambda$ is not a sum over the sheets of a function of $|n_s|$.
\end{enumerate}
\end{corollary}
\begin{proof}
(a) Every column of the step matrix $\tilde A$ of $\tilde D$ has two entries $\pm1$. In a block of $D$ with vertices $x,x'$, successors $y,y'$ and entries $u_i=\zeta^{k_i}$ as in the proof of Theorem~\ref{thm:conj}, the relation $\overline{u_1}u_2+\overline{u_3}u_4=0$ reads $k_2-k_1\equiv k_4-k_3+h$ modulo $q$. The columns $(x,s)$ and $(x',s')$ of $\tilde A$ have a common row over $y$ exactly when $s+k_1\equiv s'+k_2$ modulo $h$, and over $y'$ exactly when $s+k_3\equiv s'+k_4$ modulo $h$; by the relation the two conditions are the same. When they hold, the products of the two entries in the common rows are the signs $X^{(s'+k_2)-(s+k_1)}$ and $X^{(s'+k_4)-(s+k_3)}$, and their quotient is $X^{(k_2-k_1)-(k_4-k_3)}=-1$, an odd power of $X^h$. So the two columns are orthogonal. Columns over vertices of different blocks have disjoint supports, and two columns over one vertex of $D$ on different sheets also. Hence $\tilde A^{\mathsf T}\tilde A=2$. (b) A connected component of $\tilde D$ is closed under successors and predecessors, so it is a finite gain digraph with gains $\pm1$, norm-preserving of degree two, with connected underlying graph, and Theorem~\ref{thm:deg2} applies. (c) The first equality is the description of the walks of $\tilde D$ given above. For $0\le s,s'<h$ the sum of $\zeta^{j(s-s')}$ over the odd residues $j$ is $h$ if $s=s'$ and $0$ otherwise, so $\sum_{j\ \mathrm{odd}}|\sigma_j(\sum_sn_sX^s)|^2=h\sum_sn_s^2$. (d) The elements $1+X$ and $1+X^2$ have the same absolute values of coordinates up to their order, and $|\sigma_j(1+X)|^2=2+2\cos(2\pi j/q)$ differs from $|\sigma_j(1+X^2)|^2=2+2\cos(4\pi j/q)$: equality would need $q$ to divide $j$ or $3j$, which an even $q$ does not for odd $j$. For $q=4$ the element $1+X^2$ is zero and the statement fails: there $|\sigma_1(\cdot)|^2$ is the sum of the squares of the two sheet coordinates.
\end{proof}

With the sheets as records nothing is left to select. On each connected component the rule is unique up to a factor for every even order; when the lift has several components the factors of different components are independent, and (c) takes them equal. Part (c) is stated for walks lifted from the sheet $0$; multiplication by $X$ carries it to any other sheet. With that normalization the rule is unique, its values are integers, and no embedding of the amplitudes enters; for $q=2$ and $q=4$ it is the rule of Theorem~\ref{thm:conj}, and for $q\ge6$ it is none of the single-character rules. Its value \eqref{eq:lift} is the mean of the conjugate rules together with those of the characters that are not faithful; for $q=10$ the residue $5$, in which $\zeta$ acts as $-1$, carries one fifth of the weight. Proposition~4.4 of \cite{g8} identifies the coincidence count on the covering graph with the dimension-weighted mean over the characters; Corollary~\ref{cor:lift} adds that on the signed lift the mean over the characters in which $-1$ acts as $-1$ is the only consistent rule of the record-sum form. The freedom of Theorem~\ref{thm:conj} therefore belongs to the reduction by a character, in which the end vertex of $D$ is a record and the sheet is not. An event from one vertex of $\tilde D$ has an integer vector and enters every character with the weight fixed by \eqref{eq:lift}; multiplication by $X$ is an orthogonal map on the vectors of $\tilde D$ that commutes with the step and not with the sheet records, and the characters $\sigma_j$ are its sectors.

Director~G obtained first, for the signed walk with a sheet coordinate and odd $h$, that every protocol of steps and coordinate records from a localized sheet has the same rational law $2^{-t}\sum_im_i^2$ under each of the embedding laws of \cite{dg}, and that an integer branch vector is not supported in a single faithful plane when $\varphi(h)>2$ \cite{dgn}. Corollary~\ref{cor:lift}(b) and (c) reach the same law from the event structure, without assuming that the rule is one of the embedding laws; the assumption used instead is the record-sum form with the sheets as records. The instance on five sheets ($q=10$) shows both readings: $\tfrac14$ is the mean of $(3+\sqrt5)/16$ and $(3-\sqrt5)/16$, each with weight $\tfrac25$, and of $\tfrac12$ with weight $\tfrac15$. Which reading applies, that is, whether the gain of a walk is recorded at its end, is an assumption about records and is not decided here. Records of the sheet modulo a divisor of $h$ are not treated. Instances of the corollary for $q=4,6,8,10,12$, with the lifts of $B_q$ and of $25$ random carriers, are in \cite{comp4s}.

\subsection{Pair rules with a kernel}
The two preceding subsections fix what is recorded and classify the rules of the form \eqref{eq:phi}. The pair rules of \cite{g8} need no such choice. By Theorem~4.1 of \cite{g8} a pair rule that is unchanged by switching is given by a class function on the group. Let $C$ be a finite abelian group, written multiplicatively, and $f$ a real function on $C$ with $f(g^{-1})=f(g)$, the \emph{kernel}; put
\begin{equation}\label{eq:kernel}
\mu_f(S)=\sum_{\substack{w,w'\in S\\ t(w)=t(w')}}f\bigl(\varphi(w')\varphi(w)^{-1}\bigr).
\end{equation}
For a character $\chi$ of $C$ let $v^{\chi}_S$ be the vector of $S$ and $A_\chi$ the step matrix. The kernel $\operatorname{Re}\chi$ gives $\mu_f(S)=\lVert v^{\chi}_S\rVert^2$. For a gain digraph $D$ in which every vertex has $d$ successors let
\[
X(D)=\{\chi : A_\chi^*A_\chi=d\}
\]
be the set of characters for which $D$ is norm-preserving of degree $d$. It is closed under complex conjugation and under the automorphisms of the field generated by the values of the characters, since the entries of $A_\chi^*A_\chi$ are sums of values of $\chi$.

\begin{theorem}[Kernel rules]\label{thm:kernel}
Let $C$ be a finite abelian group and $D$ a gain digraph with gains in $C$ in which every vertex has $d$ successors.
\begin{enumerate}[label=(\roman*)]
\item If $f$ is a real combination of the functions $\operatorname{Re}\chi$, $\chi\in X(D)$, then $\mu_f$ is consistent under continuation with factor $d$.
\item Suppose that for every two distinct vertices $y_1,y_2$ with a common successor there is a starting vertex such that the quotients $\varphi(w_2)\varphi(w_1)^{-1}$, for walks $w_1$ to $y_1$ and $w_2$ to $y_2$ of one common length, take every value in $C$. If $\mu_f$ is not identically zero and is consistent under continuation with a factor $c>0$, then $c=d$ and $f$ is a real combination of the functions $\operatorname{Re}\chi$, $\chi\in X(D)$.
\end{enumerate}
\end{theorem}
\begin{proof}
(i) Then $\mu_f(S)=\sum_\chi a_\chi\lVert v^{\chi}_S\rVert^2$ over $\chi\in X(D)$, and each term is consistent with factor $d$ by Lemma~3.1 of \cite{g9}.

(ii) Write $f=\sum_\chi c_\chi\chi$ over all characters, with real $c_\chi=c_{\bar\chi}$. A single walk has $\mu_f=f(1)$, and its $d$ continuations end at different vertices, so $d\,f(1)=c\,f(1)$. Let $y_1\ne y_2$ have the set $Z\ne\emptyset$ of common successors, let $w_1,w_2$ be walks of one length to $y_1,y_2$ with quotient $g$, and $S=\{w_1,w_2\}$, so $\mu_f(S)=2f(1)$. Among the $2d$ walks of $\operatorname{Ext}S$ the pairs with a common end vertex are each walk with itself and, for $z\in Z$, the continuations of $w_1$ and $w_2$ to $z$, with quotient $g\,r_z$, $r_z=u_{zy_2}u_{zy_1}^{-1}$, where $u_{zy}$ is the gain of the arc from $y$ to $z$. Each such pair is counted in both orders, so
\[
\mu_f(\operatorname{Ext}S)=2d\,f(1)+2\sum_{z\in Z}f(g\,r_z)=2d\,f(1)+2\sum_\chi c_\chi\,\rho_\chi\,\chi(g),\qquad\rho_\chi=\sum_{z\in Z}\chi(r_z)=(A_\chi^*A_\chi)_{y_1y_2}.
\]
Consistency says that $\sum_\chi c_\chi\rho_\chi\chi(g)=(c-d)f(1)$ for every $g\in C$. The characters are linearly independent functions on $C$, so $c_\chi\rho_\chi=0$ for every character other than the trivial one, and $c_1\lvert Z\rvert=(c-d)f(1)$ for the trivial one. If $f(1)\ne0$, then $c=d$ and $c_1=0$. If $f(1)=0$, again $c_1=0$. In both cases $c_\chi(A_\chi^*A_\chi)_{y_1y_2}=0$ for every $\chi$ and every such pair of vertices. The diagonal entries of $A_\chi^*A_\chi$ are $d$ and its entries for two vertices without a common successor are zero; so $c_\chi\ne0$ implies $\chi\in X(D)$, and $f$ is a combination of the $\operatorname{Re}\chi$ with $\chi\in X(D)$. By (i) $\mu_f$ then multiplies by $d$ under continuation, and since it is not identically zero, $c=d$. If no two distinct vertices have a common successor, every character is in $X(D)$, two different walks from one vertex never end at one vertex, $\mu_f(S)=\lvert S\rvert f(1)$, and $c=d$.
\end{proof}

Within the pair rules the theorem leaves no freedom except the one it names: under the hypothesis on the quotients, the consistent rules are the real combinations of the squared-norm rules of the characters for which the step is norm-preserving. On a tree every character qualifies and every kernel is consistent, which is the example of Director~G after Theorem~\ref{thm:deg2}. Without the hypothesis on the quotients more kernels are consistent.

\begin{corollary}[Degree two]\label{cor:kernel2}
Let $C$ be a finite abelian group, $\epsilon\in C$ of order two, and $D$ a gain digraph with gains in $C$ whose vertices have two successors and two predecessors and fall into blocks, as in Theorem~4.3 of \cite{g9}: pairs $\{x,x'\}$ with a common pair of successors $\{y,y'\}$, the pairs of successors of different blocks being disjoint, with $u_4u_3^{-1}=\epsilon\,u_2u_1^{-1}$ in every block. Then $X(D)$ is the set of characters with $\chi(\epsilon)=-1$, and under the hypothesis of Theorem~\ref{thm:kernel}(ii) a kernel $f$ with $\mu_f$ not identically zero is consistent if and only if $f(\epsilon g)=-f(g)$ for all $g\in C$. For $C$ the group of the $q$-th roots of unity, $\epsilon=-1$ and $D$ norm-preserving of degree two for the faithful character:
\begin{enumerate}[label=(\alph*)]
\item under the hypothesis of Theorem~\ref{thm:kernel}(ii) the consistent kernels are the combinations of $\operatorname{Re}\chi_j$, $\chi_j(\zeta)=\zeta^j$, $j$ odd, one from each pair $\pm j$, a space of dimension $\lceil q/4\rceil$;
\item the rules of Theorem~\ref{thm:conj} are those with coefficient zero unless $j$ is prime to $q$;
\item the rule of Corollary~\ref{cor:lift} is the kernel with $f(1)=1$, $f(-1)=-1$ and $f=0$ elsewhere, that is, $f=\frac2q\sum_{j\ \mathrm{odd}}\chi_j$;
\item on a carrier that satisfies the hypotheses of Theorem~\ref{thm:conj} and of Theorem~\ref{thm:kernel}(ii), a consistent rule \eqref{eq:kernel} that is multiplicative on products, where gains are multiplied in $C$ as in Proposition~\ref{prop:compo}, is zero or $\lVert v^{\chi_j}_S\rVert^2$ for one odd $j$, which need not be prime to $q$.
\end{enumerate}
\end{corollary}
\begin{proof}
The entry of $A_\chi^*A_\chi$ for the two vertices of a block is $\chi(u_2u_1^{-1})(1+\chi(\epsilon))$, which vanishes exactly when $\chi(\epsilon)=-1$; the other entries off the diagonal are zero, because vertices of different blocks have no common successor. A kernel is a combination of characters with $\chi(\epsilon)=-1$ exactly when $f(\epsilon g)=-f(g)$: the coefficient of $\chi$ in $f$ is $\lvert C\rvert^{-1}\sum_gf(g)\overline{\chi(g)}$, in which the terms of $g$ and $\epsilon g$ cancel if $f(\epsilon g)=-f(g)$ and $\chi(\epsilon)=1$; conversely every character with $\chi(\epsilon)=-1$ has this symmetry. For the cyclic group the faithful character turns norm-preservation into the relation $u_4u_3^{-1}=-u_2u_1^{-1}$ in $C$. (a) With $h=q/2$, $f(\zeta^{h-k})=-f(\zeta^k)$, so $f$ is determined by its values at $\zeta^k$, $0\le k<h/2$, which are free. (b) is the definition of those rules. (c) follows from \eqref{eq:lift}. (d) is the argument of Proposition~\ref{prop:compo}(b): the functions $\lVert v^{\chi_j}_S\rVert^2$, $j$ odd, on events of walks to one vertex are the squared lengths of the components of an arbitrary element of $\Lambda$ in complementary subspaces of $\Lambda\otimes\mathbb{R}$ and are linearly independent.
\end{proof}

In degree two, then, the kernel has to see the central element $\epsilon$ as a sign and is otherwise free. Consistency on histories excludes every character in which $\epsilon$ acts trivially, the constant kernel of the coincidence count on $D$ among them, and admits every real combination of the others. For the roots of unity the rule is unique up to a factor only for $q=2$ and $q=4$; for $q=6$ there are two, the faithful one and the sign character, in which the cube roots of unity are not seen. Theorem~\ref{thm:conj} and Corollary~\ref{cor:lift} are two sections of this family: the first by the requirement that the rule be a function of the faithful amplitude, the second by the requirement that it be a sum over sheets. Strong positivity in the sense of \cite{g8} is non-negativity of the coefficients (Theorem~4.1(b) of \cite{g8}). By (d), composition leaves one character with $\chi(-1)=-1$ and does not make it faithful: that the whole group of gains acts is an assumption of Proposition~\ref{prop:compo}, contained in the form \eqref{eq:phi} with the faithful character, and is not a consequence of consistency or composition.

Three instances outside the cyclic case of degree two. For two sheet registers of prime order and a sign, the carrier of Director~G's relative transport \cite{dgr}, the group is not cyclic and has no faithful character; the corollary applies with $\epsilon$ the sign, and the consistent kernels are those odd under it. In degree three, with the block of gains $\omega^{zy}$ of \cite{g9}, $X(D)$ consists of the characters that are not trivial on $\omega$; if the gains lie in the sixth roots of unity, the characters of order three belong to it and the character of order two does not, so that $-1$ acting as $-1$ is a feature of degree two and not of consistency. In degree four with the block $i^{zy}$ and gains in the eighth roots of unity, $X(D)$ is again the set of odd characters. Computed solution spaces agree with $X(D)$ on $47$ carriers that satisfy the hypothesis, in these families and in the cyclic case of degree two, and are larger on carriers that do not \cite{comp4s}. Two predicates that required a minimum number of carriers with the hypothesis failed, in the first run for the two-register family and again in a follow-up with twelve small carriers, of which one qualified; they are kept, together with the three failed predicates of the first computation for the cyclic case, which did not test the hypothesis and used a threshold without basis. A reader working without the author's notes examined the theorem and the corollary step by step and compared solution spaces with $X(D)$ on about $2500$ carriers by a program of its own; it found no counterexample and four omissions in the wording of the corollary, which are corrected here. The recursion of Theorem~\ref{thm:quad} and Lemma~\ref{lem:par} for integer values have been checked in Lean~4 \cite{comp4s}.

\subsection{The record-sum form in degrees three and four}
Theorem~\ref{thm:power} excludes every power other than the second on any norm-preserving carrier with merging walks, and Theorems~\ref{thm:deg2} and \ref{thm:conj} treat a general function of the modulus in degree two. For two blocks of higher degree the same conclusion holds.

\begin{proposition}[Blocks of degree three and four]\label{prop:deg34}
Let $d\in\{3,4\}$, $w=e^{2\pi i/d}$, and let $D$ be a gain digraph in which every vertex has $d$ successors and the vertices fall into blocks: $d$ vertices $x_0,\dots,x_{d-1}$ with the common successors $y_0,\dots,y_{d-1}$ and gains $\rho_z\kappa_y w^{zy}$ on the arc from $x_y$ to $y_z$, with $\rho_z,\kappa_y$ and all gains in the $q$-th roots of unity, $q\in\{3,6\}$ for $d=3$ and $q=4$ for $d=4$, the sets of successors of different blocks being disjoint. Suppose that for some block and some starting vertex, and for every $N$, there is a length at which every $q$-th root of unity is the gain of at least $N$ walks to each vertex of the block. Then the rule \eqref{eq:phi} is consistent under continuation with factor $d$ if and only if $\phi(|\alpha|)=c\,|\alpha|^2$ for all $\alpha\in\mathbb{Z}[\zeta]$, with a real constant $c$.
\end{proposition}
\begin{proof}
The columns of a block are orthogonal and columns of different blocks have disjoint supports, so $A^*A=d$ and the square is consistent by Lemma~3.1 of \cite{g9}. Conversely, by the hypothesis there is for any $\beta_0,\dots,\beta_{d-1}\in R=\mathbb{Z}[\zeta]$ an event whose vector has the components $\kappa_y^{-1}\beta_y$ at $x_y$ and no others; the vector of the continued event has the moduli $|\sum_yw^{zy}\beta_y|$, $z=0,\dots,d-1$. Put $f(\alpha)=\phi(|\alpha|)$; it satisfies $f(\zeta\alpha)=f(\alpha)$, and $f(-\alpha)=f(\alpha)$.

For $d=4$ take $(\beta_y)=(a,0,b,0)$. The moduli are $|a+b|,|a-b|,|a+b|,|a-b|$, and consistency reads $2f(a+b)+2f(a-b)=4f(a)+4f(b)$, the parallelogram law on $R$.

For $d=3$ take $(\beta_y)=(a,b,b)$. Since $w+w^2=-1$ the moduli are $|a+2b|,|a-b|,|a-b|$, and consistency reads
\[
f(a+2b)+2f(a-b)=3f(a)+6f(b)\qquad\text{for all }a,b\in R .
\]
This equation for the pair $(a+b,-b)$ is $f(a-b)+2f(a+2b)=3f(a+b)+6f(b)$. Eliminating $f(a+2b)$ between the two gives $f(a+b)+f(a-b)=2f(a)+2f(b)$.

In both cases Lemma~\ref{lem:par} gives $f(\alpha)=B(\alpha,\alpha)$ with $B$ symmetric, biadditive and invariant under multiplication by $\zeta$. For $q\in\{3,4,6\}$ the space $R\otimes\mathbb{R}$ is one plane on which $\zeta$ acts as a rotation of order $q\ge3$, and an invariant symmetric form on it is a multiple of the Euclidean one, as in the proof of Theorem~\ref{thm:conj}.
\end{proof}

The elimination for $d=3$ has been checked in Lean~4 on an arbitrary abelian group, for integer values with the core library and for real values with Mathlib; the evenness of $f$ also follows from the displayed equation alone. With Mathlib the following are checked as well, for real values: Lemma~\ref{lem:par}, including additivity of $B$ and $B(nx,y)=nB(x,y)$ for integers $n$; the recursion of Theorem~\ref{thm:quad}; and the step of Theorem~\ref{thm:kernel} that a combination of characters which is constant on the group has no nontrivial character \cite{comp4s}. The graph-theoretic parts of these theorems are not formalized. On eleven carriers of these two kinds a solution of the consistency equations over sampled events, with the squared moduli computed as integer norms up to $25$, returns a one-dimensional space proportional to the norm \cite{comp4s}. Other blocks, other degrees and gains of higher order are not treated by this proposition; they are the subject of Theorems~\ref{thm:fourier} and \ref{thm:vsl} and of Corollary~\ref{cor:block}.

The proposition is the case of two small degrees and small orders of a statement that holds for every degree.

\begin{theorem}[Fourier blocks of any degree]\label{thm:fourier}
Let $d\ge2$ divide $q$, $\zeta=e^{2\pi i/q}$, $w=\zeta^{q/d}$, $R=\mathbb{Z}[\zeta]$, and let $D$ be a gain digraph in which every vertex has $d$ successors and the vertices fall into blocks: $d$ vertices $x_0,\dots,x_{d-1}$ with the common successors $y_0,\dots,y_{d-1}$ and the gain $\rho_z\kappa_yw^{zy}$ on the arc from $x_y$ to $y_z$, with $\rho_z$, $\kappa_y$ $q$-th roots of unity, the sets of successors of different blocks being disjoint. Suppose that for some block and some starting vertex, and for every $N$, there is a length at which every $q$-th root of unity is the gain of at least $N$ walks to each vertex of the block. Then the rule \eqref{eq:phi} is consistent under continuation with factor $d$ if and only if it has the form \eqref{eq:conj}: $\phi(|\alpha|)=\sum_{j\in J}c_j|\sigma_j(\alpha)|^2$ for all $\alpha\in R$, with real $c_j$. The consistent rules form a space of dimension $\varphi(q)/2$ for $q>2$ and of dimension one for $q=2$. Here $q$ may be odd; $\sigma_j$, $J$ and \eqref{eq:conj} are defined as for even $q$.
\end{theorem}
\begin{proof}
\emph{The rules \eqref{eq:conj} are consistent.} For $j$ prime to $q$, $\sigma_j(w)$ is a primitive $d$-th root of unity, so the columns of a block of $\sigma_j(A)$ are orthogonal: $\sum_z\sigma_j(w)^{z(y-y')}=0$ for $y\ne y'$. Columns of different blocks have disjoint supports. Hence $\sigma_j(A)^*\sigma_j(A)=d$ and Lemma~3.1 of \cite{g9} applies.

\emph{The parallelogram law.} As in the proof of Proposition~\ref{prop:deg34} there is for any $\beta_0,\dots,\beta_{d-1}\in R$ an event whose continuation has the moduli $|\sum_yw^{zy}\beta_y|$, and $f(\alpha)=\phi(|\alpha|)$ satisfies $f(-\alpha)=f(\alpha)=f(\zeta\alpha)$ and $f(0)=0$. If $d$ is even, $(\beta_y)$ with $\beta_0=a$, $\beta_{d/2}=b$ and the others zero gives the moduli $|a+b|$ and $|a-b|$, each $d/2$ times, and consistency is the parallelogram law. Let $d$ be odd, $d\ge3$, and $k=d-2$. Take $(\beta_y)=(a,b,\dots,b)$. Since $\sum_{y\ge1}w^{zy}=-1$ for $z\ne0$, the moduli are $|a+(d-1)b|$ once and $|a-b|$ $d-1$ times:
\[
f(a+(d-1)b)+(d-1)f(a-b)=d\,f(a)+d(d-1)f(b).
\]
The same equation for the pair $(a+kb,-b)$ is $f(a-b)+(d-1)f(a+(d-1)b)=d\,f(a+kb)+d(d-1)f(b)$. Eliminating $f(a+(d-1)b)$ and dividing by $d$ gives
\begin{equation}\label{eq:Mk}
k\,f(a-b)+f(a+kb)=(k+1)f(a)+k(k+1)f(b)\qquad\text{for all }a,b\in R .
\end{equation}
Write $M(a,b)$ for the difference of the two sides of \eqref{eq:Mk}. Then
\[
(k+1)\,M(a,b)-k\,M(a+b,b)+M(a+kb,-b)=k(k+1)\bigl(f(a+b)+f(a-b)-2f(a)-f(b)-f(-b)\bigr),
\]
because the terms $f(a+kb)$ and $f(a+(k+1)b)$ cancel and the coefficient of $f(a)$ is $(k+1)^2+k^2-1=2k(k+1)$. The left side is zero and $f(-b)=f(b)$, since $f$ depends on the modulus only, so $f(a+b)+f(a-b)=2f(a)+2f(b)$. This combination was found in the course of a formal proof; the statement that \eqref{eq:Mk} implies the parallelogram law for every $k\ge1$, without assuming that $f$ is even, is checked in Lean~4 with Mathlib \cite{comp4s}.

\emph{Conclusion.} By Lemma~\ref{lem:par}, $f(\alpha)=B(\alpha,\alpha)$ with $B$ symmetric, biadditive and invariant under multiplication by $\zeta$, and the last part of the proof of Theorem~\ref{thm:conj} applies without change. For odd $q$ the only point to add is that every element of $R$ is still a combination of powers of $\zeta$ with non-negative integer coefficients, because $-1=\zeta+\zeta^2+\dots+\zeta^{q-1}$; this is what the hypothesis needs in order to realize arbitrary $\beta_y$.
\end{proof}

For $d=2$ the block is that of Theorem~\ref{thm:conj} up to the factors $\rho,\kappa$. Thus on these carriers the degree plays no part: a consistent rule of the record-sum form is of the second order, and what remains free is the combination of conjugate characters, unique exactly for $q\in\{2,3,4,6\}$. On $33$ carriers with $d=3,4,5,6$ and $q$ up to $12$ that satisfy the hypothesis the computed solution spaces have dimension $\varphi(q)/2$ and contain the conjugate rules; in a first run two carriers whose gains generate a proper subgroup had smaller spaces, and the two predicates that did not test the hypothesis failed and are kept \cite{comp4s}. The disjointness of the successor sets cannot be dropped: on four vertices with the blocks $\{0,1\}\to\{0,2\}$ and $\{2,3\}\to\{1,2\}$ and signs of the form required, $A^*A\ne2$ and no rule other than zero is consistent. This example, and the omission of the condition from the first issue of Proposition~\ref{prop:deg34}, were found by a challenger working from the text of the paper alone. Blocks that are not of this form, such as real Hadamard blocks of degree four with other sign patterns or blocks from other groups, are not treated by this theorem; Corollaries~\ref{cor:vsl} and \ref{cor:block} cover them where their hypotheses hold.

\subsection{Any norm-preserving carrier with a merging pair}
The form of the blocks is not needed. One pair of vertices with a common successor suffices, on a norm-preserving carrier of any degree with gains in the roots of unity, provided the events at that pair are rich enough.

\begin{theorem}[Vanishing sums]\label{thm:vsl}
Let $q\ge2$, $\zeta=e^{2\pi i/q}$, $R=\mathbb{Z}[\zeta]$, and let $f\colon R\to\mathbb{R}$ depend only on $|\alpha|$. Let $u_1,\dots,u_m$, $m\ge2$, be roots of unity in $R$ with $u_1+\dots+u_m=0$, and suppose
\begin{equation}\label{eq:V}
\sum_{z=1}^{m}f(a+u_zb)=m\,f(a)+m\,f(b)\qquad\text{for all }a,b\in R .
\end{equation}
Then $f(\alpha)=\sum_{j\in J}c_j|\sigma_j(\alpha)|^2$ with real $c_j$; for $q=2$, $f(n)=n^2f(1)$. Conversely every such $f$ satisfies \eqref{eq:V}. Here $q$ may be odd, with $\sigma_j$ and $J$ defined as for even $q$, and the $u_z$ may be any roots of unity of $R$, for odd $q$ also the negatives of the $q$-th roots.
\end{theorem}
\begin{proof}
Put $n=\varphi(q)$ and $V=R\otimes\mathbb{R}$, identified with $\mathbb{C}^{n/2}$ by $x=(\sigma_j)_{j\in J}$ for $q>2$; multiplication in $R$ is componentwise multiplication in $V$. For $q=2$, $R=\mathbb{Z}$, the $u_z$ are $m/2$ times $1$ and $m/2$ times $-1$, \eqref{eq:V} is the parallelogram law, and the claim follows as in Theorem~\ref{thm:quad}. Let $q>2$.

\emph{Step 1: $f$ is a polynomial.} For $h\in R$ let $X^h$ be the shift $(X^hf)(a)=f(a+h)$; the shifts span the group ring $\mathbb{C}[R]$, a ring of Laurent polynomials in $n$ variables. For fixed $b=\xi$ the difference of the two sides of \eqref{eq:V} is the function $(\sum_zX^{\xi u_z}-m)f$ minus a constant, so
\[
g_{h,\xi}=(X^h-1)\Bigl(\sum_zX^{\xi u_z}-m\Bigr)
\]
annihilates $f$ for all $h,\xi\in R$. A common zero of all $g_{h,\xi}$ is a homomorphism $\chi\colon R\to\mathbb{C}^*$. If $\chi\ne1$, take $h$ with $\chi(h)\ne1$; then $\sum_z\chi(\xi u_z)=m$ for every $\xi$. Replacing $\xi$ by $k\xi$, $k=1,\dots,m$, the numbers $w_z=\chi(\xi u_z)$ have the power sums of $(1,\dots,1)$, so $w_z=1$ for all $z$ by Newton's identities. Then $\chi(\xi u_1)=1$ for all $\xi$, and $u_1$ is a unit, so $\chi=1$, a contradiction. The ideal generated by the $g_{h,\xi}$ therefore has the single zero $\chi=1$; by the Nullstellensatz it contains a power $\mathfrak{a}^K$ of the augmentation ideal $\mathfrak{a}$, which is generated by the $X^h-1$. So every $K$-fold difference of $f$ vanishes, and $f$ is the restriction to $R$ of a real polynomial $P$ on $V$.

\emph{Step 2: $P$ is a polynomial in the $t_j=|x_j|^2$.} For $\beta,c\in R$, $|\beta c|=|\bar\beta c|$, so $P(\beta c)=P(\bar\beta c)$ on $R\times R$, hence on $V\times V$, because a lattice of full rank is Zariski dense. With $\beta=(e^{i\theta_j})_j$ and $e^{i\theta}c$ in place of $c$ this is $P(e^{2i\theta}c)=P(c)$: $P$ is invariant under the torus, and the invariant polynomials are those in the $t_j$.

\emph{Step 3: $P$ is linear in the $t_j$.} Equation \eqref{eq:V} is now a polynomial identity on $V\times V$, in which $u_z$ acts on the $j$-th coordinate as $\sigma_j(u_z)$. Let $P_d$ be the homogeneous part of degree $d\ge4$ in the real coordinates. The part of \eqref{eq:V} of degree $d-2$ in $a$ and $2$ in $b$ has no term on the right. On the left it is $\tfrac12\sum_zD^2P_d(a)[u_zb,u_zb]$, in which the monomial $b_j\bar b_j$ has the coefficient $\sum_z|\sigma_j(u_z)|^2\,\partial_{x_j}\partial_{\bar x_j}P_d=m\,\partial_{x_j}\partial_{\bar x_j}P_d$. So $\partial_{x_j}\partial_{\bar x_j}P_d=0$ for every $j$. For $P_d=g(t)$ this reads $\partial_{t_j}(t_j\,\partial_{t_j}g)=0$, so $t_j\partial_{t_j}g$ does not depend on $t_j$ and, being divisible by $t_j$, is zero. Hence $g$ is constant, and $P_d=0$. There are no parts of odd degree, and $P(0)=0$ by \eqref{eq:V} at $a=b=0$. So $P=\sum_jc_jt_j$.

\emph{Converse.} $\sum_z|\sigma_j(a+u_zb)|^2=m|\sigma_j(a)|^2+m|\sigma_j(b)|^2+2\operatorname{Re}\bigl(\overline{\sigma_j(a)}\sigma_j(b)\,\sigma_j(\textstyle\sum_zu_z)\bigr)$, and the last term is zero.
\end{proof}

\begin{corollary}[Carriers with a merging pair]\label{cor:vsl}
Let $D$ be a gain digraph with $\chi$ taking values in the $q$-th roots of unity and $A^*A=d$. Let $y_1\ne y_2$ be vertices with a common successor, and suppose that for some starting vertex and every $N$ there is a length at which every $q$-th root of unity is the value of $\chi$ on the gains of at least $N$ walks to $y_1$ and of at least $N$ walks to $y_2$. Then the rule \eqref{eq:phi} with $\phi(1)\ne0$ is consistent under continuation with a factor $c>0$ if and only if $c=d$ and it has the form \eqref{eq:conj}.
\end{corollary}
\begin{proof}
The rules \eqref{eq:conj} are consistent with factor $d$ because $\sigma_j(A)^*\sigma_j(A)=\sigma_j(A^*A)=d$. Conversely, a single walk has measure $\phi(1)$ and its $d$ continuations end at different vertices, so $c=d$. Let $Z$ be the set of the $m$ common successors of $y_1,y_2$. For $a,b\in R$ there is an event with the components $a$ at $y_1$ and $b$ at $y_2$ and no others. Its continuation has the modulus $|a|$ at the $d-m$ other successors of $y_1$, $|b|$ at those of $y_2$, and $|aA_{zy_1}+bA_{zy_2}|=|a+u_zb|$ at $z\in Z$, with $u_z=\overline{A_{zy_1}}A_{zy_2}$ and $\sum_zu_z=0$ by the orthogonality of the two columns; so $m\ge2$. Consistency is \eqref{eq:V} for $f(\alpha)=\phi(|\alpha|)$, and Theorem~\ref{thm:vsl} applies.
\end{proof}

Where its hypothesis holds, Corollary~\ref{cor:vsl} gives the conclusions of Theorems~\ref{thm:quad}, \ref{thm:deg2}, \ref{thm:conj} and \ref{thm:fourier} and of Proposition~\ref{prop:deg34}, and extends Theorem~\ref{thm:power} from powers to arbitrary functions of the modulus: on a norm-preserving carrier with gains in the roots of unity, of any degree and with blocks of any form, in which two vertices have a common successor and every gain reaches both in abundance, a consistent rule of the record-sum form is of the second order, and its only freedom is the combination of conjugate characters. It does not replace those theorems. Theorem~\ref{thm:deg2} needs no abundance of gains, which it obtains from finiteness and connectedness, and Theorem~\ref{thm:power} needs neither abundance nor roots of unity; the corollary's hypothesis does follow from that of Theorem~\ref{thm:deg2}: such a carrier is strongly connected, the gains of its closed walks generate $\{\pm1\}$ since otherwise a switching would make all gains $+1$, which $A^*A=2$ with a merging pair excludes, so $\hat D$ is strongly connected; the subgroup $C_0$ of Proposition~\ref{prop:abund} is not trivial by the first step of the proof of Corollary~\ref{cor:abund}, hence is the whole group, and part (b) of the proposition gives both gains in abundance at both vertices of a block. This argument is due to the review queue, which also found it to hold on $234$ of $234$ carriers. It does not follow from the hypothesis of Theorem~\ref{thm:power}, which has no roots of unity. The earlier proofs are kept because they are elementary and give the parallelogram law by explicit instances; the present one uses the Nullstellensatz and gives no bound on $K$. The hypothesis that $f$ depends on the modulus cannot be weakened to invariance under units, signs and conjugation: $f(a)=\operatorname{Re}(a^4)$ on $\mathbb{Z}[i]$ has these invariances and satisfies \eqref{eq:V} for $(u_z)=(1,i,-1,-i)$, and is not a function of $|a|$ ($f(5)=625$, $f(3+4i)=-527$). The theorem, the example and a number of explicit certificates for special $(u_z)$ were found with a sub-agent; exact solutions of \eqref{eq:V} on boxes of $R$ for $q=3,4,5,6,8,10,12$ and $21$ vanishing multisets have dimension $\varphi(q)/2$ once the box is large enough, and three runs on boxes that were too small failed and are kept \cite{comp4s}.

\subsection{When the gains are abundant}
The hypothesis of Corollary~\ref{cor:vsl} on the abundance of gains can be decided from the carrier. Let $D$ be finite with gains in a finite abelian group $C$, and let $\hat D$ be the digraph on $V\times C$ with an arc from $(x,g)$ to $(y,g\varphi(e))$ for every arc $e$ from $x$ to $y$; a walk of $D$ from $x$ with gain $g$ is a walk of $\hat D$ from $(x,1)$ to $(t(w),g)$. The period of a strongly connected digraph is the greatest common divisor of the lengths of its closed walks.

\begin{proposition}[Abundance]\label{prop:abund}
Let $D$ be a finite gain digraph with gains in a finite abelian group $C$ and $A^*A=d$, $d\ge2$, for some character, and let $\hat D$ be strongly connected, with period $r$; let $r_0$ be the period of $D$.
\begin{enumerate}[label=(\alph*)]
\item Every vertex of $D$ has $d$ successors and $d$ predecessors. $r_0$ divides $r$, and there is a homomorphism $\pi\colon C\to\mathbb{Z}/r$ with image $r_0\mathbb{Z}/r$ such that $(x,gh)$ lies $\pi(h)$ cyclic classes of $\hat D$ after $(x,g)$. Put $C_0=\ker\pi$; its index is $r/r_0$.
\item Let $y_1\ne y_2$ have a common successor $z$, and $u=\varphi(y_1\to z)\,\varphi(y_2\to z)^{-1}$. For every starting vertex, every coset $K$ of $C_0$ and every $N$ there is a length at which the gains of the walks to $y_1$ are exactly the elements of $K$, each at least $N$ times, and those of the walks to $y_2$ exactly the elements of $Ku$, each at least $N$ times. At no length do gains of two cosets occur at one vertex.
\item Every element of $C$ is the gain of walks of one common length to $y_1$, and likewise to $y_2$, if and only if $r=r_0$; in particular if $\hat D$ is aperiodic. For a faithful character this is the hypothesis of Corollary~\ref{cor:vsl}; for a character that is not faithful that hypothesis, which concerns the values of the character, can hold although $r\ne r_0$.
\end{enumerate}
\end{proposition}
\begin{proof}
(a) $D$ is strongly connected, being the image of $\hat D$, so $r_0$ is defined; a closed walk of $\hat D$ projects to a closed walk of $D$ of the same length, so $r_0$ divides $r$. $A/\sqrt d$ is a square matrix with orthonormal columns, hence unitary, so the rows of $A$ have squared length $d$ as well: every vertex has $d$ predecessors. The same holds in $\hat D$. The group $C$ acts on $\hat D$ by the automorphisms $(x,g)\mapsto(x,gh)$, which permute the $r$ cyclic classes cyclically by a shift $\pi(h)$ that does not depend on the vertex, because an automorphism maps walks to walks; $\pi$ is a homomorphism. A closed walk of $D$ at $x$ of length $\ell$ with gain $h$ is a walk of $\hat D$ from $(x,1)$ to $(x,h)$, so $\pi(h)\equiv\ell$ modulo $r$; the lengths of closed walks generate $r_0\mathbb{Z}$ and their gains generate $C$, since $\hat D$ is strongly connected. Hence the image of $\pi$ is $r_0\mathbb{Z}/r$.

(b) $(y_1,g)$ and $(y_2,gu)$ have the common successor $(z,g\varphi(y_1\to z))$, so they lie in one cyclic class. The vertices over $y_1$ in the class of $(y_1,g)$ are the $(y_1,gh)$ with $h\in C_0$. The $r$-th power of the arc-count matrix of $\hat D$ is primitive on each cyclic class, so for large $L$ in the right residue class modulo $r$ every vertex of the class of $(y_1,g)$ is reached from $(x_0,1)$ by walks of length $L$, and by none of another class; since every vertex has $d\ge2$ predecessors the numbers of walks grow without bound.

(c) By (b) all gains occur at $y_1$ at one length exactly when $C_0=C$.
\end{proof}

\begin{corollary}\label{cor:abund}
Under the hypotheses of Proposition~\ref{prop:abund}, with $C$ the $q$-th roots of unity and the faithful character, let $y_1\ne y_2$ have a common successor and let $q_0$ be the order of $C_0$. Then $q_0\ge2$, the amplitudes at one vertex of all events of one length from one starting vertex lie in $g\,\mathbb{Z}[C_0]$ for some $g\in C$, and the rule \eqref{eq:phi} with $\phi(1)\ne0$ is consistent under continuation with a factor $c>0$ if and only if $c=d$ and $\phi(|\alpha|)=\sum_jc_j|\sigma_j(\alpha)|^2$ for $\alpha\in\mathbb{Z}[C_0]$, the sum over the conjugate characters of the $q_0$-th roots of unity. The space of consistent rules, as functions on the moduli that occur, has dimension $\varphi(q_0)/2$ for $q_0>2$ and one for $q_0=2$.
\end{corollary}
\begin{proof}
For a second common successor $z'$ the vertices $(z',g\varphi(y_1\to z'))$ and $(z',gu\varphi(y_2\to z'))$ lie in one class, so $u\,u_{z'}\in C_0$ with $u_{z'}=\overline{A_{z'y_1}}A_{z'y_2}$; the $u\,u_{z'}$ sum to zero, so $C_0\ne1$. By Proposition~\ref{prop:abund}(b), for $a',b'\in\mathbb{Z}[C_0]$ there is an event with the components $ga'$ at $y_1$ and $gub'$ at $y_2$ and no others: every element of $\mathbb{Z}[C_0]$ is a combination of elements of $C_0$ with non-negative integer coefficients, because the elements of a nontrivial finite subgroup of the roots of unity sum to zero. Its continuation has the moduli $|a'+u\,u_{z'}b'|$ at the common successors. Since $\phi$ depends on the modulus only, consistency is \eqref{eq:V} on $\mathbb{Z}[C_0]$ for the multiset $(u\,u_{z'})$ of elements of $C_0$, and Theorem~\ref{thm:vsl} applies with $q_0$ in place of $q$. The factor is $d$ as in Corollary~\ref{cor:vsl}, and the converse holds because the rules of the $q_0$-th roots are restrictions of rules \eqref{eq:conj}.
\end{proof}

So on a finite carrier with $A^*A=d$, $d\ge2$, for the faithful character, with a merging pair and with $\hat D$ strongly connected, the classification needs no hypothesis on events: the order that matters is $q_0$, the order of the gains that can meet at one length, which equals $q$ exactly when $D$ and $\hat D$ have the same period. This accounts for carriers whose solution spaces are smaller than $\varphi(q)/2$ although their gains generate the whole group; where $D$ is strongly connected and the gains of closed walks generate a proper subgroup, $\hat D$ is not strongly connected and the proposition applies, after a switching that brings all gains into that subgroup, with the subgroup in place of $C$. The naive statement, that every gain occurs at both vertices at one length whenever $\hat D$ is strongly connected, is false: on the four vertices with arcs and exponents $(0,2,0)$, $(0,3,1)$, $(1,0,1)$, $(1,1,3)$, $(2,0,0)$, $(2,1,0)$, $(3,2,0)$, $(3,3,3)$ for $q=4$, $D$ is aperiodic, $\hat D$ is strongly connected of period $2$, and at any length only two of the four gains occur at a vertex. The proposition, its proof and this example were found with a sub-agent after the director's naive version failed; on $660$ random carriers of degrees two to four the corrected statement holds wherever its hypotheses do ($483$), and the naive one fails on exactly the $130$ with $r\ne r_0$ \cite{comp4s}. A first version of the program had the coset of $y_2$ inverted and failed two predicates; it is kept. A challenger found that part (c) as first drafted, for an arbitrary character, was false, by an example with the group $\mathbb{Z}_2\times\mathbb{Z}_2$; on $473$ further carriers it found the dimension $\varphi(q_0)/2$, among them $110$ with $q_0<q$.

\subsection{Without the modulus: the block law}
So far the rule adds over end vertices a function of the \emph{modulus} of the summed gains. On a carrier with full blocks this restriction is not needed: any function of the summed gains will do, and what consistency then forces is a quadratic form.

\begin{theorem}[Block law]\label{thm:block}
In this theorem $V=R\otimes\mathbb{R}$, a real vector space of dimension $\varphi(q)$ that contains $R$ as a lattice. Let $q\ge2$, $d\ge2$, and let $B$ be a $d\times d$ matrix whose entries are $q$-th roots of unity, with $B^*B=d$. Let $C_B$ be the group generated by the entries of $B$, of order $q_B$. For a function $\Phi\colon R\to\mathbb{R}$ the following are equivalent.
\begin{enumerate}
\item[(i)] For all $\beta\in R^d$,
\begin{equation}\label{eq:BL}
\sum_{z}\Phi\Bigl(\sum_yB_{zy}\beta_y\Bigr)=d\sum_y\Phi(\beta_y).
\end{equation}
\item[(ii)] $\Phi$ is the restriction to $R$ of a real quadratic form $Q$ on $V$ with $Q(ux)=Q(x)$ for every entry $u$ of $B$, equivalently for every $u\in C_B$.
\end{enumerate}
No invariance of $\Phi$ and no condition $\Phi(0)=0$ is assumed in (i). For $q>2$, in the coordinates $x_j=\sigma_j$, $j\in J$, the forms in (ii) are the real combinations of $|x_j|^2$, of the real and imaginary parts of $x_j\bar x_k$ for $j<k$ with $j\equiv k$ modulo $q_B$, and of the real and imaginary parts of $x_jx_k$ for $j\le k$ with $j+k\equiv0$ modulo $q_B$. In particular, if $\Phi$ satisfies \eqref{eq:BL} and $\Phi(u\alpha)=\Phi(\alpha)$ for all $q$-th roots of unity $u$, or if $C_B$ contains all of them, then $\Phi(\alpha)=\sum_{j\in J}c_j|\sigma_j(\alpha)|^2$ for $q>2$ and $\Phi(\alpha)=c\,\alpha^2$ for $q=2$.
\end{theorem}

\begin{proof}
Multiplication by a root of unity $u$ of $R$ is an orthogonal map $\rho(u)$ of $V$ for the inner product $\sum_j\operatorname{Re}(x_j\bar y_j)$. Let $T\colon V^d\to V^d$, $(T\beta)_z=\sum_yB_{zy}\beta_y$. Since $\sigma_j(B)^*\sigma_j(B)=d$ for every $j$, $T=\sqrt d\,U$ with $U$ orthogonal, the block $(z,y)$ of $U$ being $\rho(B_{zy})/\sqrt d$; and $T$ maps $R^d$ into itself, as does $T^*$, the map of the matrix $B^*$, with $TT^*=d$ because $B$ is square. Write $F(\beta)=\sum_y\Phi(\beta_y)$, so that \eqref{eq:BL} is $F(T\beta)=dF(\beta)$.

\emph{(ii) implies (i).} Let $Q(x)=\langle x,Mx\rangle$ with $M$ symmetric. Invariance under $u$ is $\rho(u)^{\mathsf T}M\rho(u)=M$, that is, $M\rho(u)=\rho(u)M$. Then $1_d\otimes M$ commutes with $U$, so $F(T\beta)=d\,\langle U\beta,(1\otimes M)U\beta\rangle=d\,\langle\beta,(1\otimes M)\beta\rangle=dF(\beta)$.

\emph{Step 1: on every coset of $dR$, $\Phi$ is a polynomial of degree at most two.} Write $\Delta_hf(a)=f(a+h)-f(a)$. For $h,h'\in R$ apply to both sides of $F(T\beta)=dF(\beta)$, as functions of $\beta$, the mixed difference with the steps $h$ in $\beta_1$ and $h'$ in $\beta_2$. It annihilates $F$, which is a sum of functions of one $\beta_y$ each, and it changes $(T\beta)_z$ by $B_{z1}h$ and $B_{z2}h'$. Hence
\[
\sum_z\bigl(\Delta_{B_{z1}h}\Delta_{B_{z2}h'}\Phi\bigr)\bigl((T\beta)_z\bigr)=0\qquad(\beta\in R^d).
\]
Replace $\beta$ by $\beta+T^*(ke_{z_0})$, $k\in R$: this changes $T\beta$ by $dk$ in the component $z_0$ only. Subtracting, $\bigl(\Delta_{dk}\Delta_{B_{z_01}h}\Delta_{B_{z_02}h'}\Phi\bigr)((T\beta)_{z_0})=0$. Here $(T\beta)_{z_0}$ takes every value in $R$ (let $\beta=\overline{B_{z_01}}\,a\,e_1$), and $B_{z_01}h$, $B_{z_02}h'$ are arbitrary elements of $R$. So
\begin{equation}\label{eq:third}
\Delta_{dk}\Delta_h\Delta_{h'}\Phi=0\qquad\text{on }R,\ \text{for all }k,h,h'\in R .
\end{equation}
In particular all threefold differences of $\Phi$ with steps in $dR$ vanish, so on each coset $c+dR$, a lattice of full rank in $V$, $\Phi$ agrees with a unique real polynomial $P_c$ on $V$ of degree at most two: if all threefold differences of a function on a lattice vanish, its twofold differences are constant, its differences are affine, and it is a polynomial of total degree at most two. Write $P_c=P_c^{2}+P_c^{1}+P_c^{0}$ for its homogeneous parts with respect to the origin of $V$, with $P^2_c(x)=\langle x,M_cx\rangle$, $P^1_c(x)=\langle\ell_c,x\rangle$, $P^0_c=\kappa_c$; there are finitely many cosets.

\emph{Step 2: a polynomial identity for every vector of cosets.} Let $c=(c_y)\in(R/dR)^d$ and $c'=Tc$. For $\beta$ in the coset $c$ of $dR^d$, $T\beta$ lies in the coset $c'$ and \eqref{eq:BL} reads $\sum_zP_{c'_z}((T\beta)_z)=d\sum_yP_{c_y}(\beta_y)$. Both sides are polynomials in $\beta\in V^d$ and a coset of a lattice of full rank is Zariski dense, so this holds on $V^d$, and for each degree separately:
\begin{equation}\label{eq:deg}
\sum_zP^i_{c'_z}((T\beta)_z)=d\sum_yP^i_{c_y}(\beta_y)\qquad(\beta\in V^d,\ i=0,1,2).
\end{equation}

\emph{Step 3: the parts of degree zero and one vanish.} Let $\mu=\max_c|\kappa_c|$, attained at $c^*$, and take $c=(c^*,\dots,c^*)$ in \eqref{eq:deg} with $i=0$: the right side has absolute value $d^2\mu$ and the left at most $d\mu$; so $\mu=0$. For $i=1$, \eqref{eq:deg} says $T^{\mathsf T}\ell'=d\,\ell$ with $\ell=(\ell_{c_y})_y$ and $\ell'=(\ell_{c'_z})_z$; as $T^{\mathsf T}=\sqrt d\,U^{\mathsf T}$, $|\ell'|=\sqrt d\,|\ell|$. With $\mu=\max_c|\ell_c|$ attained at $c^*$ and $c=(c^*,\dots,c^*)$, $|\ell|=\sqrt d\,\mu$ and $|\ell'|\le\sqrt d\,\mu$, so $d\mu\le\sqrt d\mu$ and $\mu=0$.

\emph{Step 4: the parts of degree two coincide and are invariant.} For $i=2$, \eqref{eq:deg} is $U^{\mathsf T}D'U=D$ with the block diagonal matrices $D=\operatorname{diag}(M_{c_y})$ and $D'=\operatorname{diag}(M_{c'_z})$, that is, $D'=UDU^{\mathsf T}$. For $c=0$ this gives $(1\otimes M_0)U=U(1\otimes M_0)$, whose block $(z,y)$ is $M_0\rho(B_{zy})=\rho(B_{zy})M_0$: the form $P^2_0$ is invariant under every entry of $B$. For $c=c_0e_1$ subtract the case $c=0$: $D'-1\otimes M_0=U\operatorname{diag}(M_{c_0}-M_0,0,\dots,0)U^{\mathsf T}$. The left side is block diagonal, and the block $(z,z')$ of the right side is $\tfrac1d\rho(B_{z1})(M_{c_0}-M_0)\rho(B_{z'1})^{\mathsf T}$. Taking $z\ne z'$, which is possible because $d\ge2$, and using that the $\rho$ are invertible, $M_{c_0}=M_0$ for every $c_0$.

Hence $\Phi(a)=\langle a,M_0a\rangle$ for every $a\in R$, which is (ii). The list of invariant forms follows because the monomials $x_j\bar x_k$ and $x_jx_k$ are multiplied by $\zeta^{t(j-k)}$ and $\zeta^{t(j+k)}$ under $u=\zeta^t$; for $C_B$ the full group and $q>2$, $j\equiv k$ forces $j=k$ in $J$ and $j+k\equiv0$ is impossible in $J$, since $J$ contains one of $j,-j$ and $2j\not\equiv0$.
\end{proof}


\begin{corollary}[Carriers with a full block]\label{cor:block}
Let $D$ be a gain digraph with gains in the $q$-th roots of unity and the faithful character, with a block: $d\ge2$ vertices $x_0,\dots,x_{d-1}$ whose sets of successors are one set $y_0,\dots,y_{d-1}$, the matrix $B=(A_{y_zx_y})$ satisfying $B^*B=d$. Suppose that for some starting vertex and every $N$ there is a length at which every $q$-th root of unity is the gain of at least $N$ walks to each $x_y$. Let $\Phi\colon R\to\mathbb{R}$ be any function with $\Phi(0)=0$ and $\mu(S)=\sum_y\Phi(v_S(y))$. If $\mu$ is consistent under continuation with factor $d$, then $\Phi$ is a real quadratic form invariant under the entries of $B$. If moreover $\mu$ is unchanged by switching, that is, $\Phi(u\alpha)=\Phi(\alpha)$ for every $q$-th root of unity $u$, then $\Phi(\alpha)=\sum_{j\in J}c_j|\sigma_j(\alpha)|^2$; and on a carrier with $A^*A=d$ every such rule is consistent.
\end{corollary}
\begin{proof}
By the hypothesis every $\beta\in R^d$ is the vector of an event with the components $\beta_y$ at $x_y$ and no others, and its continuation has the components $\sum_yB_{zy}\beta_y$ at $y_z$ and no others. Consistency for these events is \eqref{eq:BL}, and Theorem~\ref{thm:block} applies. The last statement is Lemma~3.1 of \cite{g9} for each $\sigma_j$.
\end{proof}

Every finite norm-preserving digraph of degree two consists of blocks \cite{g9}, so for these carriers, where the gains are abundant in the sense of Proposition~\ref{prop:abund}, the corollary removes the modulus from the assumptions: a consistent rule that adds over end vertices \emph{any} function of the summed gains vanishing at zero, and is unchanged by switching, is a combination of the conjugate squared moduli; for the block of $B_q$ the entries generate all $q$-th roots of unity and switching invariance is not even needed. The proof is elementary. It uses events with all $d$ inputs: with two inputs only the conclusion can fail, as it does for $d=q=3$: for the Fourier block with $d=q=3$ the function $\operatorname{Re}(c\,a^3)$ satisfies every two-input instance of \eqref{eq:BL} and is invariant under the cube roots of unity, while it fails with three inputs (the coefficient of $\beta_0\beta_1\beta_2$ in $\sum_z(\beta_0+\omega^z\beta_1+\omega^{2z}\beta_2)^3$ is $18$). Invariance under all roots cannot be dropped when the entries of $B$ generate a proper subgroup: for a real Hadamard block and $q=4$ every real quadratic form on $\mathbb{C}$ satisfies \eqref{eq:BL}, for instance $\operatorname{Re}(a^2)$. For $d=1$ only invariance under the single entry is forced. The condition $\Phi(0)=0$ of the corollary is needed because the rule sums over all vertices; on a finite carrier it follows from consistency for an event whose vector is zero. The theorem and its proof were found by a sub-agent that had been asked a narrower question; each step was rechecked for this paper. Two steps of linear algebra used in it are checked in Lean~4 with Mathlib; exact solutions of \eqref{eq:BL} on boxes for twelve blocks have the dimension of the invariant forms, and homogeneous polynomial solutions of degrees one to six for fourteen blocks exist only in degree two \cite{comp4s}. Four predicates that expected the two-input equations to suffice failed, as they should, and are kept.

\subsection{Without the sum: the block law on multisets}
In Theorem~\ref{thm:block} the rule is a function of the summed gains. Here it may be any function of the \emph{multiset} of gains. Let $C$ be a finite abelian group, written multiplicatively, and $\hat C$ its dual group, the group of characters $\psi\colon C\to\mathbb{C}^\times$ under pointwise multiplication, with the trivial character $1$ and $\psi^{-1}=\overline\psi$; $|\hat C|=|C|$. If $C$ is cyclic of order $q$ with generator $g$ and $\chi(g)=\zeta=e^{2\pi i/q}$, then $\hat C=\{\chi^j:j\in\mathbb{Z}/q\}$, and we write $j$ for $\chi^j$. Let $M=\mathbb{N}[C]$ be the free commutative monoid on $C$: its elements are the count vectors $m=(n_h)_{h\in C}$, $n_h\in\mathbb{N}=\{0,1,2,\dots\}$, with $|m|=\sum_hn_h$, and $C$ acts by translation, $(k\cdot m)_h=n_{k^{-1}h}$. For $\psi\in\hat C$ let $\sigma_\psi(m)=\sum_hn_h\psi(h)$ (for $C$ cyclic and $\psi=\chi^j$ with $j$ prime to $q$ this is the automorphism $\sigma_j$ of Section~\ref{sec:mod} applied to the sum of the gains; here every character is allowed), so that $\sigma_1(m)=|m|$, $\sigma_{\psi^{-1}}=\overline{\sigma_\psi}$ on $M$ and $\sigma_\psi(k\cdot m)=\psi(k)\sigma_\psi(m)$. For a $d\times d$ matrix $B$ with entries in $C$ write $B^{(\psi)}$ for the complex matrix $(\psi(B_{zy}))_{z,y}$, and
\[
J_0=\{\psi\in\hat C:\ B^{(\psi)}\text{ is the all-ones matrix}\},\qquad
\mathcal{U}=\{\psi\in\hat C:\ B^{(\psi)*}B^{(\psi)}=d\}.
\]
$J_0$ is a subgroup of $\hat C$, the annihilator of the subgroup $C_B$ of $C$ generated by the entries of $B$, so $|J_0|=|C|/|C_B|$; $\mathcal{U}=\mathcal{U}^{-1}$, $\mathcal{U}J_0=\mathcal{U}$, and $1\notin \mathcal{U}$ when $d\ge2$. Whether $\psi$ belongs to $\mathcal{U}$ depends only on the kernel of $\psi$ (Corollary~\ref{cor:multiset}).

\begin{theorem}[Block law on multisets]\label{thm:multiset}
Let $C$ be a finite abelian group, $d\ge2$, and let $B$ be any $d\times d$ matrix with entries in $C$; no condition $B^*B=d$ is assumed. For a function $F\colon M\to\mathbb{R}$ the following are equivalent.
\begin{enumerate}
\item[(i)] For all $(m_1,\dots,m_d)\in M^d$,
\begin{equation}\label{eq:MBL}
\sum_{z=1}^dF\Bigl(\sum_{y=1}^dB_{zy}\cdot m_y\Bigr)=d\sum_{y=1}^dF(m_y).
\end{equation}
\item[(ii)] $F=L+Q$ on $M$, where $L(m)=\sum_ha_hn_h$ is a real linear form with $\sum_za_{B_{zy}h}=d\,a_h$ for all $y$ and $h$, and $Q$ is a real quadratic form on $\mathbb{R}^C$, homogeneous of degree two in the counts, that satisfies \eqref{eq:MBL} for all $(m_y)\in(\mathbb{R}^C)^d$.
\end{enumerate}
In (ii), $L$ and $Q$ are determined by $F$. The forms $L$ are exactly the real-valued combinations $\sum_{\psi\in J_0}l_\psi\sigma_\psi$ with complex $l_\psi$, and the forms $Q$ are exactly the real-valued combinations
\[
\sum c_{\psi\varphi}\,\sigma_\psi\overline{\sigma_\varphi}\qquad\text{over pairs }(\psi,\varphi)\in\hat C^2\text{ with }\psi\in \mathcal{U}\text{ and }\psi\varphi^{-1}\in J_0,
\]
with complex $c_{\psi\varphi}$; the pairs $(\psi,\varphi)$ and $(\varphi^{-1},\psi^{-1})$ give the same product, so these products span and are not a basis. The real vector space of solutions of \eqref{eq:MBL} has dimension
\[
|J_0|+\#\bigl\{\{\psi,\lambda\}:\ \psi,\lambda\in\hat C,\ \psi\lambda\in J_0,\ \psi\in \mathcal{U}\bigr\},
\]
the second term counting unordered pairs, $\psi=\lambda$ allowed. Every solution has $F(0)=0$ and is determined by its values on the multisets of size one and two. The same statements hold for complex-valued $F$, with complex forms and complex dimension.
\end{theorem}

\begin{corollary}\label{cor:multiset}
Let $d\ge2$ and let the entries of $B$ generate $C$. Then the solutions $F\colon M\to\mathbb{R}$ of \eqref{eq:MBL} are exactly
\[
F(m)=a\,|m|+\sum_{\psi\in \mathcal{U}/\sim}c_\psi\,|\sigma_\psi(m)|^2,\qquad a,c_\psi\in\mathbb{R},
\]
the sum over one representative of each class $\{\psi,\psi^{-1}\}$ in $\mathcal{U}$, and these functions are linearly independent. If $\mathcal{U}$ is empty, the only solutions are the multiples of $|m|$.

For any $B$, the set $\mathcal{U}$ is a union of classes of characters with a common kernel: if $\psi\in\mathcal{U}$ and $\ker\varphi=\ker\psi$, then $\varphi\in\mathcal{U}$. Equivalently, for a subgroup $K$ of $C$ with $C/K$ cyclic let $B_K$ be the matrix of the cosets $B_{zy}K$, with entries in $C/K$; the characters of $C$ with kernel $K$ are the faithful characters of $C/K$, and they belong to $\mathcal{U}$ if and only if $B_K^{(\psi)*}B_K^{(\psi)}=d$ for one, equivalently for every, faithful character $\psi$ of $C/K$. So for cyclic $C$ of order $q$, if $B^*B=d$ for $B=B^{(\chi)}$, then $\mathcal{U}$ contains every $\chi^j$ with $j$ coprime to $q$; for general $C$ there need be no faithful character, and the condition ``$B^*B=d$ for the faithful character'' is replaced by ``$B_K^{(\psi)*}B_K^{(\psi)}=d$ for a faithful character $\psi$ of a cyclic quotient $C/K$'', one condition for each such $K$.

$\mathcal{U}$ may contain characters that are not faithful. For $C$ cyclic of order $6$ and $B=\bigl(\begin{smallmatrix}1&1&-1\\1&\omega&-\omega^2\\1&\omega^2&-\omega\end{smallmatrix}\bigr)$ with $\omega=\zeta^2$, $\mathcal{U}=\{1,2,4,5\}$, and for $C$ cyclic of order $6$, $d=4$ and the matrix with rows $(1,1,1,1)$, $(1,1,-1,-1)$, $(1,-1,\zeta,-\zeta)$, $(1,-1,-\zeta,\zeta)$, $\mathcal{U}=\{1,3,5\}$ (residues $j$ for $\chi^j$). For $C=\{1,a,b,ab\}$ of type $C_2\times C_2$, which has no faithful character, and
\[
B=\begin{pmatrix}1&1&1&1\\1&a&b&ab\\1&b&ab&a\\1&ab&a&b\end{pmatrix},
\]
the entries generate $C$ and $\mathcal{U}$ consists of all three characters of order two, so that the solutions are
\[
F(m)=\alpha\,|m|+c_1(n_1+n_a-n_b-n_{ab})^2+c_2(n_1+n_b-n_a-n_{ab})^2+c_3(n_1+n_{ab}-n_a-n_b)^2 ,
\]
a space of dimension four; for the matrix with rows $(1,1,1,1)$, $(1,1,a,a)$, $(1,a,1,a)$, $(1,a,a,b)$ over the same group the entries generate $C$ and $\mathcal{U}$ is the single character with kernel $\{1,b\}$, the solutions being $\alpha|m|+c\,(n_1+n_b-n_a-n_{ab})^2$.
\end{corollary}

So where the gains of a block generate the group, counting and the squared moduli of the summed gains, taken through the characters $\psi\in \mathcal{U}$, are the only rules, and no hypothesis that the rule depends on a sum is needed. The matrix need not satisfy $B^*B=d$ through any character for the classification to hold; such conditions decide only which characters belong to $\mathcal{U}$. Each $|\sigma_\psi|^2$ is the squared modulus of the summed gains of the cyclic quotient $C/\ker\psi$.

\begin{proof}[Proof of Theorem~\ref{thm:multiset}]
We prove the statements for complex-valued $F$; the real case is treated at the end. Throughout, $e_h\in M$ is the multiset consisting of $h$ once, $w_z=\sum_yB_{zy}\cdot m_y$, and we note that $|w_z|=\sum_y|m_y|$ for every $z$, since translation preserves size.

\emph{(ii) implies (i).} $L$ is additive, so $\sum_zL(w_z)=\sum_y\sum_zL(B_{zy}\cdot m_y)$, and $L(B_{zy}\cdot m)=\sum_hn_ha_{B_{zy}h}$; the condition on $a$ gives $\sum_zL(B_{zy}\cdot m)=d\,L(m)$. $Q$ satisfies \eqref{eq:MBL} by hypothesis.

\emph{Step 1: $F(0)=0$.} Put all $m_y=0$ in \eqref{eq:MBL}: $d\,F(0)=d^2F(0)$, and $d\ge2$.

\emph{Step 2: a polynomial solution with the same values on sizes at most two.} Let $F$ satisfy \eqref{eq:MBL}. Fix a total order on $C$ and put $a_h=F(e_h)$, $b_{hk}=F(e_h+e_k)-F(e_h)-F(e_k)$, and
\[
P(m)=\sum_ha_hn_h+\sum_hb_{hh}\frac{n_h(n_h-1)}2+\sum_{h<k}b_{hk}n_hn_k .
\]
$P$ is a polynomial of degree at most two in the counts, and $P(0)=0=F(0)$, $P(e_h)=a_h=F(e_h)$, $P(2e_h)=2a_h+b_{hh}=F(2e_h)$, $P(e_h+e_k)=a_h+a_k+b_{hk}=F(e_h+e_k)$ for $h\ne k$. So $P=F$ on all multisets of size at most two. Let
\[
\delta(m_1,\dots,m_d)=\sum_zP(w_z)-d\sum_yP(m_y).
\]
Since each $w_z$ is linear in $(m_y)$, $\delta$ is a polynomial of degree at most two in the $d|C|$ counts of $(m_1,\dots,m_d)$. If $\sum_y|m_y|\le2$, then every $m_y$ and every $w_z$ has size at most two, so $\delta(m_1,\dots,m_d)$ is the defect of \eqref{eq:MBL} for $F$, which is zero. A polynomial $p$ of degree at most two in variables $t_1,\dots,t_n$ that vanishes at $0$, at the $e_i$, the $2e_i$ and the $e_i+e_{i'}$ is zero: $p(0)=0$ removes the constant; $p(e_i)=p(2e_i)=0$ gives $\alpha_i+\beta_i=0$ and $2\alpha_i+4\beta_i=0$ for the coefficients of $t_i$ and $t_i^2$, so both vanish; then $p(e_i+e_{i'})$ is the coefficient of $t_it_{i'}$. Hence $\delta=0$ identically: $P$ satisfies \eqref{eq:MBL} on $(\mathbb{C}^C)^d$, in particular on $M^d$. Therefore $G=F-P$ satisfies \eqref{eq:MBL} on $M^d$ and vanishes on all multisets of size at most two.

\emph{Step 3: a solution that vanishes on sizes at most two vanishes.} Let $G$ be such a solution. We show by induction on $N\ge3$ that $G$ vanishes on the set $M_N$ of multisets of size $N$, assuming it vanishes on all smaller sizes. Let $A=\mathbb{C}[x_h:h\in C]$ and $A_n$ its homogeneous part of degree $n$, with the basis of monomials $x^m=\prod_hx_h^{n_h}$, $m\in M_n$. For $k\in C$ let $\tau_k$ be the algebra automorphism of $A$ with $\tau_k(x_h)=x_{kh}$; then $\tau_k(x^m)=\prod_hx_{kh}^{n_h}=x^{k\cdot m}$, and $x^{m+m'}=x^mx^{m'}$, so
\begin{equation}\label{eq:mono}
x^{w_z}=\prod_y\tau_{B_{zy}}\bigl(x^{m_y}\bigr).
\end{equation}
Let $\varphi\colon A_N\to\mathbb{C}$ be the linear map with $\varphi(x^m)=G(m)$. Take \eqref{eq:MBL} for $G$ with $|m_y|=n_y$ and $\sum_yn_y=N$.

(a) If at least two of the $n_y$ are positive, every $n_y$ is less than $N$, the right side of \eqref{eq:MBL} vanishes by the induction hypothesis, and by \eqref{eq:mono} the law reads $\varphi\bigl(\sum_z\prod_y\tau_{B_{zy}}(x^{m_y})\bigr)=0$. The left side is multilinear in the $x^{m_y}$, which run over bases of the $A_{n_y}$, so
\begin{equation}\label{eq:a}
\varphi\Bigl(\sum_z\prod_y\tau_{B_{zy}}(p_y)\Bigr)=0\qquad\text{for all }p_y\in A_{n_y}.
\end{equation}

(b) If $n_y=N$ for one $y$ and the other inputs are $0$, then, as $G(0)=0$ and $x^0=1$, the law gives $\varphi\bigl(\sum_z\tau_{B_{zy}}(p)\bigr)=d\,\varphi(p)$ for $p=x^{m_y}$, hence by linearity for all $p\in A_N$ and every $y$.

Put $X_\psi=\sum_h\psi(h)^{-1}x_h$ for $\psi\in\hat C$. Then $\tau_k(X_\psi)=\sum_h\psi(h)^{-1}x_{kh}=\psi(k)X_\psi$. The $X_\psi$ are a basis of $A_1$, because by the orthogonality of characters the square matrix $(\psi(h)^{-1})_{\psi,h}$ is invertible, so the products $X_{\psi_1}X_{\psi_2}\cdots X_{\psi_N}$ span $A_N$. Fix such a product $p=X_{\psi_1}\cdots X_{\psi_N}$ and suppose $\varphi(p)\ne0$. Let $s=\psi_1\psi_2\cdots\psi_N$ and $t=\psi_3\cdots\psi_N$, characters of $C$; here $N\ge3$ is used, so that $r=X_{\psi_3}\cdots X_{\psi_N}$ has positive degree. Define the vectors $u_y,v_y,r_y\in\mathbb{C}^d$ by
\[
u_y(z)=\psi_1(B_{zy}),\qquad v_y(z)=\psi_2(B_{zy}),\qquad r_y(z)=t(B_{zy}).
\]
From (b), $\tau_k(p)=s(k)p$ gives $\bigl(\sum_zs(B_{zy})-d\bigr)\varphi(p)=0$, so $\sum_zs(B_{zy})=d$; a sum of $d$ complex numbers of modulus one equals $d$ only if each is $1$, so
\begin{equation}\label{eq:diag}
u_y(z)v_y(z)r_y(z)=s(B_{zy})=1\quad\text{for all }z,y,\qquad\text{hence}\quad\sum_zu_y(z)v_y(z)r_y(z)=d .
\end{equation}
Now let $y,y',y''$ be indices that are not all equal, and distribute the factors of $p$ over the inputs: $X_{\psi_1}$ to input $y$, $X_{\psi_2}$ to input $y'$, $r$ to input $y''$, that is, $p_{\eta}$ is the product of those of $X_{\psi_1},X_{\psi_2},r$ assigned to $\eta$, and $p_\eta=1$ if none is. At least two inputs receive a factor of positive degree, the degrees add to $N$, and $\prod_\eta\tau_{B_{z\eta}}(p_\eta)=u_y(z)v_{y'}(z)r_{y''}(z)\,p$. So \eqref{eq:a} gives $\sum_zu_y(z)v_{y'}(z)r_{y''}(z)\,\varphi(p)=0$. With \eqref{eq:diag}, writing $\langle a,b\rangle=\sum_za(z)b(z)$ and $ab$ for the entrywise product,
\begin{equation}\label{eq:tri}
\langle u_yv_{y'},\,r_{y''}\rangle=d\,\delta_{yy'}\delta_{yy''}\qquad\text{for all }y,y',y''.
\end{equation}
The vectors $r_1,\dots,r_d$ are linearly independent: if $\sum_{y''}\gamma_{y''}r_{y''}=0$, pairing with $u_yv_y$ gives $d\gamma_y=0$. So they are a basis of $\mathbb{C}^d$. Take $y\ne y'$, possible because $d\ge2$. By \eqref{eq:tri} the vector $u_yv_{y'}$ pairs to zero with every vector of a basis, and the pairing is nondegenerate, so $u_yv_{y'}=0$. But its entries $\psi_1(B_{zy})\psi_2(B_{zy'})$ have modulus one. This contradiction shows $\varphi(p)=0$ for every product $p$, so $\varphi=0$ and $G=0$ on $M_N$.

\emph{Step 4: (i) implies (ii).} By Steps 2 and 3, $F=P$ on $M$, with $P$ a polynomial of degree at most two without constant term that satisfies \eqref{eq:MBL} on $(\mathbb{C}^C)^d$. Write $P=L+Q$ with $L$ its part of degree one and $Q$ its part of degree two. The defect of \eqref{eq:MBL} for $P$ is the sum of the defects for $L$ and for $Q$, which are homogeneous of degrees one and two in the counts of $(m_y)$, because the $w_z$ are linear in them; as the sum is the zero polynomial, each is zero. So $Q$ satisfies \eqref{eq:MBL} on $(\mathbb{C}^C)^d$, and $L$ does; with $m_y=e_h$ and the other inputs $0$ the law for $L$ is $\sum_za_{B_{zy}h}=d\,a_h$. A polynomial that vanishes on $\mathbb{N}^C$ is zero, so $P$, and with it $L$ and $Q$, is determined by $F$; by Step 3 applied to the difference of two solutions, $F$ is determined by its values on sizes one and two.

\emph{Step 5: the forms in characters.} The $\sigma_\psi$, $\psi\in\hat C$, are linear coordinates on $\mathbb{C}^C$, again by the invertibility of the matrix $(\psi(h))_{\psi,h}$. Write $L=\sum_\psi l_\psi\sigma_\psi$. Since $\sigma_\psi(k\cdot m)=\psi(k)\sigma_\psi(m)$, the law for $L$, which by additivity is $\sum_zL(B_{zy}\cdot m)=d\,L(m)$ for each $y$, reads $\sum_\psi l_\psi\bigl(\sum_z\psi(B_{zy})-d\bigr)\sigma_\psi=0$; so $l_\psi\ne0$ requires $\psi(B_{zy})=1$ for all $z,y$, that is $\psi\in J_0$, and conversely every such combination satisfies the law. Fix a total order on $\hat C$ and write $Q=\sum_{\psi\le\lambda}c'_{\psi\lambda}\sigma_\psi\sigma_\lambda$, uniquely. With $s_{\psi,y}=\sigma_\psi(m_y)$, which are $d|C|$ linear coordinates on $(\mathbb{C}^C)^d$, and $\sigma_\psi(w_z)=\sum_y\psi(B_{zy})s_{\psi,y}$,
\[
\sum_zQ(w_z)-d\sum_yQ(m_y)=\sum_{\psi\le\lambda}c'_{\psi\lambda}\sum_{y,y'}\bigl(\Gamma_{\psi\lambda}(y,y')-d\,\delta_{yy'}\bigr)s_{\psi,y}s_{\lambda,y'},\qquad
\Gamma_{\psi\lambda}(y,y')=\sum_z\psi(B_{zy})\lambda(B_{zy'}) .
\]
For $\psi<\lambda$ the monomials $s_{\psi,y}s_{\lambda,y'}$ are distinct for distinct $(y,y')$; for $\psi=\lambda$ the monomial $s_{\psi,y}s_{\psi,y'}$ occurs for $(y,y')$ and $(y',y)$ with the same coefficient, as $\Gamma_{\psi\psi}$ is symmetric; and monomials belonging to different pairs $\{\psi,\lambda\}$ are distinct. Hence $Q$ satisfies the law if and only if $\Gamma_{\psi\lambda}(y,y')=d\,\delta_{yy'}$ for all $y,y'$ whenever $c'_{\psi\lambda}\ne0$. For $y=y'$ this says $\sum_z(\psi\lambda)(B_{zy})=d$, that is $\psi\lambda\in J_0$; then $\lambda(B_{zy'})=\overline{\psi(B_{zy'})}$ and $\Gamma_{\psi\lambda}(y,y')$ is the entry $(y',y)$ of $B^{(\psi)*}B^{(\psi)}$, so the remaining condition is $\psi\in \mathcal{U}$. With $\varphi=\lambda^{-1}$ and $\sigma_\lambda=\overline{\sigma_\varphi}$ on $M$ this is the stated list, and the monomials $\sigma_\psi\sigma_\lambda$ over the unordered pairs with $\psi\lambda\in J_0$ and $\psi\in \mathcal{U}$ (a condition symmetric in $\psi$ and $\lambda$, since $\psi\lambda\in J_0$ gives $B^{(\lambda)}=\overline{B^{(\psi)}}$) are a basis of the complex solutions of degree two. This gives the complex dimension $|J_0|+\#\{\{\psi,\lambda\}\}$.

\emph{The real case.} If $F$ is real, then $P$ of Step 2 has real coefficients, so $L$ and $Q$ are real forms, which is (ii). The system \eqref{eq:MBL} has real coefficients, so the complex solution space is the complexification of the real one and the real dimension equals the complex dimension; the real solutions are the real-valued members of the complex solution space.
\end{proof}

\begin{proof}[Proof of Corollary~\ref{cor:multiset}]
If the entries of $B$ generate $C$ and $\psi$ is trivial on them, then $\psi$ is trivial, so $J_0=\{1\}$. By the theorem every complex solution is $l_1\sigma_1+\sum_{\psi\in \mathcal{U}}c_\psi\sigma_\psi\overline{\sigma_\psi}$, and $\sigma_\psi\overline{\sigma_\psi}=|\sigma_\psi|^2=|\sigma_{\psi^{-1}}|^2$ on $M$. A real solution equals its real part, which has the stated form with $a=\operatorname{Re}l_1$ and real $c_\psi$. Independence: $|m|$ has degree one, and the $|\sigma_\psi|^2=\sigma_\psi\sigma_{\psi^{-1}}$ for distinct classes $\{\psi,\psi^{-1}\}$ are distinct monomials in the coordinates $\sigma_\psi$.

Let $e$ be the exponent of $C$ and $\zeta_e=e^{2\pi i/e}$; all values of characters lie in $\mathbb{Q}(\zeta_e)$. Let $\psi\in\mathcal{U}$ and $\ker\varphi=\ker\psi=K$. Then $\psi$ and $\varphi$ are faithful characters of the group $C/K$, which is cyclic, of order $n$ say, because it embeds in $\mathbb{C}^\times$; so $\varphi=\psi^j$ with $j$ coprime to $n$, and since the reduction $(\mathbb{Z}/e)^\times\to(\mathbb{Z}/n)^\times$ is onto, $j$ may be taken coprime to $e$. Apply the automorphism $\zeta_e\mapsto\zeta_e^j$ of $\mathbb{Q}(\zeta_e)$, which commutes with complex conjugation, to the equations $\sum_z\overline{\psi(B_{zy'})}\psi(B_{zy})=d\,\delta_{yy'}$; they become $B^{(\varphi)*}B^{(\varphi)}=d$. The statement on $B_K$ is the same fact, as $\psi(B_{zy})$ depends only on the coset $B_{zy}K$.

In the first example $B^{(2)}$ is the entrywise square, the conjugate of the Fourier matrix of degree three, and $B^{(3)}$ has two equal columns; in the second, $B^{(3)}$ has the rows $(1,1,1,1)$, $(1,1,-1,-1)$, $(1,-1,-1,1)$, $(1,-1,1,-1)$, a real Hadamard matrix, and $B^{(2)}$, $B^{(4)}$ have equal first and second columns. In both, $B^{(1)*}B^{(1)}=d$ is checked directly. In the third example the entrywise quotient of any two distinct columns runs through $C$ once (the quotients of the columns $2,3,4$ by the first are the columns themselves; column $3$ by column $2$ gives $(1,ab,a,b)$, column $4$ by column $2$ gives $(1,b,ab,a)$, column $4$ by column $3$ gives $(1,a,b,ab)$), and the sum of a non-trivial character over $C$ is zero, so every non-trivial character is in $\mathcal{U}$; the three squares are $\sigma_\psi^2$ for the characters with kernels $\{1,a\}$, $\{1,b\}$, $\{1,ab\}$. In the fourth, the character with $\psi(a)=-1$, $\psi(b)=1$ gives the rows $(1,1,1,1)$, $(1,1,-1,-1)$, $(1,-1,1,-1)$, $(1,-1,-1,1)$, a real Hadamard matrix; for the two characters with $\psi(a)=1$ the first two columns of $B^{(\psi)}$ are equal; and for $\psi(a)=\psi(b)=-1$ the first and last columns have the product sum $1-1-1-1\ne0$.
\end{proof}

\begin{corollary}[Carriers with a block, without the sum]\label{cor:multicar}
Let $D$ be a gain digraph with gains in a finite abelian group $C$, with $d\ge2$ vertices $x_0,\dots,x_{d-1}$ whose sets of successors are one set $y_0,\dots,y_{d-1}$, and $B=(\varphi(x_y\to y_z))_{z,y}$; $B^*B=d$ is not assumed. Suppose that for some starting vertex and every $N$ there is a length at which every element of $C$ is the gain of at least $N$ walks to each $x_y$. For an event $S$ let $m_S(y)\in M$ be the multiset of the gains of the walks of $S$ that end at $y$, let $F\colon M\to\mathbb{R}$ be any function with $F(0)=0$, and $\mu(S)=\sum_yF(m_S(y))$. If $\mu$ is consistent under continuation with factor $d$, then $F=L+Q$ as in Theorem~\ref{thm:multiset}; if the entries of $B$ generate $C$, then $F(m)=a|m|+\sum_{\psi\in \mathcal{U}/\sim}c_\psi|\sigma_\psi(m)|^2$. On a carrier in which every vertex has $d$ successors and $A^{(\psi)*}A^{(\psi)}=d$ for the $\psi$ with $c_\psi\ne0$, every rule $a|m|+\sum_\psi c_\psi|\sigma_\psi(m)|^2$ is consistent; the consistency of the other forms $L+Q$ on a whole carrier is not examined.
\end{corollary}
\begin{proof}
Let $(m_y)\in M^d$ with all counts at most $N$. At the length given by the hypothesis choose, for each $y$ and each $h$, as many walks to $x_y$ with gain $h$ as the count of $h$ in $m_y$; they are distinct walks of one length from one vertex, and the event $S$ they form has $m_S(x_y)=m_y$ and the empty multiset elsewhere. A walk ending at $x_y$ with gain $h$ has one continuation to each $y_z$, with gain $B_{zy}h$, and no others, so the continuation of $S$ has the multiset $\sum_yB_{zy}\cdot m_y$ at $y_z$ and the empty multiset elsewhere. As $F(0)=0$, consistency for $S$ is \eqref{eq:MBL}, and Theorem~\ref{thm:multiset} and Corollary~\ref{cor:multiset} apply. For the last statement, the count of the continuation of an event is $d$ times its count, and $\sigma_\psi$ of the multisets at the end vertices are the components of the vector of the event for the character $\psi$, so Lemma~3.1 of \cite{g9} applies to each term.
\end{proof}

\begin{remark}[Groups that are not abelian]\label{rem:multisetna}
Steps 1, 2 and 4 of the proof do not use that $C$ is abelian, and Step 3 can be replaced by an argument without characters. So for \emph{any} finite group $C$, acting on $M=\mathbb{N}[C]$ by left translation, and any $d\times d$ matrix $B$ with entries in $C$, $d\ge2$: (i) and (ii) of Theorem~\ref{thm:multiset} are equivalent, every solution is a polynomial of degree at most two in the counts without constant term, and it is determined by its values on sizes one and two.

\emph{Step 3 without characters.} With $\varphi$, (a) and (b) as above, let $N\ge3$ and $W=\{r\in A_{N-2}:\varphi(qr)=0\text{ for all }q\in A_2\}$. For $q_1,q_2\in A_1$, $r\in A_{N-2}$ and all $y,y',y''$, by (b) if the three indices are equal and by (a) otherwise ($r$ has positive degree),
\[
\sum_z\varphi\bigl(\tau_{B_{zy}}(q_1)\,\tau_{B_{zy'}}(q_2)\,\tau_{B_{zy''}}(r)\bigr)=d\,\delta_{yy'}\delta_{yy''}\,\varphi(q_1q_2r).
\]
Let $\Theta\colon A_{N-2}^d\to A_{N-2}^d$, $\Theta(r)_z=\sum_{y''}\tau_{B_{zy''}}(r_{y''})$. If $\Theta(r)\in W^d$, then, summing the identity with $y'=y$ over $y''$, $d\,\varphi(q_1q_2r_y)=\sum_z\varphi(\tau_{B_{zy}}(q_1)\tau_{B_{zy}}(q_2)\Theta(r)_z)=0$ for all $q_1,q_2$, so $r\in W^d$. Thus $\Theta^{-1}(W^d)\subseteq W^d$, and counting dimensions, $\operatorname{im}\Theta+W^d=A_{N-2}^d$. For $y\ne y'$ the linear map $s\mapsto\sum_z\varphi(\tau_{B_{zy}}(q_1)\tau_{B_{zy'}}(q_2)s_z)$ on $A_{N-2}^d$ vanishes on $\operatorname{im}\Theta$ by the identity and on $W^d$ by the definition of $W$, hence everywhere; with $s$ supported in one component, $\varphi(\tau_{B_{zy}}(q_1)\tau_{B_{zy'}}(q_2)s_z)=0$ for all $q_1,q_2\in A_1$ and $s_z\in A_{N-2}$, and since the $\tau$ are bijective on $A_1$ and products span $A_N$, $\varphi=0$.

\emph{The forms.} Let $H$ be the subgroup generated by the entries. A sum of $d$ vectors of equal norm that equals $d$ times one of that norm consists of equal vectors; applied to $a$ and its translates, the condition on $L$ says $a_{kh}=a_h$ for $k\in H$, so the forms $L$ have dimension $|C|/|H|$. In the same way a symmetric bilinear form $b$ on $\mathbb{C}^C$ gives a solution $Q(m)=b(m,m)$ if and only if $b(k\cdot m,k\cdot m')=b(m,m')$ for $k\in H$ and $\sum_zb(m,(B_{zy}^{-1}B_{zy'})\cdot m')=0$ for $y\ne y'$. If $H=C$, then $b(e_g,e_h)=f(g^{-1}h)$ with $f(x^{-1})=f(x)$ and $f$ vanishing on the two-sided ideal of $\mathbb{C}[C]$ generated by the elements $\sum_zB_{zy}^{-1}B_{zy'}$, $y\ne y'$; so, with $\mathcal{U}$ the set of irreducible unitary representations $\rho$ such that $\sum_z\rho(B_{zy})^*\rho(B_{zy'})=0$ for all $y\ne y'$ (that is, $\rho(B)^*\rho(B)=d$ for the block matrix $\rho(B)$), and $\rho(m)=\sum_hn_h\rho(h)$,
\[
F(m)=a\,|m|+\sum_{\rho\in\mathcal{U}}\operatorname{tr}\bigl(A_\rho\,\rho(m)^*\rho(m)\bigr),
\]
with complex matrices $A_\rho$, and the dimension of the solution space is $1+\sum d_\rho(d_\rho+1)/2$ over the $\rho\in\mathcal{U}$ of real type, plus $\sum d_\rho(d_\rho-1)/2$ over those of quaternionic type, plus $\sum d_\rho^2$ over the pairs $\{\rho,\overline\rho\}$ in $\mathcal{U}$ of complex type. For $C$ abelian this is Corollary~\ref{cor:multiset}. For $C=S_3$ and $B=\bigl(\begin{smallmatrix}1&c\\1&t\end{smallmatrix}\bigr)$, $c$ of order three and $t$ a transposition, $\mathcal{U}$ is the sign character and the dimension is $2$; for the matrices $B_{zy}=r_zc^{zy}s_y$, $0\le z,y\le2$, the Fourier matrix of $\langle c\rangle$ with rows multiplied on the left by $r_z$ and columns on the right by $s_y$, the $r_z$ and $s_y$ generating $S_3$ with $c$, $\mathcal{U}$ is the representation of degree two and the dimension is $4$. The sign and the representation of degree two are both in $\mathcal{U}$ only if the quotient of any two columns runs through $S_3$ equally often, so $6$ divides $d$.
\end{remark}

\subsection{A function at each end vertex}
In Corollary~\ref{cor:multicar} the rule uses one function $F$ at every end vertex. Additivity over events with disjoint sets of end vertices gives less: $\mu(S)=\sum_yF_y(m_S(y))$ with a function $F_y\colon M\to\mathbb{R}$, $F_y(0)=0$, for each vertex. The next theorem treats this case on a whole carrier.

Let $D$ be a finite gain digraph with vertex set $V$ and gains in a finite abelian group $C$, with at most one arc $y\to z$ for each ordered pair $(y,z)$ (loops are allowed). Write $\operatorname{pred}(z)$ and $\operatorname{succ}(y)$ for the sets of predecessors and successors, $A_{zy}\in C$ for the gain of the arc $y\to z$, and, for $\psi\in\hat C$, $A^{(\psi)}$ for the complex $V\times V$ matrix with entry $\psi(A_{zy})$ at $(z,y)$ if $y\to z$ is an arc and $0$ otherwise. For a family $(m_y)_{y\in V}$ put $w_z=\sum_{y\in\operatorname{pred}(z)}A_{zy}\cdot m_y$. A set $W\subseteq V$ is \emph{closed} if $\operatorname{succ}(y)\subseteq W$ for all $y\in W$; a \emph{terminal class} is a closed set $T$ in which every vertex can be reached from every vertex. For a terminal class $T$ and $b\in T$ let $\Gamma_T$ be the set of gains of the closed walks of positive length from $b$ to $b$; it is closed under products, so it is a subgroup of the finite group $C$, and it does not depend on $b$ (if $p$ is a walk from $b$ to $b'$, $q$ one from $b'$ to $b$ and $w$ a closed walk at $b'$, then the gains of $pq$ and of $pwq$ lie in $\Gamma_T$ at $b$, hence so does the gain of $w$). For $s\in\hat C$ let
\[
E_s=\Bigl\{f\in\mathbb{C}^V:\ \sum_{z\in\operatorname{succ}(y)}s(A_{zy})f_z=d\,f_y\ \text{for all }y\Bigr\}=\ker\bigl(A^{(s)\mathsf T}-d\bigr),
\]
and for $\psi,\lambda\in\hat C$ let
\[
K_{\psi\lambda}=\bigl\{c\in\mathbb{C}^V:\ A^{(\psi)\mathsf T}\operatorname{diag}(c)\,A^{(\lambda)}=d\,\operatorname{diag}(c)\bigr\}.
\]

\begin{theorem}[Global law with a function at each vertex]\label{thm:vertexlaw}
Let $D$ be as above: a finite gain digraph with gains in a finite abelian group $C$ and at most one arc for each ordered pair of vertices. Let every vertex have exactly $d\ge2$ outgoing arcs; nothing is assumed on the numbers of incoming arcs, on connectivity, or on $A^{(\psi)*}A^{(\psi)}$. For a family $(F_y)_{y\in V}$ of functions $F_y\colon M\to\mathbb{R}$ with $F_y(0)=0$ the following are equivalent.
\begin{enumerate}
\item[(i)] For all $(m_y)\in M^V$,
\begin{equation}\label{eq:GL}
\sum_{z\in V}F_z\Bigl(\sum_{y\in\operatorname{pred}(z)}A_{zy}\cdot m_y\Bigr)=d\sum_{y\in V}F_y(m_y).
\end{equation}
\item[(ii)] $F_y=L_y+Q_y$ on $M$ for every $y$, where the $L_y$ are real linear forms and the $Q_y$ real quadratic forms on $\mathbb{R}^C$, and each of the families $(L_y)$ and $(Q_y)$ satisfies \eqref{eq:GL} for all $(m_y)\in(\mathbb{R}^C)^V$.
\end{enumerate}
In (ii) the forms are determined by the family $(F_y)$, and a solution is determined by its values on the multisets of size one and two. Moreover:
\begin{enumerate}
\item[(a)] Write $L_y=\sum_\psi l_{\psi,y}\sigma_\psi$. The family $(L_y)$ satisfies \eqref{eq:GL} if and only if $l_\psi=(l_{\psi,y})_y\in E_\psi$ for every $\psi$. The restriction of $E_s$ to the union of the terminal classes is injective, and $\dim E_s$ is the number of terminal classes $T$ with $s$ trivial on $\Gamma_T$; on such a class $|f_y|$ is constant and $s(A_{zy})f_z=f_y$ for every arc $y\to z$.
\item[(b)] Fix a total order on $\hat C$ and write $Q_y=\sum_{\psi\le\lambda}c_{\psi\lambda,y}\sigma_\psi\sigma_\lambda$. The family $(Q_y)$ satisfies \eqref{eq:GL} if and only if $c_{\psi\lambda}=(c_{\psi\lambda,y})_y\in K_{\psi\lambda}$ for all $\psi\le\lambda$. One has $K_{\psi\lambda}\subseteq E_{\psi\lambda}$.
\item[(c)] The real vector space of solutions of \eqref{eq:GL} has dimension
\[
\sum_{T}\frac{|C|}{|\Gamma_T|}\;+\sum_{\{\psi,\lambda\}}\dim_{\mathbb{C}}K_{\psi\lambda},
\]
the first sum over the terminal classes and the second over the unordered pairs of characters, $\psi=\lambda$ allowed.
\item[(d)] Without the condition $F_y(0)=0$, the solutions of \eqref{eq:GL} are the families $(F_y+k_y)$ with $(F_y)$ as above and real constants with $\sum_yk_y=0$.
\end{enumerate}
The same statements hold for complex-valued families.
\end{theorem}

\begin{corollary}[Carriers, with a function at each end vertex]\label{cor:vertexlaw}
Let $D$ be as in Theorem~\ref{thm:vertexlaw}, and let $\mu(S)=\sum_{y}F_y(m_S(y))$ for events $S$ from one starting vertex, with $m_S(y)$ the multiset of the gains of the walks of $S$ that end at $y$ as in Corollary~\ref{cor:multicar}, and with functions $F_y\colon M\to\mathbb{R}$, $F_y(0)=0$. Suppose that for some starting vertex and every $N$ there is a length at which every element of $C$ is the gain of at least $N$ walks to each vertex of $D$. (The hypothesis requires every vertex, the starting vertex included, to be the end of walks of positive length, and closed walks at the starting vertex with every gain; a carrier with a vertex that has no predecessor does not satisfy it.) Then $\mu$ is consistent under continuation with factor $d$ if and only if $(F_y)$ is as in Theorem~\ref{thm:vertexlaw}(ii). If moreover $D$ is strongly connected, then the consistent rules are exactly
\[
F_y(m)=a\,|m|+\sum_{\psi\in\mathcal{U}(D)/\sim}c_\psi\,|\sigma_\psi(m)|^2\quad\text{for every }y\in V,\qquad
\mathcal{U}(D)=\{\psi\in\hat C:\ A^{(\psi)*}A^{(\psi)}=d\},
\]
with real constants $a$, $c_\psi$ that do not depend on the vertex, the sum over one representative of each class $\{\psi,\psi^{-1}\}$; these functions are linearly independent, and the trivial character is not in $\mathcal{U}(D)$. The second statement holds for the solutions of \eqref{eq:GL} whenever $D$ is strongly connected and $\Gamma_V=C$.
\end{corollary}

If $D$ is not strongly connected the constants can depend on the vertex: for two disjoint copies of the digraph with two vertices and step matrix $\bigl(\begin{smallmatrix}1&1\\1&-1\end{smallmatrix}\bigr)$ over the group of order two, $|\sigma_\psi|^2$ on the first copy and $0$ on the second is a solution. The hypothesis of at most one arc for each ordered pair cannot be dropped: for one vertex with two loops of gain $1$ the law is $F(2m)=2F(m)$, and the function with $F(m)=|m|$ if $m=2^ke_h$ for some $k\ge0$ and $h\in C$, and $F(m)=0$ otherwise, satisfies it and is not a polynomial of degree at most two, since $F(3e_h)-3F(2e_h)+3F(e_h)-F(0)=-3$. Finiteness cannot be dropped either: on the infinite tree in which every vertex has $d$ successors and one or no predecessor, with all gains $1$, the family $F_y=g$ satisfies \eqref{eq:GL}, for families $(m_y)$ of finite support, for every function $g$ with $g(0)=0$. For $d=1$ the theorem is false as in Theorem~\ref{thm:multiset}.

\begin{lemma}\label{lem:maxp}
Let $D$ be as in Theorem~\ref{thm:vertexlaw}, $s\in\hat C$ and $f\in E_s$, $f\ne0$, and let $W$ be the set of vertices at which $|f_y|$ is maximal.
\begin{enumerate}
\item[(1)] $W$ is closed, and $s(A_{zy})f_z=f_y$ for all $y\in W$ and $z\in\operatorname{succ}(y)$.
\item[(2)] Every non-empty closed set contains a terminal class, and contains a vertex with at least two predecessors.
\item[(3)] If $f'\in E_s$ vanishes on every terminal class, then $f'=0$.
\end{enumerate}
\end{lemma}
\begin{proof}
(1) Let $\mu=\max|f|$ and $y\in W$. Then $d\mu=|d f_y|\le\sum_{z\in\operatorname{succ}(y)}|f_z|\le d\mu$, the sum having $d$ terms. So $|f_z|=\mu$ for every successor, and $d$ complex numbers $s(A_{zy})f_z$ of modulus $\mu$ whose sum has modulus $d\mu$ are equal to each other, hence each equals their mean $f_y$.

(2) Let $W'$ be closed and non-empty, and choose $y\in W'$ for which the set $\rho(y)$ of vertices that can be reached from $y$ (including $y$) has the least number of elements. $\rho(y)$ is closed and lies in $W'$. For $z\in\rho(y)$ we have $\rho(z)\subseteq\rho(y)$, so $\rho(z)=\rho(y)$ and $y$ can be reached from $z$; hence $\rho(y)$ is a terminal class. Next, the $d|W'|$ arcs that leave the vertices of $W'$ end in $W'$ and are distinct as ordered pairs. If every vertex of $W'$ had at most one predecessor, at most $|W'|$ arcs would end in $W'$; as $d\ge2$, some vertex of $W'$ has two different predecessors.

(3) If $f'\ne0$, its set $W$ is closed and non-empty by (1) and contains a terminal class by (2), on which $|f'|=\max|f'|>0$.
\end{proof}

\begin{proof}[Proof of Theorem~\ref{thm:vertexlaw}]
We treat complex-valued families; the real case follows as at the end of the proof of Theorem~\ref{thm:multiset}, the system \eqref{eq:GL} having real coefficients. Note $|w_z|=\sum_{y\in\operatorname{pred}(z)}|m_y|\le\sum_y|m_y|$. That (ii) implies (i) is clear.

\emph{Step 2: polynomials with the same values on sizes at most two.} Let $(F_y)$ satisfy \eqref{eq:GL}. For each $y$ let $P_y$ be the polynomial of Step~2 of the proof of Theorem~\ref{thm:multiset} formed with the values of $F_y$: it has degree at most two in the counts, no constant term, and $P_y=F_y$ on the multisets of size at most two. The defect $\delta=\sum_zP_z(w_z)-d\sum_yP_y(m_y)$ is a polynomial of degree at most two in the $|V||C|$ counts of $(m_y)_{y\in V}$, because each $w_z$ is linear in them. If $\sum_y|m_y|\le2$, every $m_y$ and every $w_z$ has size at most two, so $\delta$ is the defect of \eqref{eq:GL} for $(F_y)$, which is zero. By the statement on polynomials of degree at most two proved there, $\delta=0$. So $(P_y)$ satisfies \eqref{eq:GL} on $(\mathbb{C}^C)^V$, and $G_y=F_y-P_y$ is a solution on $M^V$ that vanishes on all multisets of size at most two.

\emph{Step 3: a solution that vanishes on sizes at most two vanishes.} Let $(G_y)$ be such a solution and $N\ge3$, and suppose every $G_y$ vanishes on all multisets of size less than $N$. Let $\mathcal{P}=\mathbb{C}[x_h:h\in C]$ with its homogeneous parts $\mathcal{P}_n$, the monomials $x^m$ and the automorphisms $\tau_k$, as in Step~3 of the proof of Theorem~\ref{thm:multiset} (where the ring is called $A$), and let $\varphi_z\colon\mathcal{P}_N\to\mathbb{C}$ be linear with $\varphi_z(x^m)=G_z(m)$. Let $Y\subseteq V$ be non-empty, $n_y\ge1$ for $y\in Y$ with $\sum_{y\in Y}n_y=N$, and take \eqref{eq:GL} with $|m_y|=n_y$ for $y\in Y$ and $m_y=0$ otherwise. For a vertex $z$ with $Y\not\subseteq\operatorname{pred}(z)$ the size of $w_z$ is less than $N$, so $G_z(w_z)=0$; for the other $z$, $x^{w_z}=\prod_{y\in Y}\tau_{A_{zy}}(x^{m_y})$.

(a) If $Y$ has at least two elements, every $n_y$ is less than $N$ and the right side vanishes. By multilinearity,
\[
\sum_{z:\,Y\subseteq\operatorname{pred}(z)}\varphi_z\Bigl(\prod_{y\in Y}\tau_{A_{zy}}(p_y)\Bigr)=0\qquad\text{for all }p_y\in\mathcal{P}_{n_y}.
\]

(b) If $Y=\{y\}$, then $\sum_{z\in\operatorname{succ}(y)}\varphi_z(\tau_{A_{zy}}(p))=d\,\varphi_y(p)$ for all $p\in\mathcal{P}_N$.

Fix a product $p=X_{\psi_1}\cdots X_{\psi_N}$ of the elements $X_\psi=\sum_h\psi(h)^{-1}x_h$, for which $\tau_k(X_\psi)=\psi(k)X_\psi$; these products span $\mathcal{P}_N$. Put $f_z=\varphi_z(p)$, $s=\psi_1\cdots\psi_N$, $t=\psi_3\cdots\psi_N$ and $r=X_{\psi_3}\cdots X_{\psi_N}$, of degree $N-2\ge1$. For $y\in V$ define $u_y,v_y,r_y\in\mathbb{C}^V$ by
\[
u_y(z)=\psi_1(A_{zy}),\qquad v_y(z)=\psi_2(A_{zy}),\qquad r_y(z)=t(A_{zy})\qquad\text{if }y\in\operatorname{pred}(z),
\]
and $u_y(z)=v_y(z)=r_y(z)=0$ otherwise; $r_y$ is the column $y$ of $A^{(t)}$. For $y,y',y''\in V$ give the factor $X_{\psi_1}$ to the input $y$, $X_{\psi_2}$ to $y'$ and $r$ to $y''$, let $Y=\{y,y',y''\}$ as a set and $p_\eta$ the product of the factors given to $\eta\in Y$. For $z$ with $Y\subseteq\operatorname{pred}(z)$, $\prod_\eta\tau_{A_{z\eta}}(p_\eta)=u_y(z)v_{y'}(z)r_{y''}(z)\,p$, and for the other $z$ the number $u_y(z)v_{y'}(z)r_{y''}(z)$ is zero. So (a) and (b) give
\begin{equation}\label{eq:vtri}
\sum_{z\in V}u_y(z)\,v_{y'}(z)\,r_{y''}(z)\,f_z=d\,\delta_{yy'}\delta_{yy''}\,f_y\qquad\text{for all }y,y',y''\in V.
\end{equation}
Let $R=\{z:f_z\ne0\}$.

\emph{Every vertex of $R$ has at most one predecessor.} Suppose $z_0\in R$ has predecessors $y\ne y'$, and let $g(z)=u_y(z)v_{y'}(z)f_z$. Then $g$ vanishes outside $R$, $g(z_0)\ne0$, and by \eqref{eq:vtri} $\sum_zg(z)r_{y''}(z)=0$ for every $y''$. Let $\mathsf{M}$ be the $R\times R$ matrix with entries $u_\eta(z)v_\eta(z)f_z$, $\eta,z\in R$, and $\mathsf{A}$ the $R\times R$ matrix with entries $r_{y''}(z)$, $z,y''\in R$. Since $u_\eta(z)v_\eta(z)f_z=0$ for $z\notin R$, \eqref{eq:vtri} with $y=y'=\eta$ says $\mathsf{M}\mathsf{A}=d\operatorname{diag}(f_\eta)_{\eta\in R}$, which is invertible; so the square matrix $\mathsf{A}$ is invertible. The row vector $(g(z))_{z\in R}$ is annihilated by $\mathsf{A}$, so it is zero, contradicting $g(z_0)\ne0$. This is the place where two arcs into one vertex are used.

\emph{$R$ is empty.} By (b), $\sum_{z\in\operatorname{succ}(y)}s(A_{zy})f_z=d\,f_y$, that is $f\in E_s$. If $f\ne0$, the set $W$ of Lemma~\ref{lem:maxp} is closed, non-empty and contained in $R$, and by part (2) of the lemma it contains a vertex with two predecessors, which is not in $R$. So $f=0$: $\varphi_z(p)=0$ for all $z$ and all products $p$, and every $G_z$ vanishes on the multisets of size $N$. By induction $G_y=0$ for all $y$. In place of ``$d\ge2$, so there are two different inputs'' in Theorem~\ref{thm:multiset}, what is used is that every non-empty closed set of vertices contains a vertex with two different predecessors; this follows from $d\ge2$, finiteness, and the absence of parallel arcs.

\emph{Step 4.} By Steps 2 and 3, $F_y=P_y$ on $M$. Write $P_y=L_y+Q_y$ with the parts of degree one and two. The defect of \eqref{eq:GL} for $(P_y)$ is the sum of the defects for $(L_y)$ and $(Q_y)$, homogeneous of degrees one and two in the counts of $(m_y)$; their sum is the zero polynomial, so each is zero. This is (ii). A polynomial that vanishes on $\mathbb{N}^C$ is zero, so the forms are determined by the family, and by Step~3 a solution is determined by its values on sizes one and two.

\emph{Step 5, linear forms.} The family $(L_y)$ is additive, so \eqref{eq:GL} for it is the set of equations $\sum_{z\in\operatorname{succ}(y)}L_z(A_{zy}\cdot m)=d\,L_y(m)$, one for each $y$; with $\sigma_\psi(k\cdot m)=\psi(k)\sigma_\psi(m)$ and the independence of the $\sigma_\psi$ this is $l_\psi\in E_\psi$ for every $\psi$. Now fix $s$. A terminal class $T$ is closed, so the restriction $f|_T$ of $f\in E_s$ satisfies the equations of $E_s$ for the digraph $T$. If $f|_T\ne0$, the set where $|f|_T|$ is maximal is closed and non-empty by Lemma~\ref{lem:maxp}(1) applied to $T$, hence equal to $T$, because every vertex of $T$ can be reached from every vertex; so $|f|$ is constant on $T$ and $f_y=s(A_{zy})f_z$ for every arc in $T$. Along a closed walk from $b$ to $b$ with gain $g$ this gives $f_b=s(g)f_b$, so $s$ is trivial on $\Gamma_T$; and along a walk from $y$ to $b$ with gain $g_y$ it gives $f_y=s(g_y)f_b$, so the solutions on $T$ form a space of dimension at most one. Conversely let $s$ be trivial on $\Gamma_T$. If two walks from $y$ to $b$ have gains $g,g'$ and a walk from $b$ to $y$ has gain $h$, then $hg,hg'\in\Gamma_T$ and $s(g)=s(h)^{-1}=s(g')$; so $f_y=s(g_y)$ is well defined on $T$, with $f_b=1$, and for an arc $y\to z$ in $T$ a walk from $y$ through $z$ to $b$ shows $f_y=s(A_{zy})f_z$; summing over the $d$ successors, $f$ solves the equations on $T$. So the solutions on $T$ have dimension one if $s$ is trivial on $\Gamma_T$ and zero otherwise. By Lemma~\ref{lem:maxp}(3) the restriction of $E_s$ to the union $V_1$ of the terminal classes is injective. It is onto the solutions on $V_1$: given such a solution $g$, the equations for $f$ on $V_0=V\setminus V_1$ are
\[
d\,f_y-\sum_{z\in\operatorname{succ}(y)\cap V_0}s(A_{zy})f_z=\sum_{z\in\operatorname{succ}(y)\cap V_1}s(A_{zy})g_z,\qquad y\in V_0,
\]
a square system; a solution of the homogeneous system, continued by zero on $V_1$, lies in $E_s$ (the successors of a vertex of $V_1$ are in $V_1$) and vanishes on the terminal classes, so it is zero by Lemma~\ref{lem:maxp}(3). Hence the system has a unique solution. This proves (a). The number of characters trivial on $\Gamma_T$ is $|C|/|\Gamma_T|$, which gives the first term of (c).

\emph{Step 5, quadratic forms.} The numbers $s_{\psi,y}=\sigma_\psi(m_y)$ are linear coordinates on $(\mathbb{C}^C)^V$, and $\sigma_\psi(w_z)=\sum_{y}A^{(\psi)}_{zy}s_{\psi,y}$. So
\[
\sum_zQ_z(w_z)-d\sum_yQ_y(m_y)=\sum_{\psi\le\lambda}\ \sum_{y,y'}\Bigl(\sum_zc_{\psi\lambda,z}A^{(\psi)}_{zy}A^{(\lambda)}_{zy'}-d\,\delta_{yy'}c_{\psi\lambda,y}\Bigr)s_{\psi,y}s_{\lambda,y'} .
\]
For $\psi<\lambda$ the monomials $s_{\psi,y}s_{\lambda,y'}$ are distinct for distinct $(y,y')$; for $\psi=\lambda$ the monomial occurs for $(y,y')$ and $(y',y)$ with the same coefficient; monomials of different pairs $\{\psi,\lambda\}$ are distinct. Hence the law holds if and only if every bracket vanishes, which is $c_{\psi\lambda}\in K_{\psi\lambda}$. The brackets with $y=y'$ are $\sum_{z\in\operatorname{succ}(y)}(\psi\lambda)(A_{zy})c_{\psi\lambda,z}=d\,c_{\psi\lambda,y}$, so $K_{\psi\lambda}\subseteq E_{\psi\lambda}$. This proves (b), and the coefficients $c_{\psi\lambda}$, $\psi\le\lambda$, being independent coordinates on families of quadratic forms, the complex dimension of the solutions of degree two is $\sum_{\psi\le\lambda}\dim K_{\psi\lambda}$; since $K_{\lambda\psi}$ is the same space as $K_{\psi\lambda}$ (transpose the defining equation), this is the second term of (c).

\emph{(d).} Let $(F_y)$ be any family and $k_y=F_y(0)$. With all $m_y=0$, \eqref{eq:GL} reads $\sum_zk_z=d\sum_yk_y$, that is $\sum_yk_y=0$; and if this holds, \eqref{eq:GL} for $(F_y)$ is \eqref{eq:GL} for $(F_y-k_y)$.
\end{proof}

\begin{proof}[Proof of Corollary~\ref{cor:vertexlaw}]
Let $(m_y)\in M^V$ with all counts at most $N$. At the length given by the hypothesis choose, for each $y$ and $h$, as many walks to $y$ with gain $h$ as the count of $h$ in $m_y$; the event $S$ they form has $m_S(y)=m_y$ for all $y$. A walk that ends at $y$ with gain $h$ has one continuation to each $z\in\operatorname{succ}(y)$, with gain $A_{zy}h$, so the continuation of $S$ has the multiset $w_z$ at $z$, and consistency for $S$ is \eqref{eq:GL}. Conversely \eqref{eq:GL} for all families is consistency for all events. This is the first statement. Let $D$ be strongly connected. The hypothesis gives closed walks at the starting vertex with every gain, so $\Gamma_V=C$, and $V$ is the only terminal class. By (a), $E_s=0$ for $s\ne1$ and $E_1$ consists of the constant vectors; so $L_y=l\,\sigma_1=l\,|m|$ with one constant, and by (b) $c_{\psi\lambda}=0$ unless $\lambda=\psi^{-1}$, in which case $c_{\psi\lambda}$ is a constant vector $c\,(1,\dots,1)$ and the condition $c_{\psi\lambda}\in K_{\psi\lambda}$ reads $c\,(A^{(\psi)\mathsf T}\overline{A^{(\psi)}}-d)=0$; the complex conjugate of the matrix equation is $A^{(\psi)*}A^{(\psi)}=d$. On $M$, $\sigma_\psi\sigma_{\psi^{-1}}=|\sigma_\psi|^2$. So the complex solutions are the combinations of $|m|$ and the $|\sigma_\psi|^2$, $\psi\in\mathcal{U}(D)$, the real ones those with real coefficients, and independence is as in Corollary~\ref{cor:multiset}. If the trivial character were in $\mathcal{U}(D)$, no two different vertices would have a common successor, against Lemma~\ref{lem:maxp}(2) applied to $V$.
\end{proof}

\subsection{Periodic carriers}
The hypothesis of Corollary~\ref{cor:vertexlaw}, that at one length every gain reaches every vertex, fails when the derived digraph $\hat D$ is periodic. The next theorem removes it.

Let $D$ be as in Theorem~\ref{thm:vertexlaw}, with exactly $d\ge2$ outgoing arcs at every vertex, and let $\hat D$ be strongly connected with period $r$. Fix a starting vertex $x_0$. For a vertex $u$ of $\hat D$ let $\kappa(u)\in\mathbb{Z}/r$ be the residue of the length of a walk of $\hat D$ from $(x_0,1)$ to $u$, and let $X_i=\{u:\kappa(u)=i\}$, $i\in\mathbb{Z}/r$, be the cyclic classes. For $y\in V$ put $S_i(y)=\{g\in C:(y,g)\in X_i\}$, the \emph{fibre} of $X_i$ over $y$, and $C_0=S_0(x_0)$. A family $(m_y)\in M^V$ is \emph{supported on $X_i$} if every $m_y$ is supported on $S_i(y)$; then $m_y=0$ where $S_i(y)$ is empty.

\begin{lemma}\label{lem:classes}
Let $D$ be a finite gain digraph with gains in a finite abelian group $C$, with at most one arc for each ordered pair of vertices and exactly $d\ge2$ outgoing arcs at every vertex, let $\hat D$ be strongly connected with period $r$, and let $x_0$, $\kappa$, $X_i$, $S_i(y)$, $C_0$ be as above.
\begin{enumerate}
\item[(1)] $\kappa$ is well defined, every arc of $\hat D$ leads from $X_i$ to $X_{i+1}$, and there is $n_0$ such that for all vertices $u,v$ of $\hat D$ and all $n\ge n_0$ with $n\equiv\kappa(v)-\kappa(u)$ modulo $r$ there is a walk of length $n$ from $u$ to $v$.
\item[(2)] For all vertices $u,v$ of $\hat D$, $k\ge0$ and $n\ge k+n_0$ with $n\equiv\kappa(v)-\kappa(u)$ modulo $r$, the number of walks of length $n$ from $u$ to $v$ is at least $d^k/(|V||C|)$.
\item[(3)] $C_0$ is a subgroup of $C$, every non-empty fibre $S_i(y)$ is a coset of $C_0$, for fixed $y$ the non-empty fibres are disjoint and are all the cosets of $C_0$, and for an arc $y\to z$ of $D$ with $S_i(y)$ non-empty, $S_{i+1}(z)=A_{zy}S_i(y)$. The index of $C_0$ in $C$ is $r/r_0$, with $r_0$ the period of $D$.
\item[(4)] The family $(m_S(y))_y$ of an event $S$ of length $n$ from $x_0$ is supported on $X_{n}$ (indices modulo $r$), and the count of $g$ in $m_S(y)$ is at most the number of walks of length $n$ of $\hat D$ from $(x_0,1)$ to $(y,g)$. Conversely every family with these two properties is the family of an event of length $n$. Hence for every $N$ and every $n\ge n_0+k$ with $d^k\ge N|V||C|$, every family supported on $X_n$ with all counts at most $N$ is the family of an event of length $n$.
\end{enumerate}
\end{lemma}
\begin{proof}
(1) Let $r_u$ be the greatest common divisor of the lengths of the closed walks at $u$. If a closed walk at $v$ has length $\ell$ and $p$, $p'$ are walks from $u$ to $v$ and from $v$ to $u$, then $|p|+|p'|$ and $|p|+\ell+|p'|$ are lengths of closed walks at $u$, so $r_u$ divides $\ell$; hence $r_u$ divides the length of every closed walk and $r_u=r$. If $p_1,p_2$ are walks from $u$ to $v$, then $|p_1|+|p'|\equiv0\equiv|p_2|+|p'|$ modulo $r$, so the residue of the length of a walk from $u$ to $v$ depends on $u$ and $v$ only; with $u=(x_0,1)$ this is $\kappa$, and for general $u$ it is $\kappa(v)-\kappa(u)$, as a walk from $(x_0,1)$ through $u$ to $v$ shows. A walk to $u$ followed by an arc from $u$ to $v$ gives $\kappa(v)=\kappa(u)+1$. The lengths of the closed walks at $u$ form a set of positive integers that is closed under addition and has greatest common divisor $r$, so it contains every sufficiently large multiple of $r$; adding a fixed walk from $u$ to $v$ gives walks from $u$ to $v$ of every sufficiently large length in the residue class $\kappa(v)-\kappa(u)$, and there are finitely many pairs $(u,v)$.

(2) Every vertex of $\hat D$ has $d$ outgoing arcs, so there are $d^k$ walks of length $k$ from $u$. They end at no more than $|V||C|$ vertices, so some vertex $u'$ is the end of at least $d^k/(|V||C|)$ of them, and $\kappa(u')=\kappa(u)+k$. By (1) there is a walk of length $n-k\ge n_0$ from $u'$ to $v$, since $n-k\equiv\kappa(v)-\kappa(u')$. Different walks to $u'$ give different walks to $v$.

(3) For $h\in C$ the map $(x,g)\mapsto(x,gh)$ is an automorphism of $\hat D$. It maps walks to walks, so $\kappa(x,gh)-\kappa(x,g)$ is the same for all $(x,g)$: it equals $\pi(h)=\kappa(x_0,h)$, and $\pi\colon C\to\mathbb{Z}/r$ is a homomorphism, with kernel $C_0$. If $g\in S_i(y)$, then $gh\in S_i(y)$ exactly when $\pi(h)=0$, so $S_i(y)=gC_0$; every $g\in C$ lies in the fibre $S_{\kappa(y,g)}(y)$, and fibres for different $i$ are disjoint. For an arc $y\to z$ and $g\in S_i(y)$ there is an arc of $\hat D$ from $(y,g)$ to $(z,A_{zy}g)$, so $A_{zy}g\in S_{i+1}(z)$, and both sets are cosets of $C_0$. A closed walk of $D$ at $x_0$ of length $\ell$ with gain $h$ is a walk of $\hat D$ from $(x_0,1)$ to $(x_0,h)$, so $\pi(h)\equiv\ell$ modulo $r$; every $h\in C$ is the gain of such a walk because $\hat D$ is strongly connected, and the lengths of the closed walks of $D$ at $x_0$ generate $r_0\mathbb{Z}$. Here $r_0$ divides $r$, since a closed walk of $\hat D$ projects to a closed walk of $D$ of the same length. So the image of $\pi$ is $r_0\mathbb{Z}/r$, of order $r/r_0$, and this is the index of the kernel $C_0$.

(4) A walk of $D$ of length $n$ from $x_0$ with gain $g$ and end $y$ is a walk of $\hat D$ from $(x_0,1)$ to $(y,g)$, and this is a bijection between the walks of $D$ from $x_0$ and those of $\hat D$ from $(x_0,1)$. So the walks of an event of length $n$ end, in $\hat D$, in $X_n$, and an event is obtained by choosing for each $(y,g)\in X_n$ a set of walks of length $n$ from $(x_0,1)$ to $(y,g)$. The last statement follows from (2) with $u=(x_0,1)$.
\end{proof}

\begin{theorem}[Periodic carriers]\label{thm:periodic}
Let $D$ be a finite gain digraph with gains in a finite abelian group $C$, with at most one arc for each ordered pair of vertices and exactly $d\ge2$ outgoing arcs at every vertex, and let $\hat D$ be strongly connected, with period $r$; nothing is assumed on the numbers of incoming arcs or on $A^{(\psi)*}A^{(\psi)}$. Let $x_0$, $X_i$, $S_i(y)$ and $C_0$ be as above, and let $\mu(S)=\sum_yF_y(m_S(y))$ for the events $S$ from $x_0$, with functions $F_y\colon M\to\mathbb{R}$, $F_y(0)=0$. The following are equivalent.
\begin{enumerate}
\item[(i)] $\mu$ is consistent under continuation with factor $d$.
\item[(ii)] The law \eqref{eq:GL} holds for every family $(m_y)\in M^V$ that is supported on one of the classes $X_i$.
\item[(iii)] There are real constants $a$ and $c_{\psi'}$ such that for every $y$, every $i$ and every multiset $m$ supported on the fibre $S_i(y)$
\[
F_y(m)=a\,|m|+\sum_{\psi'}c_{\psi'}\,|\sigma_\psi(m)|^2 ,
\]
the sum over one representative $\psi'$ of each class $\{\psi',\psi'^{-1}\}$ in the set $\mathcal{U}(D)|_{C_0}$ of the restrictions to $C_0$ of the characters in $\mathcal{U}(D)$, with $\psi\in\mathcal{U}(D)$ any character that restricts to $\psi'$; on multisets supported on a fibre the value does not depend on the choice of $\psi$, by (a) below.
\item[(iv)] With the same constants, $\mu(S)=a\,|S|+\sum_{\psi'}c_{\psi'}\,\lVert v^\psi_S\rVert^2$ for every event $S$ from $x_0$.
\end{enumerate}
The constants depend neither on the vertex nor on the class $i$, that is, on the length modulo $r$. Moreover:
\begin{enumerate}
\item[(a)] For $\psi,\lambda\in\hat C$ and $m$ supported on one fibre, $|\sigma_\psi(m)|^2=|\sigma_\lambda(m)|^2$ if $\psi|_{C_0}=\lambda|_{C_0}$ or $\psi|_{C_0}=\lambda^{-1}|_{C_0}$; so $\lVert v^\psi_S\rVert^2=\lVert v^\lambda_S\rVert^2$ on all events for such $\psi,\lambda$.
\item[(b)] $\mathcal{U}(D)$ is a union of cosets of the annihilator $C_0^\perp=\{\psi:\psi|_{C_0}=1\}$, a group of order $|C|/|C_0|=r/r_0$ with $r_0$ the period of $D$, and is disjoint from $C_0^\perp$.
\item[(c)] The functions $|S|$ and $\lVert v^\psi_S\rVert^2$, one for each class $\{\psi',\psi'^{-1}\}$ in $\mathcal{U}(D)|_{C_0}$, are linearly independent as functions on the events from $x_0$. So the consistent rules, as functions on events, form a real vector space of dimension $1+\#\bigl(\mathcal{U}(D)|_{C_0}/\!\sim\bigr)$, and the same number is the dimension of the space of the restrictions of the solutions $(F_y)$ of (ii) to the multisets supported on fibres.
\item[(d)] The values of $F_y$ on multisets that are not supported on a fibre do not enter (i) or (ii) and are arbitrary.
\end{enumerate}
\end{theorem}

\begin{proof}
\emph{(i) and (ii).} As in the proof of Corollary~\ref{cor:vertexlaw}, if $S$ has the family $(m_y)$ then $\operatorname{Ext}S$ has the family $(w_z)$, so consistency for $S$ is \eqref{eq:GL} for the family of $S$. By Lemma~\ref{lem:classes}(4) the family of an event is supported on one class, so (ii) implies (i); and a family supported on $X_i$ has bounded counts, so it is the family of an event of a large length $n\equiv i$, and (i) implies (ii).

\emph{The quotient carrier.} Let $V'$ be the set of pairs $(y,i)$ with $S_i(y)$ non-empty, and choose $b_{y,i}\in S_i(y)$ for each. Let $D'$ be the gain digraph on $V'$ with an arc from $(y,i)$ to $(z,i+1)$ for every arc $y\to z$ of $D$, with gain
\[
A'_{(z,i+1),(y,i)}=A_{zy}\,b_{y,i}\,b_{z,i+1}^{-1}.
\]
By Lemma~\ref{lem:classes}(3), $(z,i+1)\in V'$ and the gain lies in $C_0$. So $D'$ is a finite gain digraph with gains in the finite abelian group $C_0$; it has at most one arc for each ordered pair, because the arc determines $y\to z$, and $d$ outgoing arcs at every vertex. The map $(y,g)\mapsto((y,\kappa(y,g)),\,g\,b_{y,\kappa(y,g)}^{-1})$ is a bijection from $V\times C$ to $V'\times C_0$ that maps the arcs of $\hat D$ onto the arcs of the derived digraph of $D'$; so the latter is strongly connected. Hence $D'$ is strongly connected, and for $h\in C_0$ a walk of positive length in the derived digraph of $D'$ from $((y,i),1)$ to $((y,i),h)$ is a closed walk of $D'$ at $(y,i)$ with gain $h$: $\Gamma_{V'}=C_0$.

\emph{The law on $D'$.} For a family $(F_y)$ put $F'_{(y,i)}(m')=F_y(b_{y,i}\cdot m')$ for $m'\in M(C_0)$, the multisets on $C_0$; then $F'_{(y,i)}(0)=0$, and $m'\mapsto b_{y,i}\cdot m'$ is a bijection from $M(C_0)$ onto the multisets supported on $S_i(y)$. Let $(m'_{y'})_{y'\in V'}$ be a family that vanishes outside the layer $V'_i=\{(y,j)\in V':j=i\}$, and $m_y=b_{y,i}\cdot m'_{(y,i)}$, a family supported on $X_i$; every family supported on $X_i$ arises so. The predecessors of $(z,i+1)$ in $D'$ are the $(y,i)$ with $y\in\operatorname{pred}(z)$ and $S_i(y)$ non-empty, so
\[
b_{z,i+1}\cdot\sum_{(y,i)}A'_{(z,i+1),(y,i)}\cdot m'_{(y,i)}=\sum_{y\in\operatorname{pred}(z)}A_{zy}\cdot m_y=w_z ,
\]
and $w_z=0$ if $S_{i+1}(z)$ is empty. Hence the law \eqref{eq:GL} of $D'$ for $(F'_{y'})$ and the family $(m'_{y'})$ is the law \eqref{eq:GL} of $D$ for $(F_y)$ and $(m_y)$. For an arbitrary family $(m'_{y'})\in M(C_0)^{V'}$ let $m'^{(i)}$ be its part on the layer $V'_i$. The multiset at $(z,i+1)$ on the left of the law depends on $m'^{(i)}$ only, and $F'_{y'}(0)=0$, so both sides of the law for $(m'_{y'})$ are the sums over $i$ of the two sides for the $m'^{(i)}$. Therefore (ii) holds for $(F_y)$ if and only if $(F'_{y'})$ satisfies \eqref{eq:GL} on $D'$ for all families in $M(C_0)^{V'}$. No equation between different layers has been dropped: there is none.

\emph{Application of Theorem~\ref{thm:vertexlaw}.} $D'$ satisfies the hypotheses of Theorem~\ref{thm:vertexlaw} with the group $C_0$, it is strongly connected and $\Gamma_{V'}=C_0$. By the last statement of Corollary~\ref{cor:vertexlaw}, whose proof uses only these properties, the solutions on $D'$ are exactly
\[
F'_{y'}(m')=a\,|m'|+\sum_{\psi'\in\mathcal{U}(D')/\sim}c_{\psi'}\,|\sigma_{\psi'}(m')|^2\qquad\text{for every }y'\in V',
\]
with real constants that do not depend on $y'=(y,i)$, and these functions are linearly independent on $M(C_0)$. Steps 2 and 3 of the proof of Theorem~\ref{thm:vertexlaw} are thus used for $D'$ as they stand: a set $Y$ of vertices of $D'$ with a common successor lies in one layer, so every instance of \eqref{eq:GL} used in Step~3 there ($Y$ a single vertex, or a set with a common successor) is an instance for a family on one layer, and the coefficients of the defect polynomial at products of counts of different layers are zero identically.

\emph{$\mathcal{U}(D')=\mathcal{U}(D)|_{C_0}$.} Let $\psi'\in\hat C_0$ and let $\psi\in\hat C$ restrict to $\psi'$; such $\psi$ exist because $C$ is a finite abelian group. Then $\psi'(A'_{(z,i+1),(y,i)})=\overline{\psi(b_{z,i+1})}\,\psi(A_{zy})\,\psi(b_{y,i})$, so $A'^{(\psi')}=\Delta^{-1}\tilde A\Delta$ with $\Delta$ the diagonal unitary matrix of the $\psi(b_{y,i})$ and $\tilde A$ the $V'\times V'$ matrix with entry $\psi(A_{zy})$ at $((z,i+1),(y,i))$ for an arc $y\to z$ and $0$ elsewhere. So $\psi'\in\mathcal{U}(D')$ if and only if $\tilde A^*\tilde A=d$. The entry of $\tilde A^*\tilde A$ at $((y,i),(y_1,i_1))$ is $0$ if $i\ne i_1$, and is the entry of $A^{(\psi)*}A^{(\psi)}$ at $(y,y_1)$ if $i=i_1$. If $y\ne y_1$ have a common successor $z$ and $S_i(y)$ is non-empty, then $S_i(y_1)$ is non-empty, since $A_{zy_1}^{-1}S_{i+1}(z)$ is a fibre $S_j(y_1)$ and $j+1=i+1$; so every non-zero entry of $A^{(\psi)*}A^{(\psi)}$ off the diagonal occurs in $\tilde A^*\tilde A$, and the diagonal entries of both matrices are $d$. Hence $\tilde A^*\tilde A=d$ if and only if $A^{(\psi)*}A^{(\psi)}=d$: a character of $C_0$ lies in $\mathcal{U}(D')$ if and only if one, equivalently every, character of $C$ that restricts to it lies in $\mathcal{U}(D)$. This gives the first part of (b); the trivial character is not in $\mathcal{U}(D')$ by Corollary~\ref{cor:vertexlaw}, which gives the second; the order of $C_0^\perp$ is the index of $C_0$, which is $r/r_0$ by Lemma~\ref{lem:classes}(3).

\emph{(ii) and (iii).} For $m=b_{y,i}\cdot m'$ and $\psi$ restricting to $\psi'$, $|m|=|m'|$ and $\sigma_\psi(m)=\psi(b_{y,i})\sigma_{\psi'}(m')$, so $|\sigma_\psi(m)|^2=|\sigma_{\psi'}(m')|^2$. With the last three paragraphs, (ii) holds if and only if (iii) does. This also proves (a), since $|\sigma_{\psi'^{-1}}|^2=|\sigma_{\psi'}|^2$ on $M(C_0)$, and (d).

\emph{(iii) and (iv).} The component of $v^\psi_S$ at $y$ is $\sigma_\psi(m_S(y))$ and $|S|=\sum_y|m_S(y)|$, and $m_S(y)$ is supported on a fibre; so (iii) implies (iv). If (iv) holds, let $(\tilde F_y)$ be given by the formula of (iii) on all of $M$. Then $\mu(S)=\sum_y\tilde F_y(m_S(y))$, and since (iii) implies (i) for $(\tilde F_y)$, $\mu$ is consistent.

\emph{(c).} Suppose $a|S|+\sum_{\psi'}c_{\psi'}\lVert v^\psi_S\rVert^2=0$ for all events. For $m'\in M(C_0)$ there is by Lemma~\ref{lem:classes}(4) an event whose family is $b_{x_0,0}\cdot m'$ at $x_0$ and $0$ elsewhere, so $a|m'|+\sum c_{\psi'}|\sigma_{\psi'}(m')|^2=0$ on $M(C_0)$, and the constants vanish by the independence on $M(C_0)$. The same argument shows that the constants are determined by the restrictions of the $F_y$ to the fibres.
\end{proof}

\begin{corollary}\label{cor:periodic}
Under the hypotheses of Theorem~\ref{thm:periodic}:
\begin{enumerate}
\item[(1)] If $D$ and $\hat D$ have the same period, then $C_0=C$, every fibre is empty or all of $C$, and the consistent rules are $\mu(S)=a|S|+\sum_{\psi\in\mathcal{U}(D)/\sim}c_\psi\lVert v^\psi_S\rVert^2$, as in Corollary~\ref{cor:vertexlaw}, although no length has all gains at all vertices unless $r=1$. For the signed arc digraph of a ring with an even number of sites ($r=r_0=2$, $C=\{\pm1\}$, $\mathcal{U}(D)$ the sign character) the consistent rules are $a|S|+c\lVert v_S\rVert^2$, with the same two constants at even and at odd lengths.
\item[(2)] Cross terms do not survive. Let $\psi\ne\lambda$ be characters of $C$, $c$ complex, and $F_y=\operatorname{Re}\bigl(c\,\sigma_\psi\overline{\sigma_\lambda}\bigr)$ for every $y$. If the rule of this family is consistent, it is zero on every event. If $\psi|_{C_0}=\lambda|_{C_0}$, then on a multiset $m$ supported on $S_i(y)$ the term equals $\operatorname{Re}(c\,\omega_{y,i})\,|\sigma_\psi(m)|^2$, with $\omega_{y,i}$ the value of $\psi\lambda^{-1}$ on $S_i(y)$, a root of unity that depends on the vertex and on the class: on events such a term is a square $|\sigma_\psi|^2$ with a coefficient that varies, which the theorem excludes.
\end{enumerate}
\end{corollary}
\begin{proof}
(1) By Lemma~\ref{lem:classes}(3) the index of $C_0$ is $r/r_0$, so $C_0=C$ exactly when $r=r_0$; then (iii) and (iv) of the theorem are the stated formula. For the ring, the columns of $A^{(\psi)}$ for $(x,R)$ and $(x,L)$ are orthogonal for the sign character and the other pairs of columns have disjoint supports.

(2) Put $\psi'=\psi|_{C_0}$, $\lambda'=\lambda|_{C_0}$ and $\epsilon=\psi\lambda^{-1}$. For $m=b_{y,i}\cdot m'$ one has $\sigma_\psi(m)\overline{\sigma_\lambda(m)}=\epsilon(b_{y,i})\,\sigma_{\psi'}(m')\,\sigma_{\lambda'^{-1}}(m')$, so $F'_{(y,i)}=\operatorname{Re}\bigl(c\,\epsilon(b_{y,i})\,\sigma_{\psi'}\sigma_{\lambda'^{-1}}\bigr)$, a family of quadratic forms on $D'$; the rule is consistent if and only if this family solves \eqref{eq:GL} on $D'$, and it is zero on events if the family is zero. By Theorem~\ref{thm:vertexlaw}(b) and the proof of Corollary~\ref{cor:vertexlaw}, applied to $D'$, a family of quadratic forms that solves the law has coefficient vector zero at every monomial $\sigma_\alpha\sigma_\beta$ with $\beta\ne\alpha^{-1}$, and a constant coefficient vector at $\sigma_\alpha\sigma_{\alpha^{-1}}$.

Let $\psi'\ne\lambda'$. The family is a combination of the monomial $\sigma_{\psi'}\sigma_{\lambda'^{-1}}$ and its conjugate $\sigma_{\psi'^{-1}}\sigma_{\lambda'}$, and $\lambda'^{-1}\ne\psi'^{-1}$; so a solution has all coefficients zero and is the zero family.

Let $\psi'=\lambda'$. Then $\epsilon\in C_0^\perp$, $\epsilon\ne1$, $\epsilon$ is constant on each coset of $C_0$, and $F'_{(y,i)}=\operatorname{Re}(c\,\omega_{y,i})\,|\sigma_{\psi'}|^2$ with $\omega_{y,i}=\epsilon(b_{y,i})$. A solution has a constant coefficient, $\operatorname{Re}(c\,\omega_{y,i})=c_0$ for all $(y,i)\in V'$. By Lemma~\ref{lem:classes}(3) the fibres over a fixed $y$ are all the cosets of $C_0$, so the $\omega_{y,i}$ at $y$ are the numbers $\omega\eta$ with $\omega$ one of them and $\eta$ in the group $\epsilon(C)$ of roots of unity, of order $k\ge2$, whose elements sum to zero. Summing $\operatorname{Re}(c\,\omega\,\eta)=c_0$ over $\eta$ gives $k\,c_0=0$, so the family is zero.
\end{proof}

The count of independent rules can be smaller than the count $1+\#(\mathcal{U}(D)/\!\sim)$ that the formula of Corollary~\ref{cor:vertexlaw} would give. For two vertices and the step matrix $\zeta\bigl(\begin{smallmatrix}1&1\\1&-1\end{smallmatrix}\bigr)$ over the eighth roots of unity, $\hat D$ is strongly connected with $r=4$, $r_0=1$ and $C_0=\{\pm1\}$, $\mathcal{U}(D)$ consists of the four faithful characters, in two classes $\{\chi,\chi^{-1}\}$, $\{\chi^3,\chi^{-3}\}$, and all four restrict to the sign character of $C_0$: $\lVert v^{\chi}_S\rVert^2=\lVert v^{\chi^3}_S\rVert^2$ on every event, and the consistent rules have dimension two, not three. For two vertices and the step matrix $a\bigl(\begin{smallmatrix}1&1\\1&b\end{smallmatrix}\bigr)$ over $\{1,a,b,ab\}$ of type $C_2\times C_2$, $r=2$, $r_0=1$, $C_0=\{1,b\}$, $\mathcal{U}(D)$ consists of the two characters with $\psi(b)=-1$, which are not inverse to each other and coincide on events; the dimension is again two. For the four-vertex example after Corollary~\ref{cor:abund} ($q=4$, $r=2$, $r_0=1$, $C_0=\{\pm1\}$, $\mathcal{U}(D)=\{\chi,\chi^{-1}\}$) the consistent rules are $a|S|+c\lVert v_S\rVert^2$, of dimension two.

\subsection{Composition: zero, the count or one character}
Theorem~\ref{thm:periodic} leaves a real vector space of consistent rules, spanned by the count and by squares. As in Proposition~\ref{prop:compo}(b), composition removes the combinations, and now the mixtures of the count with a square as well.

For events $S_1$, $S_2$ of gain digraphs $D_1$, $D_2$ with gains in $C$, of one length, let $S_1\times S_2$ be the set of pairs $(w_1,w_2)$, $w_i\in S_i$, a pair ending at the pair of the end vertices and having as gain the product of the two gains; this is the event of the product of \cite{g8}, Section~6. Then $|S_1\times S_2|=|S_1|\,|S_2|$, and for every character $\psi$ the component of $v^\psi_{S_1\times S_2}$ at $(y_1,y_2)$ is the product of the components of $v^\psi_{S_1}$ at $y_1$ and of $v^\psi_{S_2}$ at $y_2$, so $\lVert v^\psi_{S_1\times S_2}\rVert^2=\lVert v^\psi_{S_1}\rVert^2\,\lVert v^\psi_{S_2}\rVert^2$.

\begin{corollary}[Composition]\label{cor:composition}
Let $D$ be as in Theorem~\ref{thm:periodic} and let $\mu$ be a consistent rule of that theorem, $\mu(S)=a\,|S|+\sum_{\psi'}c_{\psi'}\lVert v^\psi_S\rVert^2$. For a product event put $\mu(S_1\times S_2)=a\,|S_1\times S_2|+\sum_{\psi'}c_{\psi'}\lVert v^\psi_{S_1\times S_2}\rVert^2$, the same formula. If $\mu(S_1\times S_2)=\mu(S_1)\,\mu(S_2)$ for all events $S_1$, $S_2$ of $D$ from $x_0$ of one length, at every sufficiently large length, then $\mu$ is one of the following: zero; the count $|S|$; or $\lVert v^\psi_S\rVert^2$ for one class $\{\psi',\psi'^{-1}\}$ in $\mathcal{U}(D)|_{C_0}$. Each of these is multiplicative.
\end{corollary}
\begin{proof}
Let $f_0(S)=|S|$ and let $f_1,\dots,f_k$ be the functions $\lVert v^\psi_S\rVert^2$, one for each class, with coefficients $e_0=a$ and $e_1,\dots,e_k$. Each $f_i$ is multiplicative on product events by the remarks before the statement, so $\mu(S_1\times S_2)=\sum_ie_if_i(S_1)f_i(S_2)$, while $\mu(S_1)\mu(S_2)=\sum_{i,j}e_ie_jf_i(S_1)f_j(S_2)$. Fix a vertex $y$. By Lemma~\ref{lem:classes}(3) the non-empty fibres over $y$ are all the cosets of $C_0$, so $S_n(y)$ is non-empty for $n$ in some residue class modulo $r$; choose such an $n$ so large that the hypothesis holds at length $n$ and every family supported on $X_n$ with counts at most $2$ is the family of an event of length $n$ (Lemma~\ref{lem:classes}(4)). For a multiset $m'$ on $C_0$ of size at most two there is then an event of length $n$ all of whose walks end at $y$, with $m_S(y)=b\cdot m'$ for a fixed $b\in S_n(y)$, and on it $f_0(S)=|m'|$ and $f_i(S)=|\sigma_{\psi'}(m')|^2$. These $k+1$ functions of $m'$ are polynomials of degree at most two in the counts without constant term, linearly independent on $M(C_0)$ by the proof of Theorem~\ref{thm:periodic}; a polynomial of degree at most two without constant term that vanishes on the multisets of size at most two is zero (Step~2 of the proof of Theorem~\ref{thm:multiset}), so they are linearly independent on the multisets of size at most two. Hence $f_0,\dots,f_k$ are linearly independent as functions on the events of length $n$. The hypothesis reads $\sum_ie_if_i(S_1)f_i(S_2)=\bigl(\sum_ie_if_i(S_1)\bigr)\bigl(\sum_je_jf_j(S_2)\bigr)$. For fixed $S_2$ this is a linear relation among the $f_i(S_1)$, so by independence $e_if_i(S_2)=e_i\sum_je_jf_j(S_2)$ for every $i$ and every $S_2$; for fixed $i$ this is a linear relation among the $f_j(S_2)$, so $e_ie_j=\delta_{ij}e_i$ for all $i,j$. Hence every $e_i$ is $0$ or $1$ and at most one is $1$. This is the argument of Proposition~\ref{prop:compo}(b) in the form in which it was checked in Lean.
\end{proof}

So, on these carriers, a rule that adds over end vertices a function of the multiset of gains, is consistent under continuation and is multiplicative under composition is zero, the count of walks, or the squared norm for one character, up to the identification of characters that agree on events; it is not a mixture, and its normalization is fixed. Consistency and composition together leave a choice between finitely many rules: zero, one that ignores the gains, and one for each class of characters in which the step is norm-preserving. The length must be large: at length two on $B_4$ and on the ring of four sites every event has $|S|=\lVert v_S\rVert^2$, and the mixtures are multiplicative there. Nothing here makes that choice. The corollary concerns the product events $S_1\times S_2$ only, and assumes the same formula for $\mu$ on the product; it does not examine consistency on the product carrier, nor products of two different carriers.

Corollary~\ref{cor:composition} and its proof were written by the author; a challenger found no failing step and, on four carriers, by enumeration of all pairs of events of length three (up to $65{,}536$ pairs), that exactly the rules named are multiplicative, and that at length two some mixtures are multiplicative as well, which is why the hypothesis asks for large lengths; the algebraic step $e_ie_j=\delta_{ij}e_i$ is checked in Lean~4 with Mathlib over any commutative ring, and its consequence that each $e_i$ is $0$ or $1$ over a domain \cite{comp4s}. Carriers with $C_0\ne C$ were not computed for this corollary.

Theorem~\ref{thm:periodic} was asked for after a check showed that the hypothesis of Corollary~\ref{cor:vertexlaw} holds on rings with an odd number of sites and on $B_q$ and fails on rings with an even number; statement and proof were written by a sub-agent and reread for this paper. Exact ranks for $107$ qualifying carriers agree with the dimension $1+\#(\mathcal{U}(D)|_{C_0}/\!\sim)$, a challenger found the same on $110$ further carriers with no failing step, six boundary cases agree, and the fact on sums of roots of unity used in Corollary~\ref{cor:periodic}(2) is checked in Lean~4 with Mathlib \cite{comp4s}.

Theorem~\ref{thm:vertexlaw} and its proof were found by a sub-agent asked whether the function could depend on the vertex; each step was reread for this paper. Exact ranks for $146$ carriers agree with the dimension formula, a challenger found the same on $68$ further carriers ($190$ of $190$ instances) with no failing step, ten boundary cases agree, and three steps (the equality case in Lemma~\ref{lem:maxp}(1), the matrix step, the pigeonhole of Lemma~\ref{lem:maxp}(2)) are checked in Lean~4 with Mathlib \cite{comp4s}. On a carrier as in Theorem~\ref{thm:vertexlaw} that is strongly connected and in which every gain reaches every vertex in abundance, Corollary~\ref{cor:vertexlaw} sharpens Corollary~\ref{cor:multicar}: it needs no block and no condition on the entries of a block, and its characters are those for which the whole step is norm-preserving, a subset of those that qualify for one block, and it can be a proper one: on the three vertices $0,1,2$ with the arcs $0\to1$, $0\to2$, $1\to1$, $2\to0$, $2\to2$ of gain $+1$ and $1\to2$ of gain $-1$ the sign character qualifies for the block of the vertices $0,1$ and not for the step ($A^{\mathsf T}A$ has the entries $1$ and $-1$ off the diagonal), so by Corollary~\ref{cor:vertexlaw} only the count is consistent there; the carrier is strongly connected, and at lengths $8$ and $12$ every gain reaches every vertex at least $18$ and $448$ times from every starting vertex. Corollary~\ref{cor:multicar} needs abundance at the vertices of one block only and applies to carriers that are not strongly connected. Not treated: parallel arcs together with merging, unequal out-degrees, groups that are not abelian, and a closed form for the spaces $K_{\psi\lambda}$ when the carrier is not strongly connected.

Corollary~\ref{cor:block} concerns the case of $C$ the $q$-th roots of unity and $F=\Phi\circ\sigma_\chi$, with $\chi$ the faithful character. For $q\ge2$ the map $\sigma_\chi$ sends $M$ onto $R$, because $-1$ is a sum of roots of unity, so for such $F$ the law \eqref{eq:MBL} is the law \eqref{eq:BL} for $\Phi$, and Theorem~\ref{thm:multiset} shows by a second method that $\Phi\circ\sigma_\chi$ is a polynomial of degree at most two in the counts; the count itself is excluded because it is not a function of the sum (for even $q$ the multiset of $1$ and $-1$ has sum zero and count two; for any $q\ge2$ the multiset of all $p$-th roots of unity, $p$ a prime dividing $q$, has sum zero). The description of $\Phi$ as a quadratic form on $V$ is not derived again here. Theorem~\ref{thm:multiset} joins Theorem~\ref{thm:block} to Proposition~\ref{prop:add}: on a carrier with a block whose entries generate the group, and with abundant gains, a consistent rule that sees a walk only through its end vertex and its gain, and adds over end vertices a function of what it sees that vanishes on the empty multiset, is a combination of the count and of the squares $|\sigma_\psi|^2$, $\psi\in \mathcal{U}$. Each of these is the squared modulus of a sum, over the walks with a common end vertex, of the gains taken through the character $\psi$; for a faithful character of a cyclic group this is a conjugate of the summed gains, in general it is the summed gains of the cyclic quotient $C/\ker\psi$. That gains with a common end vertex are added, through some character, is then a consequence for every part of the rule that is not the count, and not an assumption. Without the condition that the entries generate the group it is false: for the cyclic group of order $4$ and a real Hadamard block the solutions include $\sigma_{\chi^2}$ and $\operatorname{Re}\sigma_\chi^2$, and for $q=2$ and the all-ones matrix the number of walks with gain $1$. What is assumed is that the rule is a sum over end vertices, that it depends on the walks at a vertex only through the multiset of their gains, and the carrier with its abundance of gains. The count and the squares can be mixed with any real coefficients; positivity restricts them as in Proposition~\ref{prop:compo}, and nothing here selects between them or among the $\psi\in \mathcal{U}$. The theorem, with the observation that $B^*B=d$ is not needed, and its proof were found by a sub-agent that had been given the conjecture with exact solution spaces for five blocks; each step was rechecked for this paper. Exact ranks over the rationals for $62$ matrices, $18$ of them blocks, agree with the dimension formula in $216$ of $216$ instances \cite{comp4s}; a challenger's program found the same on $118$ of $118$ instances for $41$ other matrices. Three steps of the proof (a sum of $d$ numbers of modulus one that equals $d$; the vanishing of $u_yv_{y'}$; the vanishing of a polynomial of degree two) are checked in Lean~4 with Mathlib, the second by a shorter argument: the matrices $(r_{y''}(z)/d)$ and $(u_y(z)v_y(z))$ are inverse to each other. The theorem was first issued for cyclic groups; the present form for finite abelian groups uses the same proof with characters in place of residues, and exact ranks for $64$ matrices over five groups that are not cyclic agree with the dimension formula in $178$ of $178$ instances \cite{comp4s}. Remark~\ref{rem:multisetna} on groups that are not abelian rests on a different argument for Step~3; for the symmetric group on three letters exact ranks for $17$ matrices agree with its dimension formula, and a challenger found the same on $93$ of $93$ instances for that group and the quaternion group, where the term for representations of quaternionic type occurs. The linear-algebra step of the remark (if $\Theta^{-1}(W^d)\subseteq W^d$ then $\operatorname{im}\Theta+W^d$ is everything) is checked in Lean~4 with Mathlib, as are four steps of Theorem~\ref{thm:vsl} and Proposition~\ref{prop:abund} formalized by the review queue (the polynomial step, the step by Newton's identities, a multiplicatively closed non-empty subset of a finite group is a subgroup, and $A^*A=d$ implies $AA^*=d$ for a square matrix). For $d=1$ the theorem is false.

Taken together with Proposition~\ref{prop:add} this gives the following division, under the assumption that events are sets of histories. A switching-invariant rule of first order on a carrier without loops and of constant out-degree is the count and is blind to the gains. A rule of the form \eqref{eq:phi} that is consistent on a finite connected signed norm-preserving carrier of degree two is the squared modulus, and on a norm-preserving carrier in which two walks of equal length merge it is the only consistent power; with gains in the roots of unity and abundance of gains every consistent rule of that form is a combination of the conjugate squared moduli (Corollaries~\ref{cor:vsl}, \ref{cor:abund}), and on carriers with a full block the modulus need not be assumed (Corollary~\ref{cor:block}). Where no two walks ever meet, as on a tree, every rule of this form is the count. The second order of the rule of \cite{g8} is then a consequence and no longer a separate assumption; the choice of the character is not decided (Theorem~\ref{thm:conj}), and it does not arise when the gains are recorded (Corollary~\ref{cor:lift}). Among pair rules with a kernel, with no assumption on records beyond the end vertices, the consistent rules are, under the hypothesis on the quotients, the real combinations of the squared-norm rules of the characters for which the step is norm-preserving (Theorem~\ref{thm:kernel}), in degree two those in which $-1$ acts as $-1$ (Corollary~\ref{cor:kernel2}). On carriers with a block whose entries generate the group and with abundant gains the dependence on sums of gains is itself a consequence, up to the count: the rule is a combination of the count and of squared moduli of gains summed through the characters $\psi\in\mathcal{U}$, which need not be faithful, for rules that add over end vertices a function of the multiset of gains (Corollary~\ref{cor:multicar}). The record role of end vertices, the dependence on a history through its gain alone, and the carrier remain assumed, and rules that are not sums over end vertices are not covered.

\section{Gains of bounded order and a force}\label{sec:force}

Let $U$ be the normalized step of the signed arc digraph of the line or of a ring, $\mathcal{T}$ the translation by one site, and for a function $m$ of modulus one on the sites let $M_m$ be the multiplication by $m(x)$; $UM_m$ is the step in which every arc leaving site $x$ carries the additional gain $m(x)$. Proposition~6.3 of \cite{g9} shows that $m(x)=e^{-i\Phi x}$ translates the quasi-momentum by $-\Phi$ per step.

\begin{proposition}\label{prop:exact}
$\mathcal{T}\,UM_m\,\mathcal{T}^{-1}=\lambda\,UM_m$ for a number $\lambda$ if and only if $m(x-1)=\lambda\,m(x)$ for all $x$, that is, $m(x)=m(0)\lambda^{-x}$. If $m$ takes its values in a finite group of roots of unity, $\lambda$ belongs to that group; an exact translation of the quasi-momentum by $\Phi=2\pi/n$ per step, that is, $\lambda=e^{\pm i\Phi}$, requires gains of order at least $n$; more precisely the order of the group is a multiple of the order of $\lambda$, and on a ring of $N$ sites $\lambda^N=1$.
\end{proposition}
\begin{proof}
$\mathcal{T}$ commutes with $U$ and $\mathcal{T}M_m\mathcal{T}^{-1}$ is the multiplication by $m(x-1)$; $U$ is invertible.
\end{proof}

With gains of bounded order an exact force is therefore the linear gain of that order, a strong field. A weak mean field can only be approximate, with a gain that winds slowly: constant on long blocks and advancing by one group element from block to block, as in the staircase of \cite{g9}. On a block where the gain is the constant $\gamma$ the step is $\gamma U$. By Lemma~6.1 of \cite{g9} the spectrum of $U$ on the line is
\[
\Sigma=\{e^{i\theta}:|\theta|\le\pi/4\}\cup\{e^{i\theta}:|\theta-\pi|\le\pi/4\},
\]
two arcs of length $\pi/2$ separated by two gaps of length $\pi/2$.

\begin{proposition}[Local spectra]\label{prop:local}
Let $\zeta=e^{2\pi i/q}$. The set $\bigcap_{j=0}^{q-1}\zeta^{\,j}\Sigma$ equals $\Sigma$ for $q=1,2$, consists of the four points $e^{\pm i\pi/4},e^{\pm3i\pi/4}$ for $q=4$, and is empty for $q=3$ and for $q\ge5$.
\end{proposition}
\begin{proof}
A point lies in the intersection exactly when its orbit under multiplication by $\zeta$ avoids both open gaps. $\Sigma=-\Sigma$ settles $q\le2$. For $q\ge5$ consecutive points of an orbit are at angular distance $2\pi/q<\pi/2$, and an open arc of length $\pi/2$ contains one of them. For $q=4$ the distance is $\pi/2$ and an orbit avoids a gap only by containing both its end points. For $q=3$ two of the three points of an orbit would lie in one arc of length $\pi/2$, at distance $2\pi/3$.
\end{proof}

For $q\ge3$ no quasi-energy is thus an interior point of the spectrum on every block of a staircase that winds through the group; for $q=4$ the common points are band edges, where the group velocity vanishes. A state that propagates in one block is in a gap of another. This is a statement about the spectra of the pieces. It is not a theorem about the dynamics of the staircase, whose gains are periodic in position, so that the theorem of Grimmett, Janson and Scudo \cite{gjs} applies with the period as cell and exact confinement is not to be expected. The computed instances are as follows \cite{comp4s}.

\begin{table}[ht]\centering\small
\begin{tabular}{@{}lrrrrr@{}}\toprule
block length $L$ & 1 & 2 & 4 & 16 & 64\\\midrule
order 2 & $1108$ & $750$ & $403$ & $346$ & $347$\\
order 3 & $520$ & $520$ & $160$ & $9.3$ & $31$\\
order 4 & $750$ & $365$ & $102$ & $8.4$ & $34$\\
order 6 & $520$ & $181$ & $27$ & $11$ & $48$\\\bottomrule
\end{tabular}
\caption{Root-mean-square distance, in sites, after $2048$ steps of a walk from one vertex under a staircase of gains of order $q$ with blocks of $L$ sites; $L=1$ is the linear gain of order $q$. Ring of $8192$ sites; the walk does not reach the far side \cite{comp4s}.}\label{tab:stair}
\end{table}

\emph{Observed.} (a) For $L=1$ the distance doubles between $1024$ and $2048$ steps: a linear gain with a rational field spreads ballistically, as proved by Cedzich et al.\ \cite{cedzich}; the bounded distance reported in Section~6.2 of \cite{g9} refers to four periods. (b) For $L=16$ and $L=64$ and the orders $3,4,6$ the walk stays within about one block up to $2048$ steps; for $L=4$ it does not. Order $2$ spreads for every $L$, in accordance with $\Sigma=-\Sigma$. (c) A packet meeting one interface with gain ratio $e^{-i\delta}$ is transmitted with probability $0.9997$, $0.9994$, $0.975$ for $\delta=\pi/12,\pi/6,\pi/4$, with probability $0.0025$ for $\pi/3$, below $10^{-12}$ for $\pi/2$ and $2\pi/3$, and $0.48$ for $\pi$; the packet has quasi-energy near $-0.21$, which lies in the shifted band for $\delta<0.996$. (d) Under a staircase of order $4$ a packet is displaced, after an eighth of the period of the mean field, by $0.69$, $0.54$ and $0.16$ of its displacement under the linear gain of the same mean slope, for $L=16,64,256$: the deviation from the mean field grows with the block length. The first two values reproduce those of \cite{g9}.

\emph{Failed expectations, kept.} Two predicates fixed before their runs failed. The first expected the staircase of period $64$ to drift within $1024$ steps; it stayed within four sites of its start. The second expected confinement for $L=4$. A computation of band velocities from sorted eigenphases returned values above the largest possible speed and below the observed spread, and is discarded as an instrument. All three are in the result files.

In these instances a winding gain of order $3$, $4$ or $6$ either is a strong field, or holds the walk within a block, or leaks slowly. Where the displacement of a packet was computed, for order $4$, it does not follow the ray of the mean field; for orders $3$ and $6$ only the distances of Table~\ref{tab:stair} were computed, and no packet displacement. Whether some other arrangement of gains of bounded order has a ray limit with a force is not decided.

\section{The position of a ray under a linear gain}\label{sec:ray}

Table~4 of \cite{g9} compares the centroid of a packet with the ray $\dot k=-\Phi$, $\dot x=v(k)$ through the central quasi-momentum $k_0$, and reports a deviation of about $0.6$ site, most of it the change of the connection $\mathcal{A}(k)=\operatorname{Im}\langle u(k)|\partial_ku(k)\rangle$ of the branch eigenvector. Extending that comparison to longer periods shows that the remainder does not decrease: it is $0.075$, $0.038$, $0.038$, $0.062$, $0.117$, $0.232$ site at the periods $64$ to $2048$, and a predicate that expected a first-order decrease failed \cite{comp4s}. The cause is the width of the packet in $k$. If the ray and the connection are averaged over the quasi-momentum distribution $\rho$ of the initial packet,
\begin{equation}\label{eq:avg}
\bar x(t)=\sum_{s=1}^{t}\int v\bigl(k-(s-\tfrac12)\Phi\bigr)\,\rho(k)\,dk-\int\bigl[\mathcal{A}(k-t\Phi)-\mathcal{A}(k)\bigr]\rho(k)\,dk ,
\end{equation}
with the signs of the conventions of \cite{g9}, the remainder decreases (Table~\ref{tab:avg}); the fitted exponent against the period is $-1.1$.

\begin{table}[ht]\centering\small
\begin{tabular}{@{}lrrrrrr@{}}\toprule
period $2\pi/\Phi$ (steps) & 64 & 128 & 256 & 512 & 1024 & 2048\\\midrule
ray through $k_0$, with connection & $0.0748$ & $0.0385$ & $0.0381$ & $0.0622$ & $0.1174$ & $0.2315$\\
averaged ray \eqref{eq:avg} & $0.0735$ & $0.0273$ & $0.0135$ & $0.0071$ & $0.0029$ & $0.0016$\\\bottomrule
\end{tabular}
\caption{Largest deviation, in sites, of the centroid from the prediction over one period. Packet at $k_0=1$ on the first branch, width $0.03$ in $k$, ring of $8192$ sites. The second row is sampled at $32$ times per period. For a packet of width $0.015$ the second row gives $0.0075$ and $0.0018$ at the periods $512$ and $2048$ \cite{comp4s}.}\label{tab:avg}
\end{table}

Equation \eqref{eq:avg} is what the adiabatic approximation gives: a vector that stays on one branch has the form $c(k+t\Phi)\,e^{-i\int\omega}u(k)$ in the quasi-momentum representation, in which position is $i\partial_k$, so that the mean position is the mean derivative of the phase plus the mean of $\mathcal{A}$. For an electron in a periodic potential with slowly varying external potentials Panati, Spohn and Teufel prove that the dynamics in an isolated band is governed, to first order in the slowness parameter, by semiclassical equations in which the Berry connection of the band enters \cite{pst}. Their theorem concerns Schr\"odinger operators in continuous time. No corresponding bound is proved here for the step of this paper; the criterion of $0.02$ site of \cite{g9} is met by the averaged prediction from the period $256$ on, and the failed predicate of \cite{g9} stands for the ray through $k_0$.

\section{Assumptions and limits}\label{sec:limits}

\emph{Assumed.} (1) Events are sets of histories (Section~\ref{sec:cyl}). This replaces ``consistency under continuation'' in the list of \cite{g9} by a statement about events; it is not derived from the carrier. (2) The rule is a sum over end vertices, and it depends on the walks that end at a vertex only through their gains. Proposition~\ref{prop:add} concerns additive rules, and Theorem~\ref{thm:kernel} with Corollary~\ref{cor:kernel2} pair rules with a kernel. In the other statements up to Corollary~\ref{cor:abund} the rule is further assumed to have the form \eqref{eq:phi}: the gains at a common end vertex are summed and the rule is a function of the modulus of the sum. Corollary~\ref{cor:block} drops the modulus and Corollary~\ref{cor:multicar} drops the sum, each on the carriers named there; Corollary~\ref{cor:vertexlaw} lets the function depend on the end vertex. (3) The carrier: for Proposition~\ref{prop:add} a digraph without loops and of constant out-degree, and a rule that is invariant under switching; for Theorems~\ref{thm:quad} and \ref{thm:deg2} a signed norm-preserving digraph of degree two, for Theorem~\ref{thm:power} a norm-preserving digraph with a merging pair of walks; without merging walks every rule of the form is the count. For Theorems~\ref{thm:conj} and \ref{thm:fourier}, Proposition~\ref{prop:deg34}, Corollary~\ref{cor:vsl} and the corollaries on blocks, in addition an abundance of gains at one length, which Proposition~\ref{prop:abund} decides for finite carriers with $A^*A=d$, $d\ge2$; for Theorem~\ref{thm:kernel}(ii) a hypothesis on the quotients of gains; for the last conclusion of Corollary~\ref{cor:multicar} a block whose entries generate the group; for Corollary~\ref{cor:vertexlaw} a finite digraph of constant out-degree $d\ge2$ without parallel arcs, with every gain reaching every vertex in abundance at one length, and strongly connected for its second statement; for Theorem~\ref{thm:periodic} the same digraph with a strongly connected derived digraph and no hypothesis of abundance. (4) The group and the character, as in \cite{g8}; in Theorems~\ref{thm:multiset}, \ref{thm:vertexlaw} and \ref{thm:periodic} a finite abelian group.

\emph{Obtained under these assumptions.} The second order of the rule; the exclusion of first-order weights of returning walks; the classification of the consistent rules of the form \eqref{eq:phi} for gains in the roots of unity as combinations over the conjugate characters; the uniqueness of the rule on the signed lift; the classification of the consistent pair rules with a kernel; on carriers with a full block and abundant gains, the second order for rules that add a function of the summed gains, and the modulus for those that are also unchanged by switching; and, where the entries of the block generate the group, the addition of gains through a character for every part other than the count of a rule that adds a function of the multiset of gains; and, on finite strongly connected carriers of constant out-degree without parallel arcs, with abundant gains at every vertex, that functions vanishing on the empty multiset are the same at every end vertex.

\emph{Not obtained.} (5) A reason for (1) to (4). Nothing here selects the embedding of the amplitudes: in degree two with abundant gains, for orders other than $2,4,6$ every conjugate character, and every combination, is consistent; multiplicativity on products removes the combinations and leaves the choice among the conjugate characters. On the signed lift there is no such choice (Corollary~\ref{cor:lift}); what is then not obtained is a reason for treating one character as a system, and a reason for or against the gain being a record. Among the pair rules, composition leaves one character in which $-1$ acts as $-1$ and does not make it faithful (Corollary~\ref{cor:kernel2}). Nothing selects between the count and the squares in Corollary~\ref{cor:multicar}, or among the characters that occur there; on the carriers of Theorem~\ref{thm:periodic}, multiplicativity on the product events of the carrier with itself at large lengths excludes their mixtures (Corollary~\ref{cor:composition}) and leaves the choice among the count and the single squares open. Nor does anything here select a carrier on which a rule that depends on the gains is consistent: on a digraph of constant out-degree with trivial gains the count is consistent and classical. (6) A classification of consistent rules that are not sums over end vertices, or that depend on more of a walk than its end vertex and its gain; the description of the solutions of Theorem~\ref{thm:multiset} for groups that are not abelian beyond Remark~\ref{rem:multisetna}, whose description of the forms rests on one challenger's reading and on computed instances for two groups, with no representation of complex type among them; which of the solutions of Theorem~\ref{thm:multiset} are consistent on a whole carrier when the entries of the block do not generate the group. (7) A ray limit with a force for gains of bounded order, or a proof that there is none; Proposition~\ref{prop:local} concerns local spectra only. (8) An error bound for \eqref{eq:avg}. (9) Everything listed as not obtained in \cite{g8,g9}: a clock, a selection of complex amplitudes, composition compatible with the step, and any physical correspondence.

\section{Status of the statements}

\begin{center}\small
\begin{tabular}{@{}p{0.55\textwidth}p{0.38\textwidth}@{}}\toprule
Statement & Kind\\\midrule
Lemmas~\ref{lem:cyl}, \ref{lem:par}; Propositions~\ref{prop:add}, \ref{prop:compo}, \ref{prop:deg34}, \ref{prop:abund}, \ref{prop:exact}, \ref{prop:local}; Theorems~\ref{thm:quad}, \ref{thm:power}, \ref{thm:deg2}, \ref{thm:conj}, \ref{thm:kernel}, \ref{thm:fourier}, \ref{thm:vsl}, \ref{thm:block}, \ref{thm:multiset}, \ref{thm:vertexlaw}, \ref{thm:periodic}; Lemmas~\ref{lem:maxp}, \ref{lem:classes}; Corollaries~\ref{cor:lift}, \ref{cor:kernel2}, \ref{cor:vsl}, \ref{cor:abund}, \ref{cor:block}, \ref{cor:multiset}, \ref{cor:multicar}, \ref{cor:vertexlaw}, \ref{cor:periodic}, \ref{cor:composition} & derived results, proofs given here\\
Tables~\ref{tab:stair}, \ref{tab:avg}; the values in Sections~\ref{sec:add}, \ref{sec:mod}, \ref{sec:force}, \ref{sec:ray} & computed instances\\
Equation \eqref{eq:avg} & heuristic; compared with computed instances; not proved\\\bottomrule
\end{tabular}
\end{center}

The components are elementary or known: cylinder sets of histories, the parallelogram identity as a characterization of the square, the isometries of $\ell_p$, characters of finite abelian groups and their linear independence, the Nullstellensatz, finite differences and polynomial functions on lattices, Hadamard matrices of Butson type, periods of digraphs, band spectra of a walk with a piecewise constant phase, and the semiclassical motion in a band with its Berry connection. The paper applies them to the pair measure on a gain digraph and claims no priority. The statements have not been refereed, and the numerical tables are not error enclosures.
